% Options for packages loaded elsewhere
\PassOptionsToPackage{unicode}{hyperref}
\PassOptionsToPackage{hyphens}{url}
\documentclass[
  ukrainian,
]{article}
\usepackage{xcolor}
\usepackage[margin=2.5cm]{geometry}
\usepackage{amsmath,amssymb}
\setcounter{secnumdepth}{-\maxdimen} % remove section numbering
\usepackage{iftex}
\ifPDFTeX
  \usepackage[T1]{fontenc}
  \usepackage[utf8]{inputenc}
  \usepackage{textcomp} % provide euro and other symbols
\else % if luatex or xetex
  \usepackage{unicode-math} % this also loads fontspec
  \defaultfontfeatures{Scale=MatchLowercase}
  \defaultfontfeatures[\rmfamily]{Ligatures=TeX,Scale=1}
\fi
\usepackage{lmodern}
\ifPDFTeX\else
\setmainfont{DejaVu Serif}
\setsansfont{DejaVu Sans}
\setmonofont{DejaVu Sans Mono}
\setmathfont{STIX Two Math}
\fi
\ifPDFTeX\else
  % xetex/luatex font selection
\fi
% Use upquote if available, for straight quotes in verbatim environments
\IfFileExists{upquote.sty}{\usepackage{upquote}}{}
\IfFileExists{microtype.sty}{% use microtype if available
  \usepackage[]{microtype}
  \UseMicrotypeSet[protrusion]{basicmath} % disable protrusion for tt fonts
}{}
\makeatletter
\@ifundefined{KOMAClassName}{% if non-KOMA class
  \IfFileExists{parskip.sty}{%
    \usepackage{parskip}
  }{% else
    \setlength{\parindent}{0pt}
    \setlength{\parskip}{6pt plus 2pt minus 1pt}}
}{% if KOMA class
  \KOMAoptions{parskip=half}}
\makeatother
\usepackage{longtable,booktabs,array}
\newcounter{none} % for unnumbered tables
\usepackage{calc} % for calculating minipage widths
% Correct order of tables after \paragraph or \subparagraph
\usepackage{etoolbox}
\makeatletter
\patchcmd\longtable{\par}{\if@noskipsec\mbox{}\fi\par}{}{}
\makeatother
% Allow footnotes in longtable head/foot
\IfFileExists{footnotehyper.sty}{\usepackage{footnotehyper}}{\usepackage{footnote}}
\makesavenoteenv{longtable}
\usepackage{graphicx}
\graphicspath{{/Users/kalinicenkovitalijmikolajovic/.codex/worktrees/5c15/Script_for_translate_book/Output_Final/}}
\setkeys{Gin}{width=0.82\linewidth,height=0.82\textheight,keepaspectratio}
\makeatletter
\newsavebox\pandoc@box
\newcommand*\pandocbounded[1]{% scales image to fit in text height/width
  \sbox\pandoc@box{#1}%
  \Gscale@div\@tempa{\textheight}{\dimexpr\ht\pandoc@box+\dp\pandoc@box\relax}%
  \Gscale@div\@tempb{\linewidth}{\wd\pandoc@box}%
  \ifdim\@tempb\p@<\@tempa\p@\let\@tempa\@tempb\fi% select the smaller of both
  \ifdim\@tempa\p@<\p@\scalebox{\@tempa}{\usebox\pandoc@box}%
  \else\usebox{\pandoc@box}%
  \fi%
}
% Set default figure placement to htbp
\def\fps@figure{htbp}
\makeatother
\ifLuaTeX
\usepackage[bidi=basic,shorthands=off,provide=*]{babel}
\else
\usepackage[bidi=default,shorthands=off,provide=*]{babel}
\fi
\ifLuaTeX
  \usepackage{selnolig} % disable illegal ligatures
\fi
\setlength{\emergencystretch}{3em} % prevent overfull lines
\providecommand{\tightlist}{%
  \setlength{\itemsep}{0pt}\setlength{\parskip}{0pt}}
\usepackage{bookmark}
\IfFileExists{xurl.sty}{\usepackage{xurl}}{} % add URL line breaks if available
\urlstyle{same}
\hypersetup{
  pdflang={uk},
  hidelinks,
  pdfcreator={LaTeX via pandoc}}

\author{}
\date{}

\begin{document}

\pandocbounded{\includegraphics[keepaspectratio]{images/_page_362_Figure_7.jpeg}}

в начальные условия

\[
\begin{cases} 0, x < a, \\ 1, x \geqslant a, \end{cases}
\]

найдем сначала стоячне волны, т. е. решения вида

Подставив это выражение в уравнение (20.1) и разделив переменные,
получим два обыкновенных дифференциальных уравнения

\[
H(x-a) =
\]

Функція Хевісайду,

\[
H(a-x) = \begin{cases} 1, & x \leqslant a, \\ 0, & x > a, \end{cases}
\]

дзеркальною функцією Хевісайду.

ТАБЛИЦЯ A. Експоненціальне перетворення Фур'є

1. f'(x) \textbar{} iωF (ω) \textbar{}

\textbar------------------------------\textbar-------------------------------------------------\textbar{}

\begin{verbatim}
| 2. f''(x) | -\omega^2 F(\omega) |
\end{verbatim}

\begin{verbatim}
| 3. fn (x), (n-та похідна) | (i\omega)^n F(\omega) |
\end{verbatim}

\begin{verbatim}
| 4. f(ax), a > 0 | \frac{1}{a} F\left(\frac{\omega}{a}\right) |
\end{verbatim}

\begin{verbatim}
| 5. f(x-a) | e^{-ia\omega}F(\omega) |
\end{verbatim}

\begin{verbatim}
| 6. e^{-a^2x^2} | \frac{1}{a\sqrt{2}}e^{-\omega^2/4a^2} |
\end{verbatim}

\begin{verbatim}
| 7. e^{-a_1 x} | \sqrt{\frac{2}{\pi}} \frac{a}{a^2 + \omega^2} |
\end{verbatim}

\begin{enumerate}
\def\labelenumi{\arabic{enumi}.}
\setcounter{enumi}{19}
\item
  гріх (ax)
\item
\end{enumerate}

\[
\begin{cases} 0, & m \neq n, \\ L/2, & m = n, \end{cases}
\]

\begin{enumerate}
\def\labelenumi{\arabic{enumi}.}
\setcounter{enumi}{10}
\tightlist
\item
\item
\item
\item
\item
\item
\end{enumerate}

\begin{verbatim}
\begin{aligned} (\text{Y} \forall \Pi) & \qquad \phi_{\text{vx}} + \phi_{yy} = 0, \quad 0 < x < \infty, \quad 0 < y < \infty, \\ \left\{ \begin{array}{l} \phi(x, \ 0) = \begin{cases} 1, & 0 < x < 1, \\ 0, & 1 \leqslant x < \infty, \end{cases} \\ \phi(0, \ y) = \begin{cases} 1, & 0 < y < 1, \\ 0, & 1 \leqslant y < \infty. \end{cases} \end{aligned} \end{array}
\end{verbatim}

\begin{enumerate}
\def\labelenumi{\arabic{enumi}.}
\setcounter{enumi}{16}
\tightlist
\item
\item
\end{enumerate}

Зверніть увагу на зображення.

\begin{enumerate}
\def\labelenumi{\arabic{enumi}.}
\setcounter{enumi}{18}
\item
\item
\item
  Розв'язати змішану задачу Діріхле-Неймана з кутовим порошком за
  допомогою розчину під кутом 45°
\item
\end{enumerate}

\begin{verbatim}
ПРИМІТКА. Комплексна функція  w = \ln z = \ln |z| + i \arg(z)  відображає промінь  \theta = c_1  у площині z у пряму  v = c_1  , а коло  r = c_2  у пряму  u = \ln c_2  у площині w (див. рисунки).
\end{verbatim}

\begin{enumerate}
\def\labelenumi{\arabic{enumi}.}
\setcounter{enumi}{10}
\tightlist
\item
\end{enumerate}

\section{ТАБЛИЦІ ІНТЕГРАЛЬНИХ
ПЕРЕТВОРЕНЬ}\label{ux442ux430ux431ux43bux438ux446ux456-ux456ux43dux442ux435ux433ux440ux430ux43bux44cux43dux438ux445-ux43fux435ux440ux435ux442ux432ux43eux440ux435ux43dux44c}

Таблиця A: Експоненціальне перетворення Фур'є,

Таблиця B: Перетворення синуса Фур'є.

Таблиця St Cosine Фур'є-перетворення.

Таблиця Dr скінченне перетворення Фур'є, таблиця та скінченне косинусне
перетворення Фур'є.

Таблиця D: Трансформація Лапласа.

\begin{verbatim}
# Прийняті позначення  \delta\left(x\right)  — дельта-функція,
\end{verbatim}

\begin{enumerate}
\def\labelenumi{\arabic{enumi}.}
\setcounter{enumi}{11}
\tightlist
\item
\item
\item
\end{enumerate}

\[
\begin{array}{llll} (\mathrm{Y} \Pi \Pi) & u_{tt} = \alpha^2 u_{xx} - \beta u_t, & 0 < x < 1, & 0 < t < \infty, \\ (\Gamma \mathrm{Y}) & \begin{cases} u \ (0, \ t) = 0, \\ u \ (1, \ t) = 0, \end{cases} & 0 < t < \infty, \\ (\mathrm{H} \mathrm{Y}) & \begin{cases} u \ (x, \ 0) = f \ (x), \\ u_1 \ (x, \ 0) = 0. \end{cases} & 0 \leqslant x \leqslant 1, \end{array}
\]

\begin{enumerate}
\def\labelenumi{\arabic{enumi}.}
\setcounter{enumi}{17}
\tightlist
\item
\end{enumerate}

Воспользовавшись соотношениями ортогональности

\[
\begin{array}{ll} (\mathrm{Y} \mathrm{H} \Pi) & u_{tt} = \alpha^2 u_{xx} + Kx, \quad 0 < x < 1, \ 0 < t < \infty, \\ (\mathrm{\Gamma} \mathrm{Y}) & \begin{cases} u \ (0, \ t) = 0, \\ u \ (1, \ t) = 0, \end{cases} & 0 < t < \infty, \\ (\mathrm{H} \mathrm{Y}) & \begin{cases} u \ (x, \ 0) = f \ (x), \\ u_t \ (x, \ 0) = 0. \end{cases} & 0 \leqslant x \leqslant 1, \end{array}
\]

\begin{enumerate}
\def\labelenumi{\arabic{enumi}.}
\setcounter{enumi}{8}
\tightlist
\item
\item
\item
\end{enumerate}

\pandocbounded{\includegraphics[keepaspectratio]{images/_page_154_Figure_5.jpeg}}

находим коэффициенты \(a_n\) и \(b_n\) (20.5)

\section{\texorpdfstring{ТАБЛИЦЯ B. Перетворення Спнуса-Фур'є
\(i \sqrt{\frac{\pi}{2}} [\delta(\omega+a) - \delta(\omega-a)]\)}{ТАБЛИЦЯ B. Перетворення Спнуса-Фур'є i \textbackslash sqrt\{\textbackslash frac\{\textbackslash pi\}\{2\}\} {[}\textbackslash delta(\textbackslash omega+a) - \textbackslash delta(\textbackslash omega-a){]}}}\label{ux442ux430ux431ux43bux438ux446ux44f-b.-ux43fux435ux440ux435ux442ux432ux43eux440ux435ux43dux43dux44f-ux441ux43fux43dux443ux441ux430-ux444ux443ux440ux454-i-sqrtfracpi2-deltaomegaa---deltaomega-a}

Итак, задача решена. Ее решение записывается в виде

(20.6)

\[
F(\omega) = \frac{?}{\pi} \int_{0}^{\infty} f(x) \sin(\omega x) dx
\]

а коэффициенты \(a_n\) и \(b_n\) находятся по формулам (20.5). Прежде
чем закончить лекцию, сделаем несколько замечаний.

\section{ЗАМЕЧАНИЯ}\label{ux437ux430ux43cux435ux447ux430ux43dux438ux44f}

\begin{enumerate}
\def\labelenumi{\arabic{enumi}.}
\tightlist
\item
  Если начальная \emph{скорость} струны равна нулю, то решение (20.6)
  принимает вид
\end{enumerate}

\[
0 < \omega < x
\]

І це можна інтерпретувати так. Припустимо, ми розклали початкове
зміщення струни в

\pandocbounded{\includegraphics[keepaspectratio]{images/_page_155_Figure_8.jpeg}}

ряд по синусам

f'\,'(0)

Тогда каждый член такого разложения должен совершать колебания

\[
\begin{aligned} & (\text{Y} \cdot \Pi) \quad u_{tt} = -u_{xxxx}, \quad 0 < x < 1, \quad 0 < t < \infty, \\ & \begin{cases} u \cdot (0, \ t) = 0, \\ u_{xx} \cdot (0, \ t) = 0, \\ u \cdot (1, \ t) = 0, \\ u_{xx} \cdot (1, \ t) = 0, \end{cases} \quad 0 < t < \infty, \\ & \begin{cases} u \cdot (x, \ 0) = f(x), \\ u_{1}(x, \ 0) = g(x). \end{cases} \quad 0 \leqslant x \leqslant 1, \end{aligned}
\]

Якщо ми підсумуємо всі ці коливання, отримаємо розв'язок нашої проблеми.
Наприклад, припустимо, що

\begin{verbatim}
начальное смещение струны задается функцией  f(x) = \sin(\pi x/L) + 0.5 \sin(3\pi x/ L) + 0,25 \sin(5\pi x/L)  .
\end{verbatim}

Полный отклик системы на начальное условие такого вида равен сумме
откликов на каждое слагаемое, т. е.

\begin{enumerate}
\def\labelenumi{\arabic{enumi}.}
\setcounter{enumi}{3}
\tightlist
\item
\item
  Слагаемое с номером n в разложении (20.6) называется \textbf{п-й модой
  колебаний} или \textbf{п-й гармоникой}. Используя тригонометрические
  формулы, можно записать эти гармоники в виде
\end{enumerate}

\[
[2/\pi\omega]^{1/2}
\]

\begin{verbatim}
где  R_n  и  \delta_n  — новые произвольные постоянные (амплитуда и фазовый угол). Эта новая форма представления n-й моды оказывается более полезной при анализе колебаний. Частота колебаний од (рад/с) и-й моды находится по формуле
\end{verbatim}

\[
\frac{2}{\pi\omega} [1-\cos(\omega a)]
\]

де T --- натяг струни, p --- лінійна густина струни. Зверніть увагу, що
частота n-ї гармоніки у n разів більша за основну частоту (n=1). Ця
властивість частоти не характерна для всіх типів коливань. Приємний звук
струни скриму або гітари зумовлений множинністю обертонів основної
частоти. У барабані частоти гармонік високого порядку не є кратними
основної частоти коливань.

\section{ЗАДАЧИ}\label{ux437ux430ux434ux430ux447ux438}

\begin{enumerate}
\def\labelenumi{\arabic{enumi}.}
\tightlist
\item
  Найдите решение задачи (20.1), если начальные условия задаются
  формулами
\end{enumerate}

1

\(u_t\)(x, 0) = 0. Постройте графики решения для различных моментов
времени. Является ли это решение периодическим по времени? Каков его
период?

\begin{enumerate}
\def\labelenumi{\arabic{enumi}.}
\setcounter{enumi}{1}
\tightlist
\item
  Найдите решение задачи о колебаниях струны (20.1) при начальных
  условиях вида
\end{enumerate}

\(u_t(x, 0) = \sin(3\pi x/L).\) Как выглядит график решения в различные
моменты времени?

\begin{itemize}
\tightlist
\item
  \begin{enumerate}
  \def\labelenumi{\arabic{enumi}.}
  \setcounter{enumi}{2}
  \tightlist
  \item
    Показати це, якщо на рис. 20.2 значення \(\lambda \geqslant 0\) ,
    тоді розв'язки \(X(x)\) \(T(t)\) будуть або необмеженими, або
    ідентичними нулю.
  \end{enumerate}
\item
  Знайти розв'язок проблеми коливань струн за початкових умов вигляду
\end{itemize}

\[
\frac{1}{2}e^{-\alpha}\sin\omega
\]

\(u_1(x, 0) = (3\pi \alpha/L)\sin(3\pi x/L).\) 5. Гитарная струна длины
L=1 оттягивается за среднюю точку на величину h (см. рис.) Начальное
положение струны можно задать в виде

\[
\frac{1-e^{-a\omega}}{\omega}
\]

\begin{enumerate}
\def\labelenumi{\arabic{enumi}.}
\setcounter{enumi}{13}
\tightlist
\item
\item
\end{enumerate}

\[
F \sinh \sqrt{\omega} = 0,
\]

\begin{enumerate}
\def\labelenumi{\arabic{enumi}.}
\setcounter{enumi}{16}
\tightlist
\item
\end{enumerate}

\[
f(x) = \int_{0}^{\infty} F(\omega) \cos(\omega x) d\omega \qquad \qquad F(x) = \frac{2}{\pi} \int_{0}^{\infty} f(x) \cos(\omega x) dx
\]

\begin{enumerate}
\def\labelenumi{\arabic{enumi}.}
\setcounter{enumi}{18}
\tightlist
\item
\end{enumerate}

\[
D \sin \sqrt{\omega} = 0.
\]

\begin{enumerate}
\def\labelenumi{\arabic{enumi}.}
\setcounter{enumi}{18}
\tightlist
\item
\item
\end{enumerate}

\section{КОЛЕБАНИЯ БАЛКИ (УРАВНЕНИЕ С ЧАСТНЫМИ ПРОИЗВОДНЫМИ ЧЕТВЕРТОГО
ПОРЯДКА)}\label{ux43aux43eux43bux435ux431ux430ux43dux438ux44f-ux431ux430ux43bux43aux438-ux443ux440ux430ux432ux43dux435ux43dux438ux435-ux441-ux447ux430ux441ux442ux43dux44bux43cux438-ux43fux440ux43eux438ux437ux432ux43eux434ux43dux44bux43cux438-ux447ux435ux442ux432ux435ux440ux442ux43eux433ux43e-ux43fux43eux440ux44fux434ux43aux430}

ЦЕЛЬ ЛЕКЦИИ: Показать, каким образом задача о колебании балки приводит к
уравнению четвертого порядка и как эта задача решается в случае свободно
опирающихся концов. Проводится сравнение колебаний балки с колебаниями
скрипичной струны.

\[
F = 0
\]

Головна різниця між поперечними коливаннями тонкої балки та поперечними
коливаннями струни полягає в тому, що промінь чинить опір згину. Не
заглиблюючись у механіку тонких балок, припустимо, що, враховуючи опір
згинання (замість рівняння хвилі) веде до рівняння четвертого порядку

(21.1)

\section{ТАБЛИЦЯ D. Скінченне перетворення синуса
Фур'є}\label{ux442ux430ux431ux43bux438ux446ux44f-d.-ux441ux43aux456ux43dux447ux435ux43dux43dux435-ux43fux435ux440ux435ux442ux432ux43eux440ux435ux43dux43dux44f-ux441ux438ux43dux443ux441ux430-ux444ux443ux440ux454}

У скінченному перетворенні синуса Фур'є функція \(f(x)\), визначена в
інтервалі \(0 \le x \le \pi\) , перетворюється у послідовність \(S_n\) ,
n = 1, 2 \ldots{} За формулою

\[
rge  \alpha^2 = K/\rho  ,
\]

K---модуль сдвига (чем больше K, тем жестче балка, тем выше частота
колебаний),

р --- линейная плотность балки (масса/ед. длины).

Поскольку читатель впервые встречается в нашей книге с уравнением выше
чем второго порядка, будет полезно решить типовую задачу о колебаниях
балки. Позже мы познакомимся и с другими задачами теории балок.

\section{Свободно опирающаяся
балка}\label{ux441ux432ux43eux431ux43eux434ux43dux43e-ux43eux43fux438ux440ux430ux44eux449ux430ux44fux441ux44f-ux431ux430ux43bux43aux430}

Рассмотрим малые колебания тонкой балки, концы которой свободно
опираются на две опоры. Говоря «свободно опираются», мы подразумеваем,
что концы балки не перемещаются, но на-

клоны балки в концевых точках могут изменяться (концы балки вакреплены с
помощью штифтовых устройств, рис. 21.1).

Совершенно ясно, что на концах балки должны выполняться граничные
условия

\[
\omega_n = (n\pi)^2,
\]

\(u(1, t)\) = 0, но совсем не очевидно, что на концах балки должны
удовлетворяться еще два граничных условия:

Коли йдеться про таблиці, зручно розглядати змінну x як належну до
інтервалу {[}0, l{]}. Будь-який інший чисельний інтервал {[}a, b{]}
можна перетворити на \([0, \pi]\) за допомогою формули

\(u_{xx}\)(1, t) = 0. Використовуючи теорію тонких балок, можна
показати, що момент згинання в балці пропорційний значенню \(u_{xx}\) ,
а момент згинання на вільно лежачому кінці має дорівнювати P I C

\[
u_n(x, t) = X_n(x) T_n(t) = [a_n \sin(n\pi)^2 t + b_n \cos(n\pi)^2 t] \sin(n\pi x).
\]

.21.1. Свободно оппрающаяся балка.

вен нулю. Следовательно, изображенная на рис. 21.1 колеблющаяся балка
описывается следующей смешанной задачей (для простоты величина о принята
равной единице):

(21.2)

У скінченному синусоїдальному перетворенні немає нічого загадкового:
елементи послідовності \(\mathcal{S}_n\) збігаються з коефіцієнтами в
\(\sin(nx)\) у розкладі функції \(f(x)\) у ряд Фур'є, тобто

\begin{verbatim}
2.  f(\alpha x) 
\end{verbatim}

\[
u(x, t) = \sum_{n=1}^{\infty} \left[ a_n \sin(n\pi)^2 t + b_n \cos(n\pi)^2 t \right] \sin(n\pi x)
\]

\begin{verbatim}
3.  e^{-\alpha x} 
\end{verbatim}

\[
f(x) = \sum_{n=1}^{\infty} S_n \sin(nx)
\]

Для решения этой задачи воспользуемся методом разделения переменных.
Будем искать только периодические решения, т. е. колебания вида

(21.3)

\[
u(x, 0) = f(x) = \sum_{n=1}^{\infty} b_n \sin(n\pi x),
\]

Зверніть увагу, що вибір розв'язку вигляду (21.3) фактично означає, що
ми взяли константу розділення у методі розділення змінних як від'ємну.

\begin{verbatim}
Подставим теперь (21.3) в уравнение колебаний балки и получим уравнение для  X\left(x\right)
\end{verbatim}

\[
0 \le x \le \pi
\]

общее решение которого записывается в виде

\[
u_t(x, 0) = g(x) = \sum_{n=1}^{\infty} (n\pi)^2 a_n \sin(n\pi x).
\]

Для определения коэффициентов C, D, E и F подставим общее решение в
граничные условия

\begin{enumerate}
\def\labelenumi{\arabic{enumi}.}
\tightlist
\item
\end{enumerate}

\[
a_n = \frac{2}{(n\pi)^2} \int_0^1 g(x) \sin(n\pi x) dx,
\]

\begin{enumerate}
\def\labelenumi{\arabic{enumi}.}
\setcounter{enumi}{8}
\tightlist
\item
\end{enumerate}

\[
b_n = 2 \int_0^1 f(x) \sin(n\pi x) dx.
\]

\(\frac{n}{n^2 + a^2} [1 - (-1)^n e^{ax}]\) 10.
\(\frac{\sin a (\pi - x)}{\sin a\pi}\) \(\frac{2n}{\pi (n^2 + a^2)}\)
ТАБЛИЦЯ E. Скінченне перетворення Косніуса Фур'є

У скінченному косинусному перетворенні Фур'є функція \(f(x)\), визначена
на інтервалі \(0 \le x \le \pi\) , перетворюється у послідовність
\(C_n\) , \(n=0,\ 1,\ 2,\ \ldots\) , за формулою

\[
u(x, 0) = \sin(\pi x) + 0.5 \sin(3\pi x),
\]

Коли йдеться про таблиці, зручно розглядати змінну x як таку, що
належить до інтервалу \([0,\pi]\) . Будь-який інший числовий інтервал
{[}a,b{]} можна перетворити на \([0,\pi]\) за допомогою формули

\[
a_n=0
\]

У скінченному косинусному перетворенні немає нічого загадкового:
елементи послідовності \(C_n\) збігаються з коефіцієнтами, коли
\(\cos{(nx)}\) у косинусному розкладі функції \(f(x)\) у ряд Фур'є,
тобто

\[
u(x, t) = \cos(\pi^2 t) \sin(\pi x) + 0.5 \cos(9\pi^2 t) \sin(3\pi x)
f(x) = \frac{C_0}{2} + \sum_{n=1}^{\infty} C_n \cos(nx) \qquad C_n = \frac{2}{\pi} \int_{0}^{\pi} f(x) \cos nx \, dx
u(x, t) = \cos(\pi t) \sin(\pi x) + 0.5 \cos(3\pi t) \sin(3\pi x)
\]

\begin{enumerate}
\def\labelenumi{\arabic{enumi}.}
\tightlist
\item
\end{enumerate}

\[
u(0, t) = 0,
-n^{2}C_{n} - \frac{2}{\pi} \{f'(0) - (-1)^{n} f'(n)\}
u(x, t) = \sum_{n=1}^{\infty} X_n(x) \left[ a_n \sin(\omega_n t) + b_n \cos(\omega_n t) \right],
\begin{cases} 2, &n=0 \\ 0, &n=1, 2, \dots \end{cases}
\begin{cases} 1, & n=m \\ 0, & n\neq m \end{cases}
\]

x

\[
\quad 0 \leqslant x \leqslant 1,
u(x, 0) = \sin(\pi x),
\]

\begin{enumerate}
\def\labelenumi{\arabic{enumi}.}
\setcounter{enumi}{3}
\tightlist
\item
  Пусть левый конец (x=0) балки жестко закреплен в стену, а правый конец
  (x=1) свободно опирается, как показано на рис. 21.3. Решите задачу для
  балки с такими граничными условиями и объясните, как найти собственные
  частоты колебаний в этом случае. Знать собственные частоты балок очень
  важно, чтобы избежать резонанса с внешними воздействиями.
\end{enumerate}

Лекция 22

\section{ПЕРЕХОД К БЕЗРАЗМЕРНЫМ
ПЕРЕМЕННЫМ}\label{ux43fux435ux440ux435ux445ux43eux434-ux43a-ux431ux435ux437ux440ux430ux437ux43cux435ux440ux43dux44bux43c-ux43fux435ux440ux435ux43cux435ux43dux43dux44bux43c}

ЦЕЛЬ ЛЕКЦИИ: Показать, каким образом краевые, начальные и другие виды
физических задач можно записать в безразмерном виде. В конкретных
областях знания (физике, химии, биологии или экономике) уравнения
записываются по-разному.

Переход к безразмерным переменным позволяет привести уравнения к одному
и тому же виду. Именно по этой причине при изучении уравнений с частными
производными в математике стараются отвлекаться от физического смысла
входящих в них параметров. Вот почему химик, физик или биолог должен
преобразовать свое уравнение к той форме, которая принята в нашей книге.

Основная идея введения новых (безразмерных) переменных состоит в том,
что после перехода к безразмерным переменным вадача становится чисто
математической и уже не содержит

карактерных физических констант. Именно таким способом многие различные
уравнения физики, биологии и химии, содержащие всякие нюансы, связанные
с физическими параметрами, приводятся к одной и той же простой форме
(см. рис. 22.1).

Познакомимся с этим преобразованием на простом примере.

\pandocbounded{\includegraphics[keepaspectratio]{images/_page_161_Figure_4.jpeg}}

РИС. 22.1. Некоторые задачи, приводящиеся к одной и той же безразмерной
форме.

\section{Приведение диффузионной задачи к безразмерному
виду}\label{ux43fux440ux438ux432ux435ux434ux435ux43dux438ux435-ux434ux438ux444ux444ux443ux437ux438ux43eux43dux43dux43eux439-ux437ux430ux434ux430ux447ux438-ux43a-ux431ux435ux437ux440ux430ux437ux43cux435ux440ux43dux43eux43cux443-ux432ux438ux434ux443}

\begin{verbatim}
Начнем со смешанной задачи для отрезка с начальным условием на температуру  u\left(x,\,0\right)=\sin\left(\pi x/L\right)  и граничными температурами на концах  T_1  и  T_2 . Другими словами, мы решаем задачу
\end{verbatim}

(22.1)

\begin{verbatim}
\begin{array}{ccccccccccccccccccccccccccccccccccc
\end{verbatim}

\pandocbounded{\includegraphics[keepaspectratio]{images/_page_161_Figure_10.jpeg}}

\[
(\text{УЧП}) \quad u_{tt} + u_{xxxx} = 0, \quad 0 < x < 1, \quad 0 < t < \infty,
u \to U
(\text{ГУ}) \quad
U(x, t) = \frac{u(x, t) - T_1}{T_2 - T_1}
\begin{array}{ll} \text{(УЧП)} & U_t = \alpha^2 U_{xx}, \quad 0 < x < L, \quad 0 < t < \infty, \\ U_t = 0, \quad 0 < t < \infty, \\ U_t = 0, \quad 0 < t < \infty, \\ U_t = 0, \quad 0 < t < \infty, \\ U_t = 0, \quad 0 < t < \infty, \\ U_t = 0, \quad 0 < t < \infty, \\ U_t = 0, \quad 0 < t < \infty, \\ U_t = 0, \quad 0 < t < \infty, \\ U_t = 0, \quad 0 < t < \infty, \\ U_t = 0, \quad 0 < t < \infty, \\ U_t = 0, \quad 0 < t < \infty, \\ U_t = 0, \quad 0 < t < \infty, \\ U_t = 0, \quad 0 < t < \infty, \\ U_t = 0, \quad 0 < t < \infty, \\ U_t = 0, \quad 0 < t < \infty, \\ U_t = 0, \quad 0 < t < \infty, \\ U_t = 0, \quad 0 < t < \infty, \\ U_t = 0, \quad 0 < t < \infty, \\ U_t = 0, \quad 0 < t < \infty, \\ U_t = 0, \quad 0 < t < \infty, \\ U_t = 0, \quad 0 < t < \infty, \\ U_t = 0, \quad 0 < t < \infty, \\ U_t = 0, \quad 0 < t < \infty, \\ U_t = 0, \quad 0 < t < \infty, \\ U_t = 0, \quad 0 < t < \infty, \\ U_t = 0, \quad 0 < t < \infty, \\ U_t = 0, \quad 0 < t < \infty, \\ U_t = 0, \quad 0 < t < \infty, \\ U_t = 0, \quad 0 < t < \infty, \\ U_t = 0, \quad 0 < t < \infty, \\ U_t = 0, \quad 0 < t < \infty, \\ U_t = 0, \quad 0 < t < \infty, \\ U_t = 0, \quad 0 < t < \infty, \\ U_t = 0, \quad 0 < t < \infty, \\ U_t = 0, \quad 0 < t < \infty, \\ U_t = 0, \quad 0 < t < \infty, \\ U_t = 0, \quad 0 < t < \infty, \\ U_t = 0, \quad 0 < t < \infty, \\ U_t = 0, \quad 0 < t < \infty, \\ U_t = 0, \quad 0 < t < \infty, \\ U_t = 0, \quad 0 < t < \infty, \\ U_t = 0, \quad 0 < t < \infty, \\ U_t = 0, \quad 0 < t < \infty, \\ U_t = 0, \quad 0 < t < \infty, \\ U_t = 0, \quad 0 < t < \infty, \\ U_t = 0, \quad 0 < t < \infty, \\ U_t = 0, \quad 0 < t < \infty, \\ U_t = 0, \quad 0 < t < \infty, \\ U_t = 0, \quad 0 < t < \infty, \\ U_t = 0, \quad 0 < t < \infty, \\ U_t = 0, \quad 0 < t < \infty, \\ U_t = 0, \quad 0 < t < \infty, \\ U_t = 0, \quad 0 < t < \infty, \\ U_t = 0, \quad 0 < t < \infty, \\ U_t = 0, \quad 0 < t < \infty, \\ U_t = 0, \quad 0 < t < \infty, \\ U_t = 0, \quad 0 < t < \infty, \\ U_t = 0, \quad 0 < t < \infty, \\ U_t = 0, \quad 0 < t < \infty, \\ U_t = 0, \quad 0 < t < \infty, \\ U_t = 0, \quad 0 < t < \infty, \\ U_t = 0, \quad 0 < t < \infty, \\ U_t = 0, \quad 0 < t < \infty, \\ U_t = 0, \quad 0 < t < \infty, \\ U_t = 0, \quad 0 < t < \infty, \\ U_t = 0, \quad 0 < t < \infty, \\ U_t = 0, \quad 0 < t < \infty, \\ U_t = 0, \quad 0 < t < \infty, \\ U_t = 0, \quad 0 < t < \infty, \\ U_t = 0, \quad 0 < t < \infty, \\ U_t = 0, \quad 0 < t < \infty, \\ U_t = 0, \quad 0 < t < \infty, \\ U_t = 0, \quad 0 < t < \infty, \\ U_t = 0, \quad 0 < t < \infty, \\ U_t = 0, \quad 0 < t < \infty, \\ U_t = 0, \quad 0 < t < \infty, \\ U_t = 0, \quad 0 < t < \infty, \\ U_t = 0, \quad 0 < t < \infty, \\ U_t = 0, \quad 0 < t < \infty, \\ U_t = 0, \quad 0 < t < \infty, \\ U_t = 0, \quad 0 < t < \infty, \\ U_t = 0, \quad 0 < t < \infty, \\ U_t = 0, \quad 0 < t < \infty, \\ U_t \end{array}
\begin{cases} u(0, t) &= 0, \\ u_x(0, t) &= 0, \\ u_{xx}(1, t) &= 0, \\ u_{xxx}(1, t) &= 0, \end{cases}
u(x, t) = T_1 + (T_2 - T_1) U(x, t).
\xi = x/L
U_x = U_{\xi} \xi_x = \frac{1}{L} U_{\xi},
\begin{cases} u(x, 0) &= f(x), \\ u_t(x, 0) &= g(x). \end{cases}
U_{xx} = \frac{1}{L^2} U_{\xi\xi}
\]

\begin{verbatim}
\begin{array}{ll} (\text{YM}\Pi) & U_t = (\alpha/L)^2 \, U_{\xi\xi}, \quad 0 < \xi < 1, \quad 0 < t < \infty, \\ (22.3) & (\Gamma \text{Y}) & \begin{cases} U(0,\ t) = 0, \\ U(1,\ t) = 1, \\ U(\xi,\ 0) = \frac{\sin\left(\pi\xi\right) - T_1}{T_2 - T_1}, \quad 0 \leqslant \xi \leqslant 1. \end{cases}
\end{verbatim}

6 M 601

УКАЗАНИЕ. Хотя собственные функции \(X_n\)(x) в этой задаче уже не
обычные синусы, согласно общей теории Штурма --- Ли вилля они будут
ортогональны на отрезке {[}0, 1{]}.

\begin{enumerate}
\def\labelenumi{\arabic{enumi}.}
\setcounter{enumi}{1}
\tightlist
\item
  Найдите решение для свободно опирающейся на оба конца балки при
  начальных условиях
\end{enumerate}

\[
cu_{\mathbf{t}} = [\alpha/L]^2 u_{\xi\xi}.
\]

\begin{verbatim}
u_t(x, 0) = \sin(\pi x).   0 \le x \le 1,  3. Решите задачу 2 с начальными условиями
\end{verbatim}

\[
\tau = [\alpha/L]^2 t.
\quad 0 \le x \le 1,
\]

\begin{verbatim}
\begin{array}{cccccccccccccccccccccccccccccccccccc
\end{verbatim}

\[
u(x, t) = T_1 + (T_2 - T_1) U(x/L, \alpha^2 t/L^2).
\]

(см. рис. 22.2).

\pandocbounded{\includegraphics[keepaspectratio]{images/_page_161_Figure_10.jpeg}}

РИС. 22.2. Задача теплопроводности в плоскости переменных (х, т).

Наша цель --- дать новую эквивалентную (22.1) постановку задачи таким
образом, чтобы

в новой задаче не было никаких физических параметров (вроде α);

\begin{enumerate}
\def\labelenumi{\arabic{enumi})}
\setcounter{enumi}{1}
\tightlist
\item
  начальные и граничные условия стали проще.
\end{enumerate}

Для осуществления этой цели, мы введем три новых безразмерных переменных
U, \(\xi\) и \(\tau\) , чтобы заменить прежние u, x и t по схеме

\[
\begin{array}{llll} (\text{УЧП}) & u_{tt} = \alpha^2 u_{xx}, & 0 < x < L, & 0 < t < \infty, \\ u(0, t) = 0, & 0 < t < \infty, \\ u(L, t) = 0, & 0 < t < \infty, \\ u(X, 0) = \sin(\pi x/L) + (0.5)\sin(3\pi x/L), \\ u_t(x, 0) = 0. & \end{array}
\]

(безразмерная температура), \(x \to \xi\) (безразмерная длина),
\(t \to \tau\) (безразмерное время).

Мы проведем все три преобразования одновременно.

\section{\texorpdfstring{Преобразование зависимой переменной
\(u \rightarrow U\)}{Преобразование зависимой переменной u \textbackslash rightarrow U}}\label{ux43fux440ux435ux43eux431ux440ux430ux437ux43eux432ux430ux43dux438ux435-ux437ux430ux432ux438ux441ux438ux43cux43eux439-ux43fux435ux440ux435ux43cux435ux43dux43dux43eux439-u-rightarrow-u}

Определим функцию \(U(x, t)\) по формуле

\[
U(x, t) = \frac{u(x, t) - T_1}{T_2 - T_1}
\xi = x/L
\]

.

\begin{verbatim}
Ясно, что новая температура   U\left(x,\,t\right)   безразмерна, поскольку мы делим градусы Цельсия на градусы Цельсия. Новые граничные условия для   U\left(x,\,t\right)   при x=0 и x=L примут вид   U\left(0,\,t\right)=0   и   U\left(L,\,t\right)=1  . Давайте поставим задачу для новой функции   U\left(x,\,t\right)  . После несложных преобразований задача (22.1) переходит в новую задачу
\end{verbatim}

задачу

\[
(\text{YYII}) \qquad u_{\tau\tau} = u_{\xi\xi}, \quad 0 < \xi < 1, \quad 0 < \tau < \infty, 1
\]

\begin{verbatim}
Если угодно, то здесь можно остановиться, решить задачу для   U\left(x,\, t\right)   и затем найти   u\left(x,\,t\right)   по формуле
\end{verbatim}

\[
(\text{\GammaY}) \quad
u(x, t) = T_1 + (T_2 - T_1) U(x, t).
\]

Однако мы пойдем дальше и займемся преобразованием независимых
переменных x и t.

\section{\texorpdfstring{Преобразование пространственной координаты
\(x \rightarrow \xi\)}{Преобразование пространственной координаты x \textbackslash rightarrow \textbackslash xi}}\label{ux43fux440ux435ux43eux431ux440ux430ux437ux43eux432ux430ux43dux438ux435-ux43fux440ux43eux441ux442ux440ux430ux43dux441ux442ux432ux435ux43dux43dux43eux439-ux43aux43eux43eux440ux434ux438ux43dux430ux442ux44b-x-rightarrow-xi}

Представляется совершенно очевидным, как следует выбрать безразмерную
переменную \(\xi\) . Поскольку \(0 \leqslant x \leqslant L\) , полагаем

\[
\xi = x/L
\]

.

После вычисления производных

\[
U_x = U_{\xi} \xi_x = \frac{1}{L} U_{\xi},
U_{xx} = \frac{1}{L^2} U_{\xi\xi}
\begin{array}{ll} (\text{YM}\Pi) & U_t = (\alpha/L)^2 \, U_{\xi\xi}, \quad 0 < \xi < 1, \quad 0 < t < \infty, \\ (22.3) & (\Gamma \text{Y}) & \begin{cases} U(0,\ t) = 0, \\ U(1,\ t) = 1, \\ U(\xi,\ 0) = \frac{\sin\left(\pi\xi\right) - T_1}{T_2 - T_1}, \quad 0 \leqslant \xi \leqslant 1. \end{cases} \end{array}
cu_{\mathbf{t}} = [\alpha/L]^2 u_{\xi\xi}.
\tau = [\alpha/L]^2 t.
\]

\begin{verbatim}
\begin{array}{cccccccccccccccccccccccccccccccccccc
\end{verbatim}

\[
\quad 0 \leqslant \xi \leqslant 1,
u(x, t) = T_1 + (T_2 - T_1) U(x/L, \alpha^2 t/L^2).
u(\xi, \tau) = \cos(\pi \tau) \sin(\pi \xi) + 0.5 \cos(3\pi \tau) \sin(3\pi \xi).
\begin{array}{ll} (\text{УЧП}) & u_{tt} = \alpha^2 u_{xx}, & 0 < x < L, & 0 < t < \infty, \\ u(0, t) = 0, & 0 < t < \infty, \\ u(L, t) = 0, & 0 < t < \infty, \\ u(X, 0) = \sin(\pi x/L) + (0.5)\sin(3\pi x/L), \\ u_t(x, 0) = 0. & \end{array}
u(x, t) = \cos(\pi \alpha t/L)\sin(\pi x/L) + 0.5\cos(3\pi \alpha t/L)\sin(3\pi x/L).
\xi = x/L
\xi = x/L
(\text{YYII}) \qquad u_{\tau\tau} = u_{\xi\xi}, \quad 0 < \xi < 1, \quad 0 < \tau < \infty, 1
\begin{cases} u(0, \tau) = 0, \\ u(1, \tau) = 0, \\ u(\xi, 0) = \sin(\pi\xi) + 0.5\sin(3\pi\xi), \\ u_{\xi}(\xi, 0) = 0, \end{cases}
(\text{HY}) \quad
\begin{cases} u(\xi, 0) = \sin(\pi\xi) + 0.5\sin(3\pi\xi), \\ u_{\xi}(\xi, 0) = 0, \end{cases}
\quad 0 \leqslant \xi \leqslant 1,
u(\xi, \tau) = \cos(\pi \tau) \sin(\pi \xi) + 0.5 \cos(3\pi \tau) \sin(3\pi \xi).
\]

\begin{verbatim}
\right. \\ \end{array}
\end{verbatim}

\[
u(x, t) = \cos(\pi \alpha t/L)\sin(\pi x/L) + 0.5\cos(3\pi \alpha t/L)\sin(3\pi x/L).
U(x, t) = \frac{u(x, t) - T_1}{T_2 - T_1}
\xi = x/L
u_t + vu_x = 0?
\]

\begin{verbatim}
\begin{array}{lll} (\text{Y} \, \Pi) & u_t = \alpha^2 u_{xx}, & 0 < x < L, \\ (\Gamma \, \text{Y}) & \left\{ \begin{array}{ll} u \, (0, \, t) = T_1, \\ u \, (L, \, t) = 0, \\ u \, (x, \, 0) = T_2, & 0 \leqslant x \leqslant L. \end{array} \end{array}
\end{verbatim}

\[
Au_{xx} + Bu_{xy} + Cu_{yy} + Du_x + Eu_y + Fu = G
Au_{xx} + Bu_{xy} + Cu_{yy} + Du_x + Eu_y + Fu = G,
A=1
A=1
u_{\xi\eta} = \Phi(\xi, \eta, u, u_{\xi}, u_{\eta}).
Au_{xx} + Bu_{xy} + Cu_{yy} + Du_x + Eu_y + Fu = G
\xi = \xi(x, y),
\eta = \eta(x, y)
\overline{A}u_{\xi\xi} + \overline{B}u_{\xi\eta} + \overline{C}u_{\eta\eta} + \overline{D}u_{\xi} + \overline{E}u_{\eta} + \overline{F}u = \overline{G},
\overline{A} = A\xi_{x}^{2} + B\xi_{x}\xi_{y} + C\xi_{y}^{2},
\overline{B} = 2A\xi_{x}\eta_{y} + B(\xi_{x}\eta_{y} + \xi_{y}\eta_{x}) + 2C\xi_{y}\eta_{y},
\overline{C} = A\eta_{x}^{2} + B\eta_{x}\eta_{y} + C\eta_{y}^{2},
\overline{D} = A\xi_{xx} + B\xi_{xy} + C\xi_{yy} + D\xi_{x} + E\xi_{y},
\overline{E} = A\eta_{xx} + B\eta_{xy} + C\eta_{yy} + D\eta_{x} + E\eta_{y},
\overline{F} = F,
\overline{G} = G.
\]

так, чтобы в уравнении (23.3) осталась только одна частная производная
второго порядка \(u_{\xi\eta}\) (не следует пытаться искать такую замену
переменных, с помощью которой гиперболическое уравне-

1) Необходимо, чтобы это преобразование координат было локально обратимо
и дважды дифференцируемо. --- Прим. ред.

ние можно было бы привести и параболической или заливитической
канонической форме, это невозможно сделать).

Вычислим сначала частные производные

\[
\bar{A} = A\xi_x^2 + B\xi_x\xi_y + C\xi_y^2 = 0,
\]

Подставляя эти соотношения в исходное уравнение (23.3), после простых,
но громоздких вычислений получаем

(23.5)

\begin{verbatim}
Разрешив их относительно  [\xi_x/\xi_y]  и  [\eta_x/\eta_y] , получаем
\end{verbatim}

(23.7)

\[
= \frac{-B + \sqrt{B^2 - 4AC}}{2A}
[\eta_x/\eta_y] = \frac{-B - \sqrt{B^2 - 4AC}}{2A}
\begin{split} &A\left[\xi_{x}/\xi_{y}\right]^{2}+B\left[\xi_{x}/\xi_{y}\right]+C=0,\\ &A\left[\eta_{x}/\eta_{y}\right]^{2}+B\left[\eta_{x}/\eta_{y}\right]^{2}+C=0. \end{split}
\]

\pandocbounded{\includegraphics[keepaspectratio]{images/_page_170_Figure_3.jpeg}}

\[
u_{xx} - 4u_{yy} + u_x = 0.
\]

\begin{verbatim}
\begin{vmatrix} \frac{dy}{dx} = -[\xi_x/\xi_y] = \frac{B - \sqrt{B^2 - 4AC}}{2A} = -2, \\ \frac{dy}{dx} = -[\eta_x/\eta_y] = \frac{B^{\frac{1}{2}} + \sqrt{B^2 - 4AC}}{2A} = 2.
\end{verbatim}

\[
\begin{bmatrix} \xi_x/\xi_y \end{bmatrix}
y = -2x + c_1,
y = 2x + c_2.
\]

(характеристические

\[
\xi = y + 2x = c_1,
\]

уравнения).

Kaждое из квадратных уравнений для \([\xi_x/\xi_y]\) и
\([\eta_x/\eta_y]\) имеет два кория, но мы оставляем только по одному,
но так чтобы они были различны.

\[
\eta = y - 2x = c_2.
\]

РИС. 23.1. Характеристики \(\xi(x, y) = c\) и \(\eta(x, y) = c\) .

Задача була зведена до знаходження двох функцій \(\xi(x, y)\) і
\(\eta(x, y)\) , щоб відношення \([\xi_x/\xi_y]\) і \([\eta_x/\eta_y]\)
задовольняли рівняння (23.7). Дуже легко знайти такі функції 1) якщо
звернути увагу на рис. 23.1. Щоб зрозуміти, як це можна зробити,
розглянемо просте рівняння Е

\pandocbounded{\includegraphics[keepaspectratio]{images/_page_171_Figure_3.jpeg}}

го характеристики определяются из уравнений

\[
\overline{A}u_{\xi\xi} + \overline{B}u_{\xi\eta} + \overline{C}u_{\eta\eta} + \overline{D}u_{\xi} + \overline{E}u_{\eta} + \overline{F}u = \overline{G}
\]

Решая эти уравнения относительно у, получаем

\[
y^2u_{xx}-x^2u_{yy}=0
\frac{dy}{dx} = \frac{B + \sqrt{B^2 - 4AC}}{2A} = \frac{x}{y}
\frac{dy}{dx} = \frac{B - \sqrt{B^2 - 4AC}}{2A} = -\frac{x}{y}.
\]

\pandocbounded{\includegraphics[keepaspectratio]{images/_page_172_Figure_2.jpeg}}

\[
y^2 - x^2 = \text{const},
\]

1) См., например, Тихонов А. Н., Самарский А. А. Уравнения
математической физики. --- М.: Наука, 1972 и др. издания. --- Прим. ред.

Очевидно, введенные таким образом функции \(\xi\) и \(\eta\)
удовлетворяют уравнениям характеристик. Эти новые координаты изображены
на рис. 23.2. На этом мы заканчиваем обсуждение вопроса о том, как найти
новые координаты. Осталось последнее --- определить новый вид исходного
уравнения.

\[
y^2 + x^2 = \text{const}.
\]

РИС. 23.2. Характеристические координаты для уравнения
\(u_{xx} - 4u_{uu} + u_x = 0\) .

Это совсем просто: чтобы найти каноническую форму уравнения с частными
производными, подставим новые координаты \(\xi(x, y)\) и \(\eta(x,) y)\)
в уравнение

\[
\xi = y^2 - x^2
\]

,

где \(\overline{A}\) , \(\overline{B}\) , \(\overline{C}\) ,
\(\overline{D}\) , \(\overline{E}\) , \(\overline{F}\) и
\(\overline{G}\) определены формулами (23.6).

Прежде чем закончить лекцию, давайте на конкретном примере посмотрим,
как «работает» общий метод приведения к каноническому виду.

Приведение гиперболического уравнения \(y^2u_{xx}-x^2u_{yy}=0\) к
канонической форме

Рассмотрим уравнение

\[
\vec{B} = 2A\xi_x\eta_x + B(\xi_x\eta_y + \xi_y\eta_x) + 2C\xi_y\eta_y = -16x^2y^2
\]

, \(x>\) , \(0 y>0\) ,

которое является гиперболическим в первом квадранте. Найдем новые
координаты таким образом, чтобы привести исходное уравнение к
канонической форме.

ШАГ 1. Решаем уравнения характеристик

\[
\bar{D} = A\xi_{xx} + B\xi_{xy} + C\xi_{yy} + D\xi_x + E\xi_y = -2(x^2 + y^2)
\]

(напомним, что этот шаг

\[
\widetilde{E} = A\eta_{xx} + B\eta_{xy} + C\eta_{yy} + D\eta_x + E\eta_y = 2(y^2 - x^2),
\]

обеспечивает \(\overline{A} = \overline{C} = 0\) ),

\[
\tilde{F} = F = 0
\]

РИС, 23.3. Новые карактеристические координаты.

Проинтегрировав эти обыкновенные дифференциальные уравнения методом
разделения переменных, получаем два неявных соотнощения (при желании
можно было бы явно выразить у через х)

\[
\overline{G} = G = 0
\overline{A}u_{\xi\xi} + \overline{B}u_{\xi\eta} + \overline{C}u_{\eta\eta} + \overline{D}u + \overline{E}u_{\eta} + \overline{F}u = \overline{G}
u_{\xi\eta} = \frac{-(x^2 + y^2) u_{\xi} + (y^2 - x^2) u_{\eta}}{8x^2y^3}.
u_{\xi\eta} = \frac{\eta u_{\xi} - \xi u_{\eta}}{2(\xi^2 - \eta^2)}.
\alpha = \alpha (\xi, \eta) = \xi + \eta,
u_{\alpha\alpha} - u_{\beta\beta} = \frac{-\beta u_{\alpha} - \alpha u_{\beta}}{2\alpha\beta}.
\alpha = \xi + \eta = (y^2 - x^2) + (y^2 + x^2) = 2y^3,
\beta = \xi - \eta = (y^2 - x^2) - (y^2 + x^2) = -2x^2.
u_{\xi\xi} - u_{\eta\eta} = \Psi(\xi, \, \eta, \, u_{\xi}, \, u, u_{\eta}).
3u_{xx} + 7u_{xy} + 2u_{yy} = 0
u_{\xi\eta} = \Phi(\xi, \eta, u, u_1, u_2).
u_{\alpha\alpha} - u_{\beta\beta} = \Psi(\alpha, \beta, u, u_{\alpha}, u_{\beta}).
u_{xx} + 4u_{xy} = 0.
u_{tt} = e^{2} (u_{xx} + u_{yy} + u_{zz}),
\begin{array}{l} u_{\xi} = u_{\alpha}\alpha_{\xi} + u_{\beta}\beta_{\xi} = u_{\alpha} + u_{\beta}. \\ u_{\eta} = u_{\alpha}\alpha_{\eta} + u_{\beta}\beta_{\eta} = u_{\alpha} - u_{\beta}, \\ u_{\xi\eta} = u_{\alpha\alpha}\alpha_{\eta} + u_{\alpha\beta}\beta_{\eta} + u_{\beta\alpha}\alpha_{\eta} + u_{\beta\beta}\beta_{\eta} = u_{\alpha\alpha} - u_{\beta\beta}, \end{array}
\]

(НУ)

\begin{itemize}
\item
  Читач може запитати, чому необхідно класифікувати та зводити їх до
  канонічної форми диференціальних рівнянь у похідних похідних.
\item
  \begin{enumerate}
  \def\labelenumi{\alph{enumi})}
  \tightlist
  \item
    Поділ рівнянь на класи гіперболічних, параболічних та еліптичних
    відповідає поділу фізичних процесів на три основні класи: хвильовий,
    дифузійний і стаціонарний. Математичні особливості розв'язків цих
    трьох типів рівнянь є абсолютно різними.
  \end{enumerate}
\item
  \begin{enumerate}
  \def\labelenumi{\alph{enumi})}
  \setcounter{enumi}{1}
  \tightlist
  \item
    У більшості праць, присвячених розв'язанню гіперболічних задач,
    припускається, що початкове рівняння записується у канонічній формі,
    тобто у вигляді
  \end{enumerate}
\end{itemize}

Якщо у нас є рівняння і ми хочемо отримати його розв'язок, ми повинні
перетворити його у канонічну форму і використати відомі результати.

в) Для численного решения канонического гиперболического уравнения
создано много различных программ для ЭВМ.

Функція \(\Phi\) ( \(\xi\) , \(\xi \eta\) , u, \(u_{\xi}\) ,
\(u_{\eta}\) ) вводиться у VM \(\Theta\) як підпрограма, тому рівняння
спочатку має бути канонічним. Як тільки розв'язок знайдено в нових
координатах, його завжди можна перетворити назад на початкові
координати.

\section{ЗАДАЧИ}\label{ux437ux430ux434ux430ux447ux438-1}

\begin{itemize}
\tightlist
\item
  \begin{enumerate}
  \def\labelenumi{\arabic{enumi}.}
  \tightlist
  \item
    Выясните, являются ли приведенные ниже уравнения гыперболическими,
    параболическими или эллиптическими:
  \end{enumerate}
\item
  : \textbf{a)} \(u_{xx} u_{xy} = 0\) ,

  \begin{itemize}
  \tightlist
  \item
    \begin{enumerate}
    \def\labelenumi{\arabic{enumi})}
    \setcounter{enumi}{5}
    \tightlist
    \item
      \(u_{tt} = u_{xx} + u_x + hu\) ,
    \end{enumerate}
  \item
    \textbf{B)} \(u_{xx} + 3u_{yy} = \sin x\) ,
  \item
    \begin{enumerate}
    \def\labelenumi{\alph{enumi})}
    \setcounter{enumi}{17}
    \tightlist
    \item
      \(u_{xx} + u_{yy} = f(x, y)\) ,
    \end{enumerate}
  \item
    \(\mathbf{D} ) \ u_{rr} + \frac{1}{r} u_r + \frac{1}{r^2} u_{\theta \theta} = f(r, \ \theta).\)
    - 2. Получите соотношения (23.4), (23.5) и (23.6).
  \end{itemize}
\item
  \begin{enumerate}
  \def\labelenumi{\arabic{enumi}.}
  \setcounter{enumi}{2}
  \tightlist
  \item
    Убедитесь в том, что уравнение
  \end{enumerate}
\end{itemize}

\begin{verbatim}
гиперболического типа при всех x и y и найдите  x a p a \kappa m e p u c  -  m u + e c \kappa u + e c \kappa u + e c \kappa u + e c \kappa u + e c \kappa u + e c \kappa u + e c \kappa u + e c \kappa u + e c \kappa u + e c \kappa u + e c \kappa u + e c \kappa u + e c \kappa u + e c \kappa u + e c \kappa u + e c \kappa u + e c \kappa u + e c \kappa u + e c \kappa u + e e c \kappa u + e c \kappa u + e c \kappa u + e c \kappa u + e c \kappa u + e c \kappa u + e c \kappa u + e c \kappa u + e c \kappa u + e c \kappa u + e c \kappa u + e c \kappa u + e c \kappa u + e c \kappa u + e c \kappa u + e c \kappa u + e c \kappa u + e c \kappa u + e c \kappa u + e c \kap u + e c \kap
\end{verbatim}

\begin{verbatim}
pa u + e c \kappa u + e c \kappa u + e c \kappa u + e c \kappa u + e c \kappa u + e c \kappa u + e c \kappa u + e c \kappa u + e c \kappa u + e c \kappa u + e c \kappa u + e c \kappa u + e c \kappa u + e c \kappa u + e c \kappa u + e c \kappa u + e c \kappa u + e c \kappa u + e c \kappa u + e c \kappa u + e c \kappa u + e c \kappa u + e e c \kappa u + e c \kappa u + e c \kappa u + e c \kappa u + e c \kappa u + e c \kappa u + e c \kappa u + e c \kappa u + e c \kappa u + e c \kappa u + e c \kappa u + e c \kappa u + e c \kappa u + e c \kappa u + e c \kappa u + e c \kappa u + e c \kappa u + e c \kappa u + e c \kappa u + e c \kappa u + e c \kappa u + e c \kappa u + e c \kappa u + e c \kappa u + e c \kappa u + e c \kappa u + e c \kappa u + e c \kappa u + e c \kappa u + e c \kappa u + e c \kappa u + e c \kappa u + e c \kappa u + e c \kappa u + e c \kappa u + e c \kappa u + e c \kappa u + e c \kappa u + e c \kappa u + e c \kappa u + e u + e c \kappa u + e c \kappa u + e c \kappa u + e c \kappa u + e c \kappa u + e c \kappa u + e c \kappa u + e c \kappa u + e c \kappa u + e c \kappa u + e c \kappa u + e c \kappa u + e c \kappa u + e c \kappa u + e c \kappa u + e c \kappa u + e c \kappa u + e c \kappa u + e c \kappa u + e c \ kappa u + e c \kappa u + e c \kappa u + e c \kappa u + e c \kappa u + e c \kappa u + e c \kappa u + e c \kappa u + e c \kappa u + e c \kappa u + e c \kappa u + e c \kappa u + e c \kappa u + e c \kappa u + e c \kappa u + e c \kappa u + e c \kappa u + e c \kappa u + e c \kappa u + e e c \kappa u + e c \kappa u + e c \kappa u + e c \kappa u + e c \kappa u + e c \kappa u + e c \kappa u + e c \kappa u + e c \kappa u + e c \kappa u + e c \kappa u + e c \kappa u + e c \kappa u + e c \kappa u + e c \kappa u + e c \kappa u + e c \kappa u + e c \kappa u + e c \kappa u + e c \kappa u + e c \kappa u + e c \kappa u + e c \kappa u + e c \kappa u + e c \kappa u + e c \kappa u + e c \kappa u + e c \kappa u + e c \kappa u + e c \kappa u + e c \kappa u + e c \kappa u + e c \kappa u + e c \kappa u + e c \kappa u + e c \kappa u + e c \kappa u + e c \kappa u + e c \kappa u + e c \kappa u + e u + e c \kappa u + e c \kappa u + e c \kappa u + e c \kappa u + e c \kappa u + e c \kappa u + e c \kappa u + e c \kappa u + e c \kappa u + e c \kappa u + e c \kappa u + e c \kappa u + e c \kappa u + e c \kappa u + e c \kappa u + e c \kappa u + e c \kappa u + e c \kappa u + e c \kappa u + e c \ kappa u + e c \kappa u + e c \kappa u + e c \kappa u + e c \kappa u + e c \kappa u + e c \kappa u + e c \kappa u + e c \kappa u + e c \kappa u + e c \kappa u + e c \kappa u + e c \kappa u + e c \kappa u + e c  4. Приведите уравнение из предыдущей задачи к каноническому виду
\end{verbatim}

\[
u_{tt} = c^2 (u_{xx} + u_{yy}), \quad
\]

\begin{enumerate}
\def\labelenumi{\arabic{enumi}.}
\setcounter{enumi}{4}
\item
  Продолжите вадачу 4 и приведите уравнение ко второй канонической форме
\item
  Найдите характеристики уравнения
\end{enumerate}

Преобразуйте уравнение к новым координатам, решите и, вернувшись к
координатам x и y, получите решение исходной вадачи.

\section{ВОЛНОВОЕ УРАВНЕНИЕ В СВОБОДНОМ ПРОСТРАНСТВЕ (ДВУМЕРНЫЕ И
ТРЕХМЕРНЫЕ
ЗАДАЧИ)}\label{ux432ux43eux43bux43dux43eux432ux43eux435-ux443ux440ux430ux432ux43dux435ux43dux438ux435-ux432-ux441ux432ux43eux431ux43eux434ux43dux43eux43c-ux43fux440ux43eux441ux442ux440ux430ux43dux441ux442ux432ux435-ux434ux432ux443ux43cux435ux440ux43dux44bux435-ux438-ux442ux440ux435ux445ux43cux435ux440ux43dux44bux435-ux437ux430ux434ux430ux447ux438}

ЦЕЛЬ ЛЕКЦИИ: Решить задачу с начальными условиями

(УЧП)

\[
\begin{cases} -\infty < x < \infty, \\ -\infty < y < \infty, \\ -\infty < z < \infty, \end{cases}
\begin{cases} -\infty < x < \infty, \\ -\infty < y < \infty, \\ -\infty < z < \infty, \end{cases}
\begin{cases} u(x, y, z, 0) = \varphi(x, y, z), \\ u_{t}(x, y, z, 0) = \psi(x, y, z) \end{cases}
P_{2}
\begin{aligned} u_{tt} &= c^2 \Delta u, \\ (HY) & \begin{cases} u(x, y, z, 0) = 0, \\ u_t(x, y, z, 0) = \psi(x, y, z), \end{cases} \end{aligned}
P_{1}
\]

(НУ)

\[
u(x,y) = g(x,y) =
\]

використовується метод спуску, і доведено, що двовимірні розв'язки не
задовольняють принцип Гюйгенса. І нарешті, використовуючи метод спуску,
показано, що розв'язок одномірного рівняння виражається формулою
д'Аламбера, з якою ми познайомилися раніше.

У попередніх лекціях ми розглядали задачу з початковими умовами для
нескінченної вібруючої струни і показали, що її розв'язок задається
формулою д'Аламбера. Читачеві має бути зрозуміло, що в тривимірному
просторі одномірне хвильове рівняння описує плоскі хвилі. Наприклад,
звукові або електромагнітні хвилі на достатньо великих відстанях від
джерел можна вважати плоскими хвилями і, отже, описуватися одномірним
рівнянням. Існує така термінологія:

\begin{itemize}
\tightlist
\item
  \begin{enumerate}
  \def\labelenumi{\arabic{enumi})}
  \tightlist
  \item
    одномерные волны называются плоскими волнами,
  \end{enumerate}
\item
  \begin{enumerate}
  \def\labelenumi{\arabic{enumi})}
  \setcounter{enumi}{1}
  \tightlist
  \item
    двумерные волны называются цилиндрическими волнами.
  \end{enumerate}
\end{itemize}

\begin{enumerate}
\def\labelenumi{\arabic{enumi})}
\setcounter{enumi}{2}
\tightlist
\item
  трехмерные волны называются сферическими волнами.
\end{enumerate}

Другими словами, одномерное волновое уравнение может описывать либо
плоские волны в пространствах большего числа измерений, либо колебания
струны. Задача лекции --- обобщить формулу Даламбера на случай двух и
трех измерений.

\section{Волны в трехмерном
пространстве}\label{ux432ux43eux43bux43dux44b-ux432-ux442ux440ux435ux445ux43cux435ux440ux43dux43eux43c-ux43fux440ux43eux441ux442ux440ux430ux43dux441ux442ux432ux435}

Рассмотрим сферические волны в трехмерном пространстве при ваданных
начальных условиях, т. е. решим задачу с начальными условиями

(24.1)

\begin{verbatim}
где  \overline{\psi}  — *среднее* значение начального распределения  \psi  *по сфере* радиуса ct с центром в точке (x, y, z), т. е.
\end{verbatim}

(24.1)

\pandocbounded{\includegraphics[keepaspectratio]{images/_page_177_Picture_3.jpeg}}

\pandocbounded{\includegraphics[keepaspectratio]{images/_page_177_Picture_7.jpeg}}

\begin{verbatim}
где экаком  \Delta  обозначен дифференциальный оператор
\end{verbatim}

\[
u_h(r, \theta) = \sum_{n=0}^{\infty} r^n [a_n \cos(n\theta) + b_n \sin(n\theta)].
u = \frac{\partial}{\partial t} [t\overline{\varphi}].
\]

\(+ ct \cos \varphi\) ) \((ct)^2 \sin \varphi d\varphi d\theta\) .

Аргументи функції \(\psi\) проходити поверхнею сфери, якщо значення
\(\theta\) і \(\phi\) змінюються відповідно в межах \([0, 2\pi]\) і
\([0, \pi]\) (рисунок 24.1). Цю формулу можна інтерпретувати так:

См Владимиров В. С. Уравнения математической физики. --- М.: Наука, 1971
и др. издания. --- Прим ред.

дующим образом: каждая точка пространства излучает расходящуюся (со
скоростью c) сферическую волну, через t секунд точка с координатами (x,
y, z) окажется под воздействием начального Р И С

\[
u_p(r, \theta) = Ar^4 + Br^4 \cos(2\theta).
\]

. 24 1. Решение является средним значением начатьного распределения на
сфере.

возмущения, сосредоточенного на сфере раднуса ct с центром в заданной
точке (см. рис. 24.2).

Формула (24.3) дозволяє обчислити розв'язок задачі для більшості
початкових умов на комп'ютері. Можливо, читачу буде цікаво знайти
розв'язок для деяких простих функцій f.

\[
u(x, t) = \frac{1}{2c} \int_{x-ct}^{x+ct} \varphi(s) ds.
\]

РИС. 24.2. Начальное возмущение \(\psi\) распространяется из каждой
точки во все стороны.

Чтобы получить полное решение, рассмотрим вторую половину задачи

(24.4)

\[
u_p(r, \theta) = -\frac{r^4}{32} - \frac{r^4}{24} \cos(2\theta).
u_t(x, t) = \frac{1}{2} [\varphi(x + ct) + \varphi(x - ct)].
\sum_{n=0}^{\infty} [a_n \cos(n\theta) + b_n \sin(n\theta)] - \frac{1}{32} - \frac{1}{24} \cos(2\theta) = 0.
u_1(r, \theta) = \frac{1}{32} + \frac{1}{24}r^2\cos(2\theta) - \frac{r^4}{24}\cos(2\theta) - \frac{r^4}{32} = \frac{(r^4 - 1)}{32} - \frac{(r^4 - r^2)}{24}\cos(2\theta).
u = u_0 + u_1 = r \cos \theta - \frac{(r^4 - 1)}{32} - \frac{(r^4 - r^2)}{24} \cos (2\theta).
\Delta u = 0, \quad 0 < r < 1 + \frac{1}{4} \sin \theta,
(\text{УЧП}) \quad u_{tt} = c^2 \Delta u, \quad (x, y, z) \in \mathbb{R}^3,
u \left( 1 + \frac{1}{4} \sin \theta, \ \theta \right) = \cos \theta, \quad 0 \le \theta \le 2\pi.
(\text{HV}) \quad
u\left(1+\frac{1}{4}\sin\theta,\ \theta\right)=\cos\theta
u(1+\epsilon\sin\theta, \theta) = \cos\theta, \quad 0 \le \epsilon \le 1/4.
\begin{cases} u = \varphi, \\ u_t = \psi \end{cases}
\]

\pandocbounded{\includegraphics[keepaspectratio]{images/_page_350_Figure_2.jpeg}}

\[
u(1+\varepsilon\sin\theta,\ \theta)=u(1,\ \theta)+u_r(1,\ \theta)(\varepsilon\sin\theta)+u_{rr}(1,\ \theta)\frac{(\varepsilon\sin\theta)^2}{21}+\ldots
\Delta u = 0, \quad 0 < r < 1 + \varepsilon \sin \theta,
u(1, \theta) + u_r(1, \theta) (\varepsilon \sin \theta) + u_{rr}(1, \theta) \frac{(\varepsilon \sin \theta)^2}{2!} + \dots = \cos \theta.
(46.8) u = u_0 + \varepsilon u_1 + \varepsilon^2 u_2 + \dots.
P_{0}
\begin{cases} \Delta u_{0} = 0, & 0 < r < 1 \text{ (внутри круга),} \\ u_{0}(1, \theta) = \cos \theta, & u_{0}(r, \theta) = r \cos \theta, \end{cases}
\begin{cases} -\infty < x < \infty, \\ -\infty < y < \infty, \end{cases}
P_{1}
\begin{cases} \Delta u_{1} = 0, & 0 < r < 1 \text{ (внутри круга),} \\ u_{1}(1, \theta) = -\sin \theta \frac{\partial u_{0}(1, \theta)}{\partial r} = -\sin \theta \cos \theta \end{cases}
\]

(НУ)

Тепер сформульуємо головну мету наших лекцій: знайти функцію, яка дає
мінімум (або максимум) функціоналу (44.1). Стратегія пошуку буде такою
ж, як при пошуку мінімуму функції в диференціальному численні. Там ми
знайшли критичні точки функції за умовою f'(x) = 0. У варіаційному
численні все буде трохи складніше, оскільки аргумент більше не є
числовою змінною, а функцією. Однак загальний підхід залишається
незмінним: ми обчислимо так звану функціональну похідну функції \(y(x)\)
і прирівняємо її до нуля. Це нове рівняння буде аналогічним

Розглянемо розв'язок тривимірної задачі за початкових умов, які залежать
лише від двох змінних X і Y.

Сделав это, получим, что трехмерная формула

для розв'язку та опише циліндричні хвилі і, отже, дасть розв'язок
двовимірної задачі. Цей метод називається методом спуску. Після всіх
розрахунків ми отримуємо 1) (зовсім не тривіальне)

Перед нами стоїть розв'язок двовимірного хвильового рівняння, і хоча,
найімовірніше, його практичне застосування передбачає використання
комп'ютера, цікаво аналізувати його з точки зору принципу Гюйгенса. У
цьому розв'язці обидва інтеграли початкових збурень розглядаються вздовж
внутрішньої частини (тут слово «внутрішнє» відіграє вирішальну роль)
кола з радіусом cte e центр у точці (x, y). Іншими словами, якщо ми
проаналізуємо це рішення так само, як у тривимірному випадку, то
побачимо, що початкове збурення поширюється з чітко визначеним переднім
краєм і розмитим заднім краєм. Отже, принцип Гюйгенса не виконується у
просторі двох вимірів.

Нарешті, припустимо, що початкові збурення \(\phi\) і \(\phi\) залежать
лише від однієї змінної. У цьому випадку виникають плоскі хвилі, і після
застосування методу спуску отримуємо

\&lt;sup\textgreater2) См., например, Тихонов А. Н., Самарский А. А.
Уравнения математической физики. --- М.: Наука, 1972 и др. издания. ---
Прим. ред.

хорошо известную формулу Даламбера

\[
u_u + \frac{1}{4} u_t
\]

Слід зазначити, що в цьому випадку використання методу спуску вимагає
нетривіальних розрахунків. За формулою д'Аламбера початкове зміщення f
призводить до утворення гострого заднього краю, а початкова швидкість f
його розмиває. Іншими словами, одномірний випадок є незвичайним тим, що
принцип Гюйгенса виконується для початкових зсув і не застосовується для
початкових швидкостей. Можна стверджувати, що загалом принцип Гюйгенса
не виконується в одномірному випадку.

\section{ЗАМЕЧАНИЯ}\label{ux437ux430ux43cux435ux447ux430ux43dux438ux44f-1}

Метод спуску ми не викладали детально, оскільки неможливо зупинитися на
питанні, як можна зменшити розмірність інтегралів, включених у формули.
Усі ці розрахунки можна знайти в навчальній літературі, перелік якої
наведено наприкінці книги. Загальна ідея полягає в тому, що при
розв'язанні задач у низьковимірному просторі можна використати відоме
розв'язання подібної задачі у просторі вищої вимірності, а потім
спростити його, припускаючи, що початкові та граничні умови не залежать
від частини змінних. Читачеві має бути зрозуміло, що метод спуску можна
застосувати не лише до розглянутої проблеми.

\section{ЗАДАЧИ}\label{ux437ux430ux434ux430ux447ux438-2}

\begin{enumerate}
\def\labelenumi{\arabic{enumi}.}
\tightlist
\item
  Покажите, что в одномерном случае решение задачи
\end{enumerate}

\[
u_t = (1+x) u_{xx}, \\ u(x, 0) = q(x), \quad -\infty < x < \infty,
u_1 = (1 + \varepsilon_1) u_{xx}
P_{0}
\begin{aligned} & (\forall \mathsf{Y}\Pi) \quad u_{tt} = c^2 u_{xx}, \\ & (\mathsf{H} \mathsf{Y}) \quad \begin{cases} u\left(x, \ 0\right) = 0, \\ u_t\left(x, \ 0\right) = \phi\left(x\right). \end{cases} \end{aligned}
\]

\begin{verbatim}
,  (x,  y, z) \in R^3  ,  u = 0   u_t =
\end{verbatim}

\begin{enumerate}
\def\labelenumi{\arabic{enumi}.}
\setcounter{enumi}{3}
\tightlist
\item
  Решите аналогичную (см. задачу 3) задачу для двумерного уравнения
\end{enumerate}

(УЧП)

Зверніть увагу, що в перших двох задачах коефіцієнти є сталими. Читач
повинен розуміти, що параметр b має бути достатньо малим, а необережний
ряд може відрізнятися

\section{ЗАДАЧИ}\label{ux437ux430ux434ux430ux447ux438-3}

\begin{verbatim}
- 1. Подставьте разложение (46.8) в задачу (46.7) и получите последовательность задач P_0,\ P_1,\ P_2,\ \dots . 2. Покажите, что нелинейную задачу
\end{verbatim}

(НУ)

\[
u(r, \theta) = u_0(r, \theta) + \frac{1}{4}u_1(r, \theta)
\Delta u = 0,
q_{xx} + \varphi_{uu} = 0
q_{uu} + q_{vv} = 0
\]

\pandocbounded{\includegraphics[keepaspectratio]{images/_page_353_Figure_9.jpeg}}

\[
\begin{cases} u(x, 0) = x, \\ u_t(x, 0) = 0. \end{cases}
w = f(z)
\]

\pandocbounded{\includegraphics[keepaspectratio]{images/_page_355_Figure_5.jpeg}}

\begin{enumerate}
\def\labelenumi{\arabic{enumi}.}
\setcounter{enumi}{2}
\tightlist
\item
  Проиллюстрируйте и прокомментируйте решение трехмерной задачи
\end{enumerate}

(УЧП)

\pandocbounded{\includegraphics[keepaspectratio]{images/_page_355_Figure_7.jpeg}}

\[
u + iv = (x + iy)^2 = x^2 - y^2 + 2ixy
\]

\begin{verbatim}
\left\{
\end{verbatim}

(НУ)

Чи знаєте ви, як виглядає графік цієї функції з сегментом {[}0, 1{]}?
Подумай про це. Зверніть увагу, що синусоїдальне перетворення функції
\(f(x)\) визначається лише для додатних цілих чисел n.~Іншими словами,
скінченний синус і косинус перетворюють функції у чисельні
послідовності.

\pandocbounded{\includegraphics[keepaspectratio]{images/_page_185_Figure_3.jpeg}}

РИС. 25.1. Графики функции \(f(x)\) = 1 и ее преобразования.

\section{Свойства
преобразований}\label{ux441ux432ux43eux439ux441ux442ux432ux430-ux43fux440ux435ux43eux431ux440ux430ux437ux43eux432ux430ux43dux438ux439}

Прежде чем приступить к решению задач, мы должны получить некоторые
свойства этих преобразований.

\begin{verbatim}
Якщо you(x, t) — функція  \partial \theta yx  змінних, і ми виконуємо перетворення змінної x, то
\end{verbatim}

\begin{enumerate}
\def\labelenumi{\arabic{enumi}.}
\setcounter{enumi}{4}
\tightlist
\item
  Решите аналогичную (см. задачу 3) задачу для одномерных плоских волн
\end{enumerate}

(УЧП)

(Зверніть увагу, що перетворення було виконано на змінній x, і отримана
послідовність залежить лише від часу t.)

А как быть с производными? Ниже приводятся несколько полезных формул

\begin{itemize}
\tightlist
\item
  \begin{enumerate}
  \def\labelenumi{\arabic{enumi}.}
  \setcounter{enumi}{5}
  \tightlist
  \item
    Какова физическая природа того факта, что принцип Гюйгенса не
    справедлив в двумерном случае?
  \end{enumerate}
\item
  \begin{enumerate}
  \def\labelenumi{\arabic{enumi}.}
  \setcounter{enumi}{6}
  \tightlist
  \item
    Воспользуйтесь формулой Лейбница
  \end{enumerate}
\end{itemize}

и продифференцируйте интеграл

\[
w = \ln\left\{\frac{z-1}{z+1}\right\}
\]

по времени t.

\section{КОНЕЧНЫЕ ПРЕОБРАЗОВАНИЯ ФУРЬЕ (СИНУС- И
КОСИНУС-ПРЕОБРАЗОВАНИЯ)}\label{ux43aux43eux43dux435ux447ux43dux44bux435-ux43fux440ux435ux43eux431ux440ux430ux437ux43eux432ux430ux43dux438ux44f-ux444ux443ux440ux44cux435-ux441ux438ux43dux443ux441--ux438-ux43aux43eux441ux438ux43dux443ux441-ux43fux440ux435ux43eux431ux440ux430ux437ux43eux432ux430ux43dux438ux44f}

ЦЕЛЬ ЛЕКЦИИ: Ввести два новых интегральных преобразования (конечные
спиус- и косинус-преобразования):

(конечное синус-преобразование),

(конечное косинус-преобразование),

(обратное синус-преобразование),

\[
(\text{УЧП}) \quad u_{tt} = u_{xx} + \sin(\pi x), \quad 0 < x < 1, \quad 0 < t < \infty,
\]

(обратное косинус-преобразование),

и показать, как с помощью этих преобразований решать краевые задачи (в
частности, неоднородные). Ранее мы познакомились с регулярными
преобразованиями Фурье и Лапласа и увидели, как применение этих
преобразований к уравнению с частными производными сводит задачу к
решению обыкновениого дифференциального уравнения. Обычное
преобразование Фурье подразумевает, что преобразуемая переменная
изменяется от \(-\infty\) до \(+\infty\) , и, следовательно, оно
применяется для решения задач в свободном пространстве (без границ). В
этой лекции мы покажем, как решать краевые задачи (с границами),
преобразуя ограниченные переменные (с таким преобразованием мы
встречаемся впервые).

Давайте на время забудем о наших задачах, а займемся самими
преобразованнями. Дадим определение преобразований, приведем формулы
обращения и основные свойства. Для чего нужны преобразования и как ими
пользоваться, мы узнаем несколько позже. Впрочем, если говорить кратко,
то суть интегральных преобразований в разложении входящих в задачу
функций по

частотам, решении семейства задач для каждой частоты и после-

дующем суммировании всех полученных результатов.

Начнем с функции \(f(x)\), определенной на интервале {[}0, L{]}.
Конечные синус- и косинус-преобразования этой функции определяются
формулами

Конечное синус-преобразование
\(S[f] = S_n = \frac{2}{L} \int_{\sigma}^{L} f(x) \sin(n\pi x/L) dx\)
\(n = 1, 2, \dots\)
\(C[f] = C_n = \frac{2}{L} \int_{\sigma}^{L} f(x) \cos(n\pi x/L) dx\)
\(n = 0, 1, \dots\) Конечное косинус-преобразование

Читатель, наверное, заметил, что эти преобразования не отличаются от
формул для определения коэффициентов Фурье при разложении их в ряд по
синусам или косинусам. Формулы обращения для этих преобразований---это
обычные ряды Фурье по синусам и косинусам.

Обратное синус-преобразование

\[
(\text{ГУ}) \quad
\]

Обратное косинус-преобразование

Отметим, что в обратном косинус-преобразовании суммирование начинается с
n=0, а в синус-преобразовании --- с n=1.

\section{Примеры
синус-преобразования}\label{ux43fux440ux438ux43cux435ux440ux44b-ux441ux438ux43dux443ux441-ux43fux440ux435ux43eux431ux440ux430ux437ux43eux432ux430ux43dux438ux44f}

Пусть

тогда

\[
= \ln\left[\frac{z-1}{z+1}\right] + i \arg\left[\frac{z-1}{z+1}\right]
\]

График функции \(f(x)\) и ее преобразования изображены на рис. 25.1.
Обратное преобразование записывается в виде

\[
= \arg\left[\frac{x^2+y^2-1+i2y}{(x+1)^2+y^2}\right] =
\]

(действительная форма отображения w=z\^{}2 ).

Эта форма понадобится нам позже. В этом примере она интересна тем, что
показывает, как гиперболы в плоскости z переходят в координатные линии u
= const и v = const в плоскости w.

Теперь мы перейдем к изучению одного частного вида функций комплексного
переменного, которые осуществляют конформные отображения.

\section{Определение конформного
отображения}\label{ux43eux43fux440ux435ux434ux435ux43bux435ux43dux438ux435-ux43aux43eux43dux444ux43eux440ux43cux43dux43eux433ux43e-ux43eux442ux43eux431ux440ux430ux436ux435ux43dux438ux44f}

\begin{verbatim}
Відображення w=f(z) комплексної площини z у комплексну площину w називається конформним у точці  z_{\bullet}  z-площини, якщо похідна  f'(z_{\bullet}) \neq 0  . Відображення f(z) називається конформним у регіоні D, якщо  f'(z) \neq 0  в кожній точці регіону D.
\end{verbatim}

\begin{verbatim}
Наприклад, відображення  f(z) = z^2  конформно всюди, крім z = 0, оскільки  z'(z) = 2z \neq 0  для всіх  z \neq 0  . З іншого боку, відображення  e^z  конформно по всій z-площині, оскільки скрізь  f'(z) = e^z \neq 0  . Яке призначення конформних відображень? Відповідь полягає в тому, що при розв'язанні рівняння Лапласа  \phi_{xx} + \phi_{yy} = 0  в деякій області площини змінних x і y можна розглядати цю площину як площину комплексної змінної z. Розглянемо комплексне відображення w = f(z) площини z у площину w, в якій введені координати (u, v). Відображуючи w = f(z  \phi_{xy} + \phi_{uy} = 0  ), рівняння Лапласа  \varphi_{xx} + \varphi_{yy} = 0  перетворюється на нове диференціальне рівняння в похідних похідних, яке залежить від нових координат і v. To — це рівняння  \phi_{xx} + \phi_{uu} = 0  яке повернеться у рівняння  \varphi_{nn} + \varphi_{nn} = 0  . Таким чином, ми придумаємо ідею знайти конформне відображення, яке перетворює область з комплексною межею на область з простою межею (нагадаємо, що відображення  w=z^*  переносить межу першого квадранта площини z у дійсну вісь площини w).
\end{verbatim}

Приведем несколько примеров, в которых будет показано, как можно решать
уравнения с частными производными с помощью конформных отображений.

Уравнение Лапласа в верхней полуплоскости.

Предположим, что нам необходимо решить следующую задачу Дирихле в
верхней полувлоскости (рис. 47.3):

(47.1)

\[
\begin{cases} u(x, 0) = 1, \\ u_t(x, 0) = 0, \end{cases}
x^{2} + y^{2} = 1,
\]

\begin{verbatim}
\begin{aligned} (Y & \Pi) & \varphi_{xx} + \varphi_{yy} = 0, & -\infty < x < \infty, & 0 < y < \infty, \\ (47.1) & \varphi(x, 0) = \begin{cases} 0, & |x| > 1, \\ 1, & |x| \le 1. \end{cases}
\end{verbatim}

Решение задачи выполним за три следующих шага:

ШАГ І (Определение подходящего преобразования.)

Поскольку переменная x меняется от 0 до 1, мы воспользуемся конечным
преобразованием, а именно синус-преобразованием; позже станет ясно
почему. Мы могли бы решить эту задачу с помощью преобразования Лапласа
по переменной t (по трудоемкости этот метод не отличается от конечного
синус-преобразования).

ШАГ 2 (Выполнение преобразования).

Для удобства будем использовать обозначение \(S_n\)(t) = S{[}u{]} .
Применим синус-преобразование к исходному уравнению

\[
w = \ln\left\{\frac{z-1}{z+1}\right\}
\]

и воспользуемся тождеством для синус-преобразования, получим

\[
\begin{split} S\left[u_{t}\right] &= \frac{dS\left[u\right]}{dt}, \quad S\left[u_{tt}\right] = \frac{d^{2}S\left[u\right]}{dt^{2}}, \\ S\left[u_{xx}\right] &= -\left[n\pi/L\right]^{2}S\left[u\right] + \frac{2n\pi}{L^{2}}\left[u\left(0,\ t\right) + (-1)^{n+2}u\left(L,\ t\right)\right], \\ C\left[u_{xx}\right] &= -\left[n\pi/L\right]^{2}C\left[u\right] - \frac{2}{L}\left[u_{x}\left(0,\ t\right) + (-1)^{n+1}u_{x}\left(L,\ t\right)\right]. \end{split}
\begin{aligned} & (\text{УЧП}) & \phi_{xx} + \phi_{yy} = 0, & \text{внутри } D, \\ & (\text{ГУ}) & \begin{cases} \varphi(x, y) = 1 & \text{на } x^2 + y^2 = 1, \\ \varphi(x, y) = 2 & \text{на } (x - 1)^2 + y^2 = 9. \end{cases} \end{aligned}
\]

M A

\pandocbounded{\includegraphics[keepaspectratio]{images/_page_357_Figure_11.jpeg}}

THBLK1213X M A

\pandocbounded{\includegraphics[keepaspectratio]{images/_page_358_Figure_2.jpeg}}

THBLK1214X

\begin{verbatim}
\begin{cases} (\text{УЧП}) & \varphi_{uu} + \varphi_{vv} = 0, \quad -\infty < u < \infty, \quad 0 < v < \pi, \\ (\text{47.2}) & \begin{cases} \varphi(u, 0) = 0, \\ \varphi(u, \pi) = 1. \end{cases}
\end{verbatim}

где

\[
w = u + iv = \ln\left[\frac{z-1}{z+1}\right] =
u = \gamma [(x-s)(x-t)-\gamma^2],
\begin{cases} u(0, t) = 0, \\ u(1, t) = 0, \end{cases}
\gamma = 2t/[(x-t)^2 + y^2].
= \ln\left[\frac{z-1}{z+1}\right] + i \arg\left[\frac{z-1}{z+1}\right]
\]

Если теперь преобразовать начальные условия краевой задачи, то получим
начальные условия для обыкновенного дифференциального уравнения

\[
v = \arg\left[\frac{z-1}{z+1}\right] = \arg\left[\frac{x+iy-1}{x+1+iy}\right] =
= \arg\left[\frac{x^2+y^2-1+i2y}{(x+1)^2+y^2}\right] =
= \arg\left[\frac{2y}{x^2+y^2-1}\right]
\varphi(r) = a \ln r + b, \quad r = u^2 + v^2.
\begin{cases} u(x, 0) = 1, \\ u_t(x, 0) = 0, \end{cases}
\varphi(x, y) = 0.57 \ln(u^2 + v^2) + 1,
\varphi(x, y) = \frac{1}{\pi} \arctan\left[\frac{2y}{x^2 + y^2 - 1}\right].
\]

Итак, решим теперь новое семейство задач Кощи:

(ОДУ)

\[
x^{2} + y^{2} = 1,
w = \ln \left[ \frac{z-1}{z+1} \right].
(x-1)^{2} + y^{2} = 9,
\begin{cases} 1, & n = 1, \\ 0, & n = 2, 3, \dots, \end{cases}
\begin{aligned} & (\text{УЧП}) & \phi_{xx} + \phi_{yy} = 0, & \text{внутри } D, \\ & (\text{ГУ}) & \begin{cases} \varphi(x, y) = 1 & \text{на } x^2 + y^2 = 1, \\ \varphi(x, y) = 2 & \text{на } (x - 1)^2 + y^2 = 9. \end{cases} \end{aligned}
\begin{split} \frac{d^2S_n\left(t\right)}{dt^2} &= -\left(\pi n\right)^2S_n\left(t\right) + 2n\pi\left[u\left(0,\ t\right) + \left(-1\right)^{n+1}u\left(1,\ t\right)\right] + D_n\left(t\right) = \\ &= -\left(n\pi\right)^2S_n\left(t\right) + D_n\left(t\right), \end{split}
\right.
\]

M A

\pandocbounded{\includegraphics[keepaspectratio]{images/_page_359_Figure_10.jpeg}}

THBLK1228X

\[
w = 2t \left[ \frac{z-s}{z-t} \right],
u = \gamma [(x-s)(x-t)-\gamma^2],
v = \gamma [y(x-t)],
\]

\pandocbounded{\includegraphics[keepaspectratio]{images/_page_362_Figure_2.jpeg}}

\[
\gamma = 2t/[(x-t)^2 + y^2].
\]

\pandocbounded{\includegraphics[keepaspectratio]{images/_page_362_Figure_5.jpeg}}

\begin{verbatim}
\begin{aligned} \varphi_{uu} + \varphi_{vv} &= 0, & \text{в кольче,} \\ (\Gamma Y) & & & & & & & \\ (\Gamma Y) & & & & & & \\ (\Gamma Y) & & & & & \\ & & & & & \\ (\Gamma Y) & & & & & \\ & & & & \\ & & & & \\ & & & & \\ & & & & \\ & & & \\ & & & \\ & & & \\ & & & \\ & & & \\ & & & \\ & & \\ & & \\ & & \\ & & \\ & & \\ & & \\ & & \\ & & \\ & & \\ & & \\ & & \\ & & \\ & & \\ & & \\ & & \\ & & \\ & & \\ & & \\ & & \\ & & \\ & & \\ & & \\ & & \\ & & \\ & & \\ & & \\ & & \\ & & \\ & & \\ & & \\ & & \\ & & \\ & & \\ & & \\ & & \\ & & \\ & & \\ & & \\ & & \\ & & \\ & & \\ & & \\ & & \\ & & \\ & & \\ & & \\ & & \\ & & \\ & & \\ & & \\ & & \\ & & \\ & & \\ & & \\ & & \\ & & \\ & & \\ & & \\ & & \\ & & \\ & & \\ & & \\ & & \\ & & \\ & & \\ & & \\ & & \\ & & \\ & & \\ & & \\ & & \\ & & \\ & & \\ & & \\ & & \\ & & \\ & & \\ & & \\ & & \\ & & \\ & & \\ & & \\ & & \\ & & \\ & & \\ & & \\ & & \\ & & \\ & & \\ & & \\ & & \\ & & \\ & & \\ & & \\ & & \\ & & \\ & & \\ & & \\ & & \\ & & \\ & & \\ & & \\ & & \\ & & \\ & & \\ & & \\ & & \\ & & \\ & & \\ & & \\ & & \\ & & \\ & & \\ & & \\ & & \\ & & \\ & & \\ & & \\ & & \\ & & \\ & & \\ & & \\ & & \\ & & \\ & & \\ & & \\ & & \\ & & \\ & & \\ & & \\ & & \\ & & \\ & & \\ & & \\ & & \\ & & \\ & & \\ & & \\ & & \\ & & \\ & & \\ & & \\ & & \\ & & \\ & & \\ & & \\ & & \\ & & \\ & & \\ & & \\ & & \\ & & \\ & & \\ & & \\ & & \\ & & \\ & & \\ & & \\ & & \\ & & \\ & & \\ & & \\ & & \\ & & \\ & & \\ & & \\ & & \\ & & \\ & & \\ & & \\ & & \\ & & \\ & & \\ & & \\ & & \\ & & \\ & & \\ & & \\ & & \\ & & \\ & & \\ & & \\ & & \\ & & \\ & & \\ & & \\ & & \\ & & \\ & & \\ & & \\ & & \\ & & \\ & & \\ & & \\ & & \\ & & \\ & & \\ & & \\ & & \\ & & \\ & & \\ & & \\ & & \\ & & \\ & & \\ & & \\ & & \\ & & \\ & & \\ & & \\ & & \\ & & \\ & & \\ & & \\ & & \\ & & \\ & & \\ & & \\ & & \\ & & \\ & & \\ & & \\ & & \\ & & \\ & & \\ & & \\ & & \\ & & \\ & & \\ & & \\ & & \\ & & \\ & & \\ & & \\ & & \\ & & \\ & & \\ & & \\ & & \\ & & \\ & & \\ & & \\ & & \\ & & \\ & & \\ & & \\ & & \\ & & \\ & & \\ & & \\ & & \\ & & \\ & & \\ & & \\ & & \\ & & \\ & & \\ & & \\ & & \\ & & \\ & & \\ & & \\ & & \\ & & \\ & & \\ & & \\ & & \\ & & \\ & & \\ & & \\ & & \\ & & \\ & & \\ & & \\ & & \\ & & \\ & & \\ & & \\ & & \\ & & \\ & & \\ & & \\ & & \\ & & \\ & & \\ & & \\ & & \\ & & \\ & & \\ & & \\ & & \\ & & \\ & & \\ & & \\ & & \\ & & \\ & & \\ & & \\ & & \\ & & \\ & & \\ & & \\ & &
\end{verbatim}

\pandocbounded{\includegraphics[keepaspectratio]{images/_page_362_Figure_7.jpeg}}

\begin{verbatim}
\begin{cases} 4/n\pi, & n = 1, 3, \dots, \\ 0, & n = 2, 4, \dots, \end{cases} \end{aligned}
\end{verbatim}

\[
\varphi(r) = a \ln r + b, \quad r = u^2 + v^2.
\varphi(u, v) = 0.57 \ln(u^2 + v^2) + 1
H(x-a) =
\varphi(x, y) = 0.57 \ln(u^2 + v^2) + 1,
\begin{cases} 0, x < a, \\ 1, x \geqslant a, \end{cases}
w = \ln \left[ \frac{z-1}{z+1} \right].
\]

В результате получаем

\[
z = re^{i\theta}
H(x-a) =
\begin{cases} 1, & n = 1, \\ 0, & n = 2, 3, \dots, \end{cases}
\begin{cases} 0, & x < a, \\ 1, & x \geqslant a, \end{cases}
\]

\begin{verbatim}
\begin{aligned} (\text{Y} \forall \Pi) & \qquad \phi_{\text{vx}} + \phi_{yy} = 0, \quad 0 < x < \infty, \quad 0 < y < \infty, \\ \left\{ \begin{array}{l} \phi(x, \ 0) = \begin{cases} 1, & 0 < x < 1, \\ 0, & 1 \leqslant x < \infty, \end{cases} \\ \phi(0, \ y) = \begin{cases} 1, & 0 < y < 1, \\ 0, & 1 \leqslant y < \infty. \end{cases} \end{aligned} \end{array}
\end{verbatim}

\[
f(x) = \mathcal{F}^{-1}[F] = \frac{1}{\sqrt{2\pi}} \int_{-\infty}^{+\infty} F(\omega) e^{-i\omega x} d\omega \qquad F(\omega) = \mathcal{F}[f] = \frac{1}{\sqrt{2\pi}} \int_{-\infty}^{+\infty} f(x) e^{-i\omega x} dx
\begin{cases} 1, |x| < a \\ 0, |x| > a \end{cases}
\]

\begin{verbatim}
\begin{cases} S_n(0) = \begin{cases} 4/n\pi, & n = 1, 3, \dots, \\ 0, & n = 2, 4, \dots, \end{cases}
\end{verbatim}

\section{ЗАМЕЧ
АНИЯ}\label{ux437ux430ux43cux435ux447-ux430ux43dux438ux44f}

\begin{enumerate}
\def\labelenumi{\arabic{enumi}.}
\tightlist
\item
  Для решения задачи методом конечного синус- или косинуспреобразования
  краевые условия при x=0 и x=L должны иметь вид
\end{enumerate}

\[
\begin{cases} 1, |x| < 1 \\ 0, |x| > 1 \end{cases}
\begin{cases} \frac{dS_n(0)}{dt} = 0. \end{cases}
\]

\(f(x)\) * \(g(x)\)

\[
(1+x^2)^{-1}
\]

,

где

\[
\frac{a}{x^2+a^2}
\frac{2ax}{(x^2+a^2)^2}
\]

\pandocbounded{\includegraphics[keepaspectratio]{images/_page_189_Picture_10.jpeg}}

\begin{enumerate}
\def\labelenumi{\arabic{enumi}.}
\setcounter{enumi}{17}
\tightlist
\item
\end{enumerate}

\[
\begin{cases} 1-|x|, |x| < 1 \\ 0, |x| > 1 \end{cases}
\]

\begin{enumerate}
\def\labelenumi{\arabic{enumi}.}
\setcounter{enumi}{8}
\tightlist
\item
\end{enumerate}

\[
\begin{cases} u(0, t) = f(t) \\ u(L, t) = g(t) \end{cases}
\frac{2}{\pi} \frac{\sin a\omega}{\omega}
\frac{2}{\pi} \frac{\sin a\omega}{(\omega^2+a^2)^2}
\]

\(f(x)\) * \(g(x)\)

\begin{enumerate}
\def\labelenumi{\arabic{enumi}.}
\setcounter{enumi}{1}
\tightlist
\item
  Решите в общем виде задачу
\end{enumerate}

\[
\frac{\pi}{2} e^{-a+\omega}
\frac{\pi}{2} e^{-a+\omega}
\frac{2ax}{(x^2+a^2)^2}
\frac{\cos(ax), |x| > \pi/2a}{0, |x| > \pi/2a}
\]

\section{ПРИНЦИП СУПЕРПОЗИЦИИ--- ОСНОВА ТЕОРИИ ЛИНЕЙНЫХ
СИСТЕМ}\label{ux43fux440ux438ux43dux446ux438ux43f-ux441ux443ux43fux435ux440ux43fux43eux437ux438ux446ux438ux438-ux43eux441ux43dux43eux432ux430-ux442ux435ux43eux440ux438ux438-ux43bux438ux43dux435ux439ux43dux44bux445-ux441ux438ux441ux442ux435ux43c}

МЕТА ЛЕКЦІЇ: Ознайомити читача з принципом суперпозиції та показати, як
його можна використати для розкладання початкової задачі на підзадачі,
потім розв'язання всіх підзадач одночасно і, підсумувавши отримані
результати, знаходження рішення початкової задачі. Ми також хочемо
показати, що два основні методи розв'язання лінійних рівнянь --- метод
розділення змінних і метод інтегрального перетворення --- використовують
принцип суперпозиції.

Припустимо, що інженеру потрібно знайти відгук лінійної системи на
вхідний ефект f.~Загальний підхід до розв'язання задачі такий:

\begin{verbatim}
- 1) разложить f на элементарные воздействия  f = \sum f_k  ;
\end{verbatim}

\begin{itemize}
\tightlist
\item
  \begin{enumerate}
  \def\labelenumi{\arabic{enumi})}
  \setcounter{enumi}{1}
  \tightlist
  \item
    найти отклик системы \(u_k\) на воздействие \(f_k\) ;
  \end{enumerate}
\end{itemize}

\begin{verbatim}
3) сложить отклики  u = \sum_{i=1}^{n} u_{i}  .
\end{verbatim}

Если система линейная, то полученная сумма и является откликом системы
на входное воздействие f.~В этом состоит принцип суперпоэиции (см. рнс.
26.1).

РИС, 26.1. Основная идея принципа суперпоэнции; L---дифференциальный
оператор; f---вход; u---выход.

Эту основную идею можно использовать при решении сложных смешанных
задач, разбивая их на простые подзадачи, ре-

шая каждую из них и затем суммируя полученные результаты (конечно,
дифференциальное уравнение и граничные условия должны быть линейны).

\section{Разложение смешанной задачи на две более
простые}\label{ux440ux430ux437ux43bux43eux436ux435ux43dux438ux435-ux441ux43cux435ux448ux430ux43dux43dux43eux439-ux437ux430ux434ux430ux447ux438-ux43dux430-ux434ux432ux435-ux431ux43eux43bux435ux435-ux43fux440ux43eux441ux442ux44bux435}

Предположим, нам необходимо решить линейную задачу (обозначим ее P)

(УЧП)

(Р) (ГУ)

\[
\cos(ax)
\begin{cases} u(0, t) = 0, \\ u(1, t) = 0, \\ u(x, 0) = \sin(2\pi x), \quad 0 \le x \le 1. \end{cases}
\frac{\pi}{2} [\delta(\omega+a) + \delta(\omega-a)]
F(\omega) = \frac{?}{\pi} \int_{0}^{\infty} f(x) \sin(\omega x) dx
\begin{cases} u(0, t) = 0, \\ u(1, t) = 0, \\ u(1, t) = 0, \end{cases}
\]

f'\,'(0)

\[
\frac{1}{\alpha} F\left(\frac{\omega}{\alpha}\right)
\frac{2\omega}{\pi (\alpha^2 + \omega^2)}
[2/\pi\omega]^{1/2}
\frac{2}{\pi\omega} [1-\cos(\omega a)]
e^{-1}a + \omega
\frac{1}{2}e^{-\alpha}\sin\omega
\begin{cases} u_t = u_{xx} + f(x, t), \\ u(0, t) = 0, \\ u(1, t) = 0, \\ u(x, 0) = \varphi(x) \end{cases}
\cos(ax)
f(x) = \int_{0}^{\infty} F(\omega) \cos(\omega x) d\omega \qquad \qquad F(x) = \frac{2}{\pi} \int_{0}^{\infty} f(x) \cos(\omega x) dx
\begin{cases} u_t = u_{xx} + f(x, t), \\ u(0, t) = 0, \\ u(1, t) = 0, \\ u(x, 0) = 0 \end{cases}
S_n = \frac{2}{\pi} \int_0^{\pi} f(x) \sin(nx) dx.
y = \frac{\pi}{b - a} (x - a).
\begin{cases} u_t = u_{xx}, \\ u(0, t) = 0, \\ u(1, t) = 0, \\ u(x, 0) = \varphi(x). \end{cases}
f(x) = \sum_{n=1}^{\infty} S_n \sin(nx).
f(x) = \sum_{n=1}^{\infty} S_n \sin(nx)
\]

f'\,'(0)

\[
S_n = \frac{2}{\pi} \int_{0}^{\pi} f(x) \sin(nx) dx
0 \le x \le \pi
n = 1, 2, ...
\]

f'\,'(x)

\[
\frac{\pi}{2} e^{ax}
C_n = \frac{2}{\pi} \int_0^{\pi} f(x) \cos(nx) dx.
y = \frac{\pi}{b-a}(x-a).
f(x) = \frac{C_0}{2} + \sum_{n=1}^{\infty} C_n \cos(nx) \qquad C_n = \frac{2}{\pi} \int_{0}^{\pi} f(x) \cos nx \, dx
\]

f'\,'(x)

\[
\begin{array}{ll} u\left(x, \ 0\right) = \sin{(\pi x)}, \\ u_{t}\left(x, \ 0\right) = 0, \end{array}
\]

\begin{enumerate}
\def\labelenumi{\alph{enumi})}
\tightlist
\item
\end{enumerate}

\[
\begin{split} A_n(t) &= 2\int_0^1 u_t \sin(n\pi x) \, dx = \frac{d}{dt} \bigg[ 2\int_0^1 u(x, t) \sin(n\pi x) \, dx \bigg] = \frac{dU_n(t)}{dt} \,, \\ B_n(t) &= 2\int_0^1 u_{xx} \sin(n\pi x) \, dx = -(n\pi)^2 \, U_n(t) \,+ \\ &\quad + 2n\pi \, \big[ u(0, t) + (-1)^{n+1} \, u(1, t) \big], \end{split}
\]

\begin{enumerate}
\def\labelenumi{\arabic{enumi}.}
\setcounter{enumi}{4}
\tightlist
\item
  \(\cos(mx), m=1, 2, \dots\)
\end{enumerate}

\begin{verbatim}
\mathbf{\Pi}  для яких рівнянь функція  u_1 + u_2  також буде розв'язком? Які висновки ви можете зробити зі своїх відповідей?
\end{verbatim}

\begin{enumerate}
\def\labelenumi{\arabic{enumi}.}
\setcounter{enumi}{3}
\item
  Найдите четыре смещанные задачи, сумма решений которых дает решение
  следующей задачи:
\item
  Решите задачу Коши
\end{enumerate}

(OAY)

\[
\begin{cases} u(0, t) = 0, \\ u(1, t) = 0 \end{cases}
\]

(HY) \(U_n(0) = 0.\) Можете ли вы проверить полученное решение?

\begin{enumerate}
\def\labelenumi{\arabic{enumi}.}
\setcounter{enumi}{5}
\tightlist
\item
  Предположим, что функции \(u_1\) и \(u_2\) удовлетворяют линейным
  однородным граничным условиям
\end{enumerate}

\begin{verbatim}
\begin{cases}  u_x(1, t) + h_2 u(1, t) = 0.  Будет ли функция   u_i + u_i   удовлетворять этим условиям? , .
\end{verbatim}

Лекция 27

\section{УРАВНЕНИЯ ПЕРВОГО ПОРЯДКА (МЕТОД
ХАРАКТЕРИСТИК)}\label{ux443ux440ux430ux432ux43dux435ux43dux438ux44f-ux43fux435ux440ux432ux43eux433ux43e-ux43fux43eux440ux44fux434ux43aux430-ux43cux435ux442ux43eux434-ux445ux430ux440ux430ux43aux442ux435ux440ux438ux441ux442ux438ux43a}

ЦЕЛЬ ЛЕКЦИИ: Ввести понятие уравнения с частными производными первого
порядка (до сих пор мы рассматривали уравнения второго порядка) и
познакомиться с важным методом решения задач с начальными
условиями---методом характеристик. Задача, которой мы теперь займемся,
имеет вид

(УЧП)

\begin{verbatim}
Сразу же отметим, что впервые мы будем решать задачу с переменными коэффициентами. Оказывается, если от координат (x, t) перейти к новым, характеристическим координатам   (s, \tau)  , то наше УЧП превратится в обыкновенное дифференциальное уравнение. Затем можно решить обыкновенное дифференциальное уравнение, т. е. найти   u(s, \tau)  , а на последнем шаге выразить s и t через x и t и получить u(x, t).
\end{verbatim}

Напомним читателю, что, когда мы решали уравнение дифф

\begin{verbatim}
константа   \alpha^2   играла роль коэффициента диффузии, а v—была скоростью среды. Очевидно, что если   \alpha=0   (т. е. диффузии нет), то решение будет смещаться как целое вдоль оси x со скоростью v (у нас осталась только конвекция). Другими словами, если начальное условие   u(x,0)=\varphi(x)  , то соответствующее решение уравнения
\end{verbatim}

\begin{verbatim}
будет иметь вид   u(x, t) = \varphi(x-vt)  .
\end{verbatim}

Это дает нам основания думать, что решение уравнения первого порядка

переносится вдоль потока со скоростью

Конечно, если a и b---константы, то решение представляет собой бегущую с
постоянной скоростью волну. Если же \(a(x, t)\) и \(b(x, t)\)

\begin{verbatim}
РИС. 27.1. Начальное значение в точке x оказывает влияние на решение только вдоль характеристики в плоскости переменных (x,t). Параметр s изменяется вдоль характеристик от нуля до бесконечности, начиная с линии начального условия;   \tau   постоянна вдоль каждой характеристики;   \tau   изменяется вдоль линии начального условия.
\end{verbatim}

зависят от x и t, то скорость потока изменяется как вдоль потока, так и
во времени (читатель сможет увидеть, что начальная кривая сильно
искажается). Здесь так много аналогий с конвекцией!

Вернемся к нашей основной задаче:

(УЧП)

\begin{verbatim}
(НУ)   u(x, 0) = \varphi(x), -\infty < x < \infty.
\end{verbatim}

Решение этого линейного уравнения первого порядка основывается на
следующем физическом факте: начальное условие в некоторой точке x
переносится в tx-плоскости вдоль линии, которая

называется характеристи-кой (см. рис. 27.1)

Этим наше уразнение отличается от других уравнений (таких, как уравнение
теплопроводности \(u_t = u_{xx} )\),для которых начальное значение в
точке x оказывает влияние на решение во всех точках пространства и во
все моменты времени. Если читатель помнит, то начальное сме- Р И С

\begin{verbatim}
. 27.2. Начальное возмущение   u\left(x,0\right)   в точке порождает две волны. Начальное возмущение струны согласно волновому уравнению распространяется вдоль двух характеристик.
\end{verbatim}

\begin{verbatim}
щение скрипичной струны в точке x оказывает влияние на решение вдоль   \partial syx   линий в tx-плоскости (соответствующих двум бегущим волнам) (см. рис. 27.2).
\end{verbatim}

\begin{verbatim}
Естественно, напрашивается идея ввести две новые координаты s и   \tau   (вместо x и t) так, чтобы
\end{verbatim}

\begin{itemize}
\item
  s изменялась вдоль характеристических кривых,
\item
  т изменялась вдоль начальной кривой (желательно вдоль линии t=0).
\end{itemize}

Рассмотрим сначала новую координату s. Если s выбрать так, что она
удовлетворяет приведенному выше условию, то уравнение

превратится в обыкновенное дифференциальное уравнение К о н

\[
0 \le x \le \pi
\]

\begin{verbatim}
ечно, остается открытым вопрос, как найти эти характеристики. Ответ прост: мы выберем характеристики   \{[x(s), t(s)]: 0 < s < \infty\}   так, чтобы они удовлетворяли системе
\end{verbatim}

f'\,'(x)

\begin{verbatim}
-\infty < y < +\infty -\infty < z < +\infty \begin{array}{llll}
 \begin{array}{ll} (\mathrm{Y} + \Pi) & u_{tt} = u_{xx} + \sin{(3\pi x)}, & 0 < x < 1, & 0 < t < \infty, \\ (\mathrm{\Gamma} \mathrm{Y}) & \left\{ \begin{array}{ll} u\left(0, \ t\right) = 0, \\ u\left(1, \ t\right) = 0, \\ \end{array}
\end{verbatim}

\[
x(s) = s + c_1, \quad t(s) = s + c_2.
\]

Для каждой из подзадач выберите метод решения по своему BKYCY,

\begin{enumerate}
\def\labelenumi{\arabic{enumi}.}
\setcounter{enumi}{2}
\tightlist
\item
  Предположим, что \(u_1\) и \(u_2\) --- решения следующих уравнений:
\end{enumerate}

;

\[
b)  u_t = u_{xx} + e^t  ;
\]

\begin{enumerate}
\def\labelenumi{\Alph{enumi})}
\setcounter{enumi}{1}
\tightlist
\item
\end{enumerate}

;

\[
F)  u_t = u_{xx} + u^2  .
\frac{du}{ds} + 2u = 0, \quad 0 < s < \infty,
\]

\begin{verbatim}
\begin{aligned} u_t &= u_{xx} + f(x, t), & 0 < x < 1, & 0 < t < \infty, \\ & \begin{cases} u(0, t) = g_1(t), \\ u(1, t) = g_2(t), \end{cases} & 0 < t < \infty, \\ & u(x, 0) = \varphi(x), & 0 \le x \le 1. \end{aligned} \end{array} \end{array} \end{cases}
\end{verbatim}

x

\begin{verbatim}
| _                                |,   -\infty < x < \infty  ,   0 < t < \infty  ,
\end{verbatim}

\[
s=t
(HV)   u(x, 0) = \varphi(x), \quad -\infty < x < \infty.
u(x, t) = \sin(x-t)e^{-xt}
\Delta u = u_{xx} + u_{yy},   -\infty < x < \infty  ,   0 < t < \infty
(\forall \Pi) \quad a(x, t) u_x + b(x, t) u_t + c(x, t) u = 0, \quad -\infty < x < \infty, \\ 0 < t < \infty,
\Delta u = u_{rr} + \frac{1}{r}u_r + \frac{1}{r^2}u_{\theta\theta}
u(x, 0) = f(x), \quad -\infty < x < \infty.
\Delta u = u_{xx} + u_{yy} + u_{zz}
\frac{dx}{ds} = a(x, t), \qquad \frac{dt}{ds} = b(x, t).
\Delta u = u_{rr} + \frac{1}{r}u_r + \frac{1}{r^2}u_{\theta\theta} + u_{zz}.
x = x (s, \tau),
\]

M A

\pandocbounded{\includegraphics[keepaspectratio]{images/_page_196_Figure_10.jpeg}}

THBLK1368X

\[
\Delta u = u_{rr} + \frac{2}{r}u_r + \frac{1}{r^2}u_{\theta\theta} + \frac{\cot\theta}{r^2}
(\text{УЧП})
\]

M A

\pandocbounded{\includegraphics[keepaspectratio]{images/_page_197_Figure_5.jpeg}}

THBLK1370X

\[
\Delta u = u_{rr} + \frac{2}{r}u_r + \frac{1}{r^2}u_{\theta\theta} + \frac{\cot\theta}{r^2}u_0 + \frac{1}{r^2\sin\theta}u_{\phi\phi}
\]

M A

\pandocbounded{\includegraphics[keepaspectratio]{images/_page_370_Figure_1.jpeg}}

THBLK1372X M A

\pandocbounded{\includegraphics[keepaspectratio]{images/_page_370_Picture_2.jpeg}}

THBLK1373X M A

\pandocbounded{\includegraphics[keepaspectratio]{images/_page_370_Figure_3.jpeg}}

THBLK1374X

\[
\begin{cases} \frac{dx}{ds} = a(x, t), \\ \frac{dt}{ds} = b(x, t). \end{cases}
u(0) = F(\tau).
u(s, \tau) = F(\tau) e^{-s} \ell^2
s=t
u(x, t) = F(xe^{-t}) e^{-t^2/2}
\]

Уравнение одномерной диффузии Уравнение конвективной диффузии

Уравнение теплопроводности с поглощением Уравнение теплопроводности с
источником

\section{Диагонализация
матрицы}\label{ux434ux438ux430ux433ux43eux43dux430ux43bux438ux437ux430ux446ux438ux44f-ux43cux430ux442ux440ux438ux446ux44b}

\begin{verbatim}
Пусть A — квадратная матрина порядка n и пусть все ее собственные значения  \lambda_1  ,  \lambda_2  ...  \lambda_n  различны. Составим матрицу P так, чтобы у нее *столбцами* были координаты собственных векторов (в k-м столбце стоит собственный вектор  X_k  , соответствующий собственному значению  \lambda_k  ), т. е.
\end{verbatim}

\[
\Lambda =
\begin{bmatrix} X_1 & X_{\tilde{2}} & \dots & X_n \\ X_{1} & X_{2} & \dots & X_n \end{bmatrix}
\begin{bmatrix} 4 & 0 \\ 0 & -4 \end{bmatrix}
\]

.

\[
\begin{bmatrix} x_{11} & x_{1\tilde{2}} & \dots & x_{1n} \\ x_{21} & x_{22} & \dots & x_{2n} \\ \vdots & \vdots & \ddots & \vdots \\ x_{n1} & x_{n\tilde{3}} & \dots & x_{nn} \end{bmatrix}
\frac{\partial u_1}{\partial t} + 8 \frac{\partial u_2}{\partial x} = 0
\]

Тогда матрица Л, определяемая соотношением

где \(P^{-1}\) --- обратная к матрице P, будет диагональной матрицей
вида

\[
\begin{bmatrix} \frac{\partial u_1}{\partial t} \\ \frac{\partial u_2}{\partial t} \end{bmatrix}
\]

\begin{itemize}
\tightlist
\item
\end{itemize}

\[
\begin{bmatrix} \lambda_1 & & & \\ & \lambda_2 & 0 & \\ & & \ddots & \\ & & & \lambda_n \end{bmatrix}
\begin{bmatrix} 0 & 8 \\ 2 & 0 \end{bmatrix}
\]

Например, у матрицы

\[
\begin{bmatrix} \frac{\partial u_1}{\partial x} \\ \frac{\partial u_2}{\partial x} \end{bmatrix}
=
\begin{bmatrix} 0 \\ 0 \end{bmatrix}
\]

in X\_2 =

.

Значит, произведение

\[
u_t + Au_s = 0
A =
\begin{bmatrix} 0 & 8 \\ 2 & 0 \end{bmatrix}
, \quad u_1 =
\]

равно диагональной матрице

\[
\begin{bmatrix} \frac{\partial u_1}{\partial t} \\ \frac{\partial u_2}{\partial t} \end{bmatrix}
, \quad u_x =
\]

(Читатель может проверить это самостоятельно.)

Теперь можно решить простую систему двух уравнений с частными
производными (с соответствующими начальными условиями).

\section{\texorpdfstring{Решение линейной системы \(u_t\) + Au\_x = 0
.}{Решение линейной системы u\_t + Au\_x = 0 .}}\label{ux440ux435ux448ux435ux43dux438ux435-ux43bux438ux43dux435ux439ux43dux43eux439-ux441ux438ux441ux442ux435ux43cux44b-u_t-au_x-0-.}

Рассмотрим задачу Коши для системы, содержащей два уравнения и два
начальных условия:

(УЧП 1)

,\\
(29.3) (УЧП 2)
\(\frac{\partial u_2}{\partial t} + 2 \frac{\partial u_1}{\partial x} = 0\)
, \(-\infty < x < \infty\) , \(0 < t < \infty\) ,\\
(HУ 1) \(u_1(x, 0) = f(x)\) ,\\
(HУ 2) \(u_2(x, 0) = g(x)\) , \(-\infty < x < \infty\) .

Эта задача может соответствовать определению давления \(u_1(x, t)\) и
плотности \(u_2(x, t)\) как функции пространственной координаты x и
времени t по известным распределениям этих величин в начальный момент.

Запишем систему уравнений в матричной форме

\[
\begin{bmatrix} \frac{\partial u_1}{\partial x} \\ \frac{\partial u_2}{\partial x} \end{bmatrix}
\]

\begin{itemize}
\tightlist
\item
\end{itemize}

\[
, \quad 0 =
\begin{bmatrix} \frac{\partial u_1}{\partial x} \\ \frac{\partial u_2}{\partial x} \end{bmatrix}
\begin{bmatrix} 0 \\ 0 \end{bmatrix}
\]

.

\[
v =
\begin{bmatrix} v_1 \\ v_2 \end{bmatrix}
\begin{bmatrix} 0 & 8 \\ 2 & 0 \end{bmatrix}
u = Pv
\begin{bmatrix} \frac{\partial u_1}{\partial t} \\ \frac{\partial u_2}{\partial t} \end{bmatrix}
u = Pv
\frac{\partial u}{\partial t} = P \frac{\partial v}{\partial t},
\begin{bmatrix} \frac{\partial u_1}{\partial x} \\ \frac{\partial u_2}{\partial x} \end{bmatrix}
\frac{\partial u}{\partial x} = P \frac{\partial v}{\partial x}.
u_t + Au_x = 0
\]

.

\[
Pv_t + APv_x = 0.
v_t + P^{-1}APv_x = 0,
(29.6) v_t + \Lambda v_x = 0.
\]

с помощью преобразования

\[
\Lambda =
\]

.

де P --- матриця, стовпці якої містять власні вектори матриці A (матриця
P, яку вже знають). Виявляється, що після такого перетворення буде
отримана дуже проста система для визначення v (два рівняння відносно
нового невідомого \(v_1\) і \(v_2\) виявляються незалежними). Це
означає, що \(v_1\) і \(v_2\) легко знайти. Після \(v_1\) і \(v_2\)
знайдено, згідно з формулою

находятся искомые функции и, и и2.

Сначала, однако, выясним, как выглядит система для определения v. Чтобы
разобраться в этом, читатель должен обратиться к матричному исчислению.
Продифференцировав обе части соотношения u = Pv, получаем

(29 5)

\[
\begin{bmatrix} 4 & 0 \\ 0 & -4 \end{bmatrix}
\]

.

\begin{verbatim}
\frac{\frac{\partial v_1}{\partial t} + 4 \frac{\partial v_1}{\partial x} = 0, \\ \frac{\partial v_2}{\partial t} - 4 \frac{\partial v_2}{\partial x} = 0.
\end{verbatim}

\[
v_1(x, t) = \varphi(x-4t),
u = Pv =
\begin{bmatrix} 2 & -2 \\ 1 & 1 \end{bmatrix}
\begin{bmatrix} v_1 \\ v_2 \end{bmatrix}
=
=
\]

.

\[
\begin{bmatrix} 2 & -2 \\ 1 & 1 \end{bmatrix}
\begin{bmatrix} \varphi(x - 4t) \\ \psi(x + 4t) \end{bmatrix}
=
=
\begin{bmatrix} v_1 \\ v_2 \end{bmatrix}
\begin{bmatrix} 2\varphi(x - 4t) - 2\psi(x + 4t) \\ \varphi(x - 4t) + \psi(x + 4t) \end{bmatrix}
\]

.

\[
u_1(x, t) = 2\varphi(x-4t) - 2\psi(x+4t), u_2(x, t) = \varphi(x-4t) + \psi(x+4t).
\phi(\xi) = \sin \xi, \ \psi(\xi) = \xi^2
u_1(x, t) = 2 \sin(x-4t) - 2(x+4t)^2,
u_1(x, 0) = f(x),
2\varphi(x) - 2\psi(x) = f(x),
\]

.

\[
\varphi(x) + \psi(x) = g(x),
\varphi(x) = \frac{1}{4} \left[ f(x) + 2g(x) \right],
\psi(x) = \frac{1}{4} \left[ 2g(x) - f(x) \right].
u_{1}(x, t) = 2\varphi(x-4t) - 2\psi(x+4t) =
= \frac{1}{2} \left[ f(x-4t) + 2g(x-4t) \right] - \frac{1}{2} \left[ 2g(x+4t) - f(x+4t) \right],
u_{2}(x, t) = \varphi(x-4t) + \psi(x+4t) =
= \frac{1}{4} \left[ f(x-4t) + 2g(x-4t) \right] + \frac{1}{4} \left[ 2g(x+4t) - f(x+4t) \right].
u_t + A(x, t) u_x + B(x, t) u = 0
Au_t + Bu_x + Cu = 0,
A =
\]

откуда функци і ф и ф легко находятся:

\[
\begin{bmatrix} 1 & 1 \\ 4 & 1 \end{bmatrix}
\]

.

\[
\begin{cases} \frac{\partial u_1}{\partial t} + \frac{\partial u_1}{\partial x} + \frac{\partial u_2}{\partial x} = 0, \\ \frac{\partial u_2}{\partial t} + 4 \frac{\partial u_1}{\partial x} + \frac{\partial u_2}{\partial x} = 0. \end{cases}
\]

(29.9)

\[
u_{tt} = c^2 \left( u_{rr} + \frac{1}{r} u_r + \frac{1}{r^2} u_{\theta\theta} \right)
u(r, \theta, t) = P(r) \Theta(\theta) T(t),
T'' + \lambda^2 c^2 T = 0
\begin{array}{lll} (\text{УЧП}) & u_{tt} = c^2 \left( u_{rr} + \frac{1}{r} u_r + \frac{1}{r^2} u_{\theta\theta} \right), & 0 < r < 1, \ 0 < t < \infty, \\ (\text{ГУ}) & u = 0 & \text{при } r = 1, & 0 < t < \infty, \\ (\text{НУ}) & \begin{cases} u = f(r, \ \theta) \\ u_t = g(r, \ \theta) \end{cases} & \text{при } t = 0. \end{array}
\]

.

\[
u(r, \theta, t) = U(r, \theta) T(t)
\Delta U + \lambda^2 U = 0
\begin{cases} \frac{\partial u_1}{\partial t} + \frac{\partial u_1}{\partial x} + \frac{\partial u_2}{\partial x} = 0, \\ \frac{\partial u_2}{\partial t} + 4 \frac{\partial u_1}{\partial x} + \frac{\partial u_2}{\partial x} = 0. \end{cases}
T'' + \lambda^2 c^2 T = 0
\Delta U = U_{rr} + \frac{1}{r} U_r + \frac{1}{r^2} U_{\theta\theta}.
u(1, \theta, t) = U(1, \theta) T(t) = 0, \quad 0 < t < \infty,
U(1, \theta) = 0.
\Delta U + \lambda^2 U = 0,
U(1, \theta) = 0.
\begin{cases} \Delta U + \mu u = 0 & \text{в круге,} \\ U = 0 & \text{на границе круга} \end{cases}
\begin{array}{ll} (\mathsf{Y}\mathsf{H}\mathsf{\Pi}) & \Delta U + \lambda^2 U = 0, \\ (\mathsf{\Gamma}\mathsf{Y}) & U(1, \theta) = 0, \end{array}
U(r, \theta) = R(r) \Theta(\theta)
r^2R'' + rR' + (\lambda^2r^2 - n^2)R = 0
r^2R'' + rR' + (\lambda^2r^2 - n^2) R = 0
r^2R'' + rR' + (\lambda^2r^2 - n^2)R = 0
R(r) = AJ_n(\lambda r) + BY_n(\lambda r).
R(r) = AJ_n(\lambda r).
\lambda = k_{nm}
U_{nm}(r, \theta) = J_n(k_{nm}r) [A \sin(n\theta) + B \cos(n\theta)].
f_{nm} = k_{nm}c/2\pi.
T_{nm}(t) = A \sin(k_{nm}ct) + B \cos(k_{nm}ct)
T'' + k_{nm}^2 c^2 T = 0.
sU(x) = -V \frac{dU}{dx}, \quad 0 < x < \infty,
U(0) = P/s.
\begin{array}{ll} (\mathsf{Y}\mathsf{H}\mathsf{\Pi}) & \Delta U + \lambda^2 U = 0, \\ (\mathsf{\Gamma}\mathsf{Y}) & U(1, \theta) = 0, \end{array}
H(\xi) =
\begin{cases} 1, & 0 \le \theta < \pi, \\ 0, & \pi \le \theta < 2\pi. \end{cases}
\begin{cases} 0, & \xi < 0, \\ 1, & \xi \geqslant 0. \end{cases}
\cos(3\varphi) = a_0 P_0(\cos\varphi) + a_1 P_1(\cos\varphi) + \dots
\Delta u = 0, \quad 0 < r < 1, \quad 0 \le \varphi \le \pi, \quad 0 \le \theta < 2\pi.
u(1, \varphi) =
\begin{cases} 1, & 0 \le \varphi \le \pi/2, \\ -1, & \pi/2 < \varphi \le \pi. \end{cases}
\Delta u = 0
u_{rr} + \frac{1}{r} u_r + \frac{1}{r^2} u_{\theta\theta} = f(r, \theta), \quad 0 \le r < 1, \quad 0 \le \theta < 2\pi,
\Delta u = -q
u(1, \theta) = 0, \quad 0 \le \theta \le 2\pi.
= -\int_{0}^{2\pi} u_{r}(r) r d\theta = -2\pi r u_{r}(r).
-2\pi r u_r(r) = q.
u(r) = \frac{-q}{2\pi} \ln r = \frac{q}{2\pi} \ln \frac{1}{r}.
u\left(r\right)=-\frac{q}{2\pi}\ln\frac{1}{r}.
(Y \cup \Pi)
u\left(r,\;\theta\right) = \int_{0}^{2\pi} \int_{0}^{1} G\left(r,\;\theta,\;\rho,\;\varphi\right) f\left(\rho,\;\varphi\right) \rho \; d\rho \; d\varphi.
\frac{1}{2\pi} \ln \left( \frac{1}{r} \right)
\frac{1}{2\pi} \ln \frac{1}{R}
u(r, \theta) = \frac{1}{2\pi} \ln 1/R - \frac{1}{2\pi} \ln 1/\overline{R}
-\frac{1}{2\pi}\ln\rho
G(r, \theta, \rho, \phi) = \frac{1}{2\pi} \ln 1/R - \frac{1}{2\pi} \ln 1/\overline{R} + \frac{1}{2\pi} \ln \rho,
R = V \frac{r^2 - 2\rho r \cos(\theta - \phi) + \rho^2}{R},
\overline{R} = V \frac{r^2 - 2\frac{r}{\rho}\cos(\theta - \phi) + 1/\rho^2}{r^2 - 2\frac{r}{\rho}\cos(\theta - \phi) + 1/\rho^2}
u(r, \theta) = \int_{0}^{2\pi} \int_{0}^{1} G(r, \theta, \rho, \varphi) f(\rho, \varphi) \rho d\rho d\varphi
u(r, \theta) = \frac{1}{2\pi} \int_{0}^{2\pi} \int_{0}^{1} \ln (\rho \overline{R}/R) f(\rho, \varphi) \rho d\rho d\varphi.
\begin{array}{llll} (\text{УЧ\Pi}) & \Delta u = 0, & 0 < r < 1, & 0 \leqslant \theta < 2\pi, \\ (\text{ГУ}) & u(1, \theta) = g(\theta), & 0 \leqslant \theta \leqslant 2\pi, \end{array}
u(r, \theta) = \int_{0}^{2\pi} \frac{\partial G}{\partial r} (r, \theta, 1, \varphi) g(\varphi) d\varphi.
u(r, \theta) = \frac{1}{2\pi} \int_{0}^{2\pi} \left[ \frac{1 - r^2}{1 - 2r\cos(\theta - \varphi) + r^2} \right] g(\varphi) d\varphi.
\begin{array}{ll} (\text{УЧП}) & \Delta u = f(r, \theta), & 0 < r < 1, & 0 \le \theta < 2\pi, \\ (\text{ГУ}) & u(1, \theta) = g(\theta), & 0 < \theta \le 2\pi, \end{array}
\Delta u = 1
u_p = (r, \theta) = Ar^2
f(x+h) = f(x) + f'(x)h + \frac{f''(x)}{2!}h^2 + \dots
f(x+h) \cong f(x) + f'(x) h
(37.1) f'(x) \cong \frac{f(x+h)-f(x)}{h}.
f'(x) \cong \frac{f(x) - f(x - h)}{h}
f'(x) \simeq \frac{1}{2h} [f(x+h) - f(x-h)].
f''(x) \cong \frac{1}{h^2} [f(x+h) - 2f(x) + f(x-h)].
u(x+h, y) = u(x, y) + u_x(x, y)h + u_{xx}(x, y)\frac{h^2}{2!} + \dots,
u(x-h, y) = u(x, y) - u_x(x, y)h + u_{xx}(x, y)\frac{h^2}{2!} - \dots,
(Y \cup \Pi)
\begin{split} u_x(x,\ y) & \cong \frac{u(x+h,\ y)-u(x,\ y)}{h}, \\ u_{xx}(x,\ y) & \cong \frac{1}{h^2} \big[ u(x+h,\ y)-2u(x,\ y)+u(x-h,y) \big], \\ u_u(x,\ y) & \cong \frac{u(x,\ y+k)-u(x,\ y)}{k}, \\ u_{uy}(x,\ y) & \cong \frac{1}{h^2} \big[ u(x,\ y+k)-2u(x,\ y)+u(x,\ y-k) \big]. \end{split}
u_{xx} + u_{yy} = 0
\begin{split} u\left(x,\,y\right) &= u_{ij}\,,\\ u\left(x,\,y+k\right) &= u_{i+1,j}\,,\\ u\left(x,\,y-k\right) &= u_{i-1,j}\,,\\ u\left(x-h,\,y\right) &= u_{i,j-1}\,,\\ u\left(x+h,\,y\right) &= u_{i,j+1}\,,\\ u_{x}\left(x,\,y\right) &= \frac{1}{2h}\left(u_{i,j+1}-u_{i,j-1}\right)\,,\\ u_{y}\left(x,\,y\right) &= \frac{1}{2k}\left(u_{i+1,j}-u_{i-1,j}\right)\,,\\ u_{xx}\left(x,\,y\right) &= \frac{1}{h^{2}}\left(u_{i,j+1}-2u_{i,j}+u_{i,j-1}\right)\,,\\ u_{yy}\left(x,\,y\right) &= \frac{1}{h^{2}}\left(u_{i+1,j}-2u_{i,j}+u_{i-1,j}\right)\,. \end{split}
u_{xx} + u_{yy} = 0
\Delta u = \frac{1}{h^2} (u_{i,j+1} - 2u_{i,j} + u_{i,j-1}) + \frac{1}{k^2} (u_{i+1,j} - 2u_{i,j} + u_{i-1,j}) = 0.
u_{i,j} = \frac{1}{4} (u_{i+1,j} + u_{i-1,j} + u_{i,j+1} + u_{i,j-1}).
-4u_{22} + 0 + \sin(\pi/3) + u_{23} + u_{32} = 0,
-4u_{23} + u_{22} + \sin(2\pi/3) + 0 + u_{33} = 0,
-4u_{32} + 0 + u_{22} + u_{33} + 0 = 0,
\]

формулы, выражающие расстояние между двумя точками в полярных
координатах. Чтобы найти решение исходной задачи, осталось провести
наложение всех импульсных откликов системы. Итак, мы переходим к
последнему шагу.

ШАГ 3. Суперпозиция импульсных откликов Этот шаг совсем прост. Из
соотношения

\[
-4u_{33} + u_{32} + u_{23} + 0 + 0 = 0.
\]

получаем

(36.4)

Формула (36.4) виражає розв'язок задачі Діріхле для рівняння Пуассона
всередині кола через функцію Гріна цієї задачі. Якщо щільність заряду
\(f(r, \theta)\) відома, то інтеграл можна обчислити, наприклад,
чисельно.

\section{ЗАМЕЧАНИЯ}\label{ux437ux430ux43cux435ux447ux430ux43dux438ux44f-2}

\begin{enumerate}
\def\labelenumi{\arabic{enumi}.}
\tightlist
\item
  Задачу
\end{enumerate}

також можна розв'язати за допомогою методу \emph{функції Гріна}. У цьому
випадку рішення фіксується у формі

\[
\begin{bmatrix} -4 & 1 & 1 \\ 1 & -4 & 0 & 1 \\ 1 & 0 & -4 & 1 \\ 0 & 1 & 1 & -4 \end{bmatrix}
\]

\begin{verbatim}
Останнє відношення можна надати більш обчислювально корисну форму, явно виставляючи вираз у  \partial G/\partial r  . Внаслідок цього ми отримуємо інтегральну формулу Пуассона, знайдену раніше в лекції 33.
\end{verbatim}

(36.5)

\begin{enumerate}
\def\labelenumi{\arabic{enumi}.}
\setcounter{enumi}{1}
\tightlist
\item
  Решение общей задачи Дирихле
\end{enumerate}

\[
\begin{bmatrix} u_{22} \\ u_{23} \\ u_{32} \\ u_{33} \end{bmatrix}
\]

находится как сумма решений (36.4) и (36.5).

\begin{enumerate}
\def\labelenumi{\arabic{enumi}.}
\setcounter{enumi}{2}
\item
  Метод функцій Гріна дозволяє отримати розв'язки багатьох задач у
  областях різних форм. Однак для кожної області і (точніше, для кожного
  оператора зліва від граничної умови) для кожного рівняння потрібно
  знайти власну функцію Гріна, і це не завжди легко зробити.
\item
  Чтобы реально найти решение по формуле (36.4), в большинстве случаев
  необходимо вычислять интегралы на ЭВМ.
\end{enumerate}

\section{ЗАДАЧИ}\label{ux437ux430ux434ux430ux447ux438-4}

Найти потенциал точечного источника в трехмерном пространстве.

\begin{enumerate}
\def\labelenumi{\arabic{enumi}.}
\setcounter{enumi}{1}
\tightlist
\item
  Найти функцию Грина \(G(x, y, \xi, \eta)\) для уравнения Лапласа в
  верхней полуплоскости y \textgreater{} 0. Другими словами, найти
  потенциал в точке (x, y) в верхней полуплоскости, если заряд
  расположен в точке \((\xi, \eta)\) , а потенциал на границе y = 0
  равен нулю (см. рис.).
\end{enumerate}

\begin{verbatim}
УКАЗАНИЕ. Если отрицательный заряд поместить в точку  \overline{Q} = (\xi, -\eta)  , очевидно, що загальний потенціал на лінії y = 0 буде нульовим. Отже, функція Гріна збігається з потенціалом цих двох зарядів.
\end{verbatim}

\begin{verbatim}
- 3. Використовуючи результати задачі 2, знайдемо розв'язок рівняння Пуассона  \Delta u = -k  у верхній півплощині з умовою нульової межі.
\end{verbatim}

\begin{itemize}
\tightlist
\item
  \begin{enumerate}
  \def\labelenumi{\arabic{enumi}.}
  \setcounter{enumi}{3}
  \tightlist
  \item
    Як би ви побудували функцію Зеленого для першого квадранта
    x\textgreater0,~y\textgreater0 ?
  \end{enumerate}
\item
  \begin{enumerate}
  \def\labelenumi{\arabic{enumi}.}
  \setcounter{enumi}{4}
  \tightlist
  \item
    Інший підхід до розв'язання рівняння Пуассона такий. Припустимо, ви
    хочете розв'язати проблему
  \end{enumerate}
\end{itemize}

(УЧП)

\[
=
,  0 < r < 1  ,  0 \le \theta < 2\pi  ,  (\Gamma Y)   u(1, \theta) = \sin \theta  ,  0 \le \theta \le 2\pi  .
\]

Сначала попытайтесь найти частное решение уравнения с помощью
подстановки

.

Затем подставьте это выражение в уравнение Пуассона и определите
константу \(A_{\cdot}\) Після того, як \(u_p\) знайдено часткове
рішення, необхідно шукати повне рішення у вигляді \(u=w+u_p\) . Яку
граничну умову має задовольняти функція w? Знайдіть \(w(r, \theta)\) .
Яке рішення початкової проблеми \(u(r, \theta)\) ? Подивись. Ретельно
проаналізуйте відповідь. Дайте інтерпретацію кожного члена.

\section{Часть 5. ЧИСЛЕННЫЕ И ПРИБЛИЖЕННЫЕ
МЕТОДЫ}\label{ux447ux430ux441ux442ux44c-5.-ux447ux438ux441ux43bux435ux43dux43dux44bux435-ux438-ux43fux440ux438ux431ux43bux438ux436ux435ux43dux43dux44bux435-ux43cux435ux442ux43eux434ux44b}

Лекция 37

\section{ЧИСЛЕННЫЕ РЕШЕНИЯ (ЭЛЛИПТИЧЕСКИЕ
ЗАДАЧИ)}\label{ux447ux438ux441ux43bux435ux43dux43dux44bux435-ux440ux435ux448ux435ux43dux438ux44f-ux44dux43bux43bux438ux43fux442ux438ux447ux435ux441ux43aux438ux435-ux437ux430ux434ux430ux447ux438}

МЕТА ЛЕКЦІЇ: Показати, як диференціальні рівняння в похідних похідних
можна звести до систем алгебраїчних рівнянь, якщо похідні з частковими
похідними замінити на наближення з скінченними різницями. Розв'язання
цієї системи алгебраїчних рівнянь ітеративним методом дозволяє отримати
наближене розв'язок

Дается информация о пакете прикладных программ ELLPACK, предназначенном
для решения эллиптических

задач общего вида на ЭВМ.

К настоящему моменту мы познакомились с несколькими методами решения
линейных уравнений с частными производными. Однако большинство
уравнений, с которыми мы имели дело, были очень простыми. Так же просты
были и формы областей, в которых решались задачи. Тем не менее многие
задачи нельзя упростить до такой степени, чтобы они свелись к заданному
набору шаблонных задач. Такие задачи следует решать приближенно
численными методами. Прогресс в области высокопроизводительных ЭВМ
привел к созданию новых численных методов. Сегодня уже решены такие
нелинейные задачи динамики жидкости, теории упругости и теории
потенциала, о решении которых десять лет назад никто даже не
задумывался.

Под общим названием «численные методы» объединены несколько различных
подходов к решению задач. Подробное описание всех этих подходов можно
найти в книге {[}4{]} из рекомендованного списка литературы. В этой и
двух последующих лекциях будет показано, как решать эллиптические,
гиперболические и параболические уравнения методом конечных разностей.

Сначала познакомимся с понятием конечных разностей, а затем покажем, как
с их помощью решить задачу Дирихле в квадрате.

\section{Конечно-разностные
аппроксимации}\label{ux43aux43eux43dux435ux447ux43dux43e-ux440ux430ux437ux43dux43eux441ux442ux43dux44bux435-ux430ux43fux43fux440ux43eux43aux441ux438ux43cux430ux446ux438ux438}

Вспомним ряд Тейлора для функции \(f(x)\):

\[
\begin{bmatrix} -\sin{(\pi/3)} \\ -\sin{(2\pi/3)} \\ 0 \\ 0 \end{bmatrix}
\]

Если оборвать этот ряд на втором члене, то получим

.

Откуда

\[
=
\]

Выражение, стоящее в правой части, называется правой разностной
производной. Она аппроксимирует первую производную f'(x) в точке x.

\begin{verbatim}
В разложении Тейлора функции  f(\mathbf{r})  можно заменить h на -h и получить левую разностную производную
\end{verbatim}

(37.2)

.

Buyutan \(f(x - h) \cong f(x) - f'(x) h\)
\nif \(f(x + h) \cong f( x) + f'(x) h\) ,

получаем центральную разностную производную

(37.3)

\[
\begin{bmatrix} -0.86 \\ -0.86 \\ 0 \\ 0 \end{bmatrix}
\]

Якщо в ряду Тейлора залишиться ще один член, то центральну різницю
похідну для наближення f'\,'(x) можна отримати точно таким самим
способом

(37.4)

.

Тепер ми можемо розширити поняття наближення скінченних різниць до
\emph{часткових похідних}. Якщо виходити з розкладу Тейлора функції двох
змінних

\[
u_{xx} + u_{yy} = f(x, y)
xu_{xx} + u_{yy} + 2u = \sin(x - y)
\sin xu_{xx} + u_{xy} + 3u = 0
\]

\pandocbounded{\includegraphics[keepaspectratio]{images/_page_285_Figure_8.jpeg}}

\[
\begin{split} u_x(x,\ y) & \cong \frac{u(x+h,\ y)-u(x,\ y)}{h}, \\ u_{xx}(x,\ y) & \cong \frac{1}{h^2} \big[ u(x+h,\ y)-2u(x,\ y)+u(x-h,y) \big], \\ u_u(x,\ y) & \cong \frac{u(x,\ y+k)-u(x,\ y)}{k}, \\ u_{uy}(x,\ y) & \cong \frac{1}{h^2} \big[ u(x,\ y+k)-2u(x,\ y)+u(x,\ y-k) \big]. \end{split}
\]

\begin{verbatim}
(УЧ\Pi)
\end{verbatim}

\[
u_{xx} + u_{yy} + 2u = 0
u_{xx} + u_{yy} = 0
\begin{cases} u = 0 & \text{квадрата,} \\ \frac{\partial u}{\partial x}(1, y) = 1, & 0 \leqslant y \leqslant 1, \end{cases}
\begin{split} u\left(x,\,y\right) &= u_{ij}\,,\\ u\left(x,\,y+k\right) &= u_{i+1,j}\,,\\ u\left(x,\,y-k\right) &= u_{i-1,j}\,,\\ u\left(x-h,\,y\right) &= u_{i,j-1}\,,\\ u\left(x+h,\,y\right) &= u_{i,j+1}\,,\\ u_{x}\left(x,\,y\right) &= \frac{1}{2h}\left(u_{i,j+1}-u_{i,j-1}\right)\,,\\ u_{y}\left(x,\,y\right) &= \frac{1}{2k}\left(u_{i+1,j}-u_{i-1,j}\right)\,,\\ u_{xx}\left(x,\,y\right) &= \frac{1}{h^{2}}\left(u_{i,j+1}-2u_{i,j}+u_{i,j-1}\right)\,,\\ u_{yy}\left(x,\,y\right) &= \frac{1}{h^{2}}\left(u_{i+1,j}-2u_{i,j}+u_{i-1,j}\right)\,. \end{split}
(\text{УЧП})
u_t = u_{xx}
\]

\pandocbounded{\includegraphics[keepaspectratio]{images/_page_288_Figure_3.jpeg}}

\[
u_t = u_{xx}, \quad 0 < x < 1, \quad 0 < t < \infty,
\begin{cases} u(0, t) = 1, \\ u_x(1, t) = -[u(1, t) - g(t)], \\ u(x, 0) = 0, \quad 0 \le x \le 1. \end{cases}
x_j = jh
u_{t} = \frac{1}{k} [u(x, t+k) - u(x, t)] = \frac{1}{k} [u_{t+1,j} - u_{t,j}],
u_{xx} = \frac{1}{h^{2}} [u(x+h, t) - 2u(x, t) + u(x-h, t)] = \frac{1}{h^{2}} [u_{t,j+1} - 2_{t,j} + u_{t,j-1}].
\]

Це потрібно вирішувати у зв'язку з \(u_{22}\) , \(u_{21}\) , \(u_{32}\)
та \(u_{33}\) . Розв'язок такої системи можна знайти за допомогою
\emph{ітеративного методу}. Метод Лібмана також є одним із ітеративних
методів.

\begin{itemize}
\tightlist
\item
  \begin{enumerate}
  \def\labelenumi{\arabic{enumi}.}
  \setcounter{enumi}{1}
  \tightlist
  \item
    Якщо зменшити кроки сітки на h і k (кількість вузлів сітки
    збільшиться), то отримаємо систему рівнянь типу (37,7), але більших
    розмірів.
  \end{enumerate}
\item
  \begin{enumerate}
  \def\labelenumi{\arabic{enumi}.}
  \setcounter{enumi}{2}
  \tightlist
  \item
    Систему рівнянь (37.7) можна переписати у матричній формі
  \end{enumerate}
\end{itemize}

\[
u_{t+1,j} = u_{t,j} + \frac{k}{h^2} \left[ u_{t,j+1} - 2u_{t,j} + u_{t,j-1} \right].
a_x(1, t) = -[u(1, t) - g(t)].
\begin{bmatrix} u_{22} \\ u_{23} \\ u_{32} \\ u_{33} \end{bmatrix}
\frac{1}{h}[u_{t,n} - u_{t,n-1}] = -[u_{t,n} - g_t].
u_{t,n} = \frac{u_{t,n-1} - hg_t}{1 + h}.
u_{2,j} = u_{1,j} + \frac{k}{h^2} [u_{1,j+1} - 2u_{1,j} + u_{1,j-1}], \quad j = 2, 3, \ldots, n-1.
u_{2,n} = \frac{u_{2,n-1} + hg_2}{1+h}.
u_{tt} = u_{xx}, \quad 0 < x < 1, \quad 0 < t < \infty,
\]

С ростом числа уравнений (например, до 1000) матрица коэффициентов
становится разреженной, т. е. содержит много нулей. Системы с
разреженными матрицами решаются специальными методами. Обычно
используются итерационные методы, такие, как метод Якоби, метод
Гаусса---Зейделя или метод последовательной верхней релаксации.

\begin{enumerate}
\def\labelenumi{\arabic{enumi}.}
\setcounter{enumi}{3}
\item
  При розв'язанні задачі Неймана ті \emph{похідні}, що входять до
  граничної умови, також слід замінити різницею апроксимаціями.
\item
  Подобным образом можно решать задачи
\end{enumerate}

\begin{enumerate}
\def\labelenumi{\alph{enumi})}
\tightlist
\item
\end{enumerate}

\begin{enumerate}
\def\labelenumi{\arabic{enumi})}
\setcounter{enumi}{5}
\tightlist
\item
\end{enumerate}

\[
\begin{cases} u(0, t) = g_1(t), \\ u(1, t) = g_2(t), \end{cases}
\]

\begin{enumerate}
\def\labelenumi{\Alph{enumi})}
\setcounter{enumi}{1}
\tightlist
\item
\end{enumerate}

(неоднородные уравнения с переменными коэффициентами), (уравнения с
переменными коэффициентами).

\begin{enumerate}
\def\labelenumi{\arabic{enumi}.}
\setcounter{enumi}{5}
\tightlist
\item
  Якщо область, в якій розв'язується задача, має неправильну форму, її
  можна покрити сіткою і наближити розв'язок у точках, найближчих до
  межі, інтерполюючи граничні умови. Після цього задачу вирішують
  звичайним способом (див. рисунок 37.2).
\end{enumerate}

\[
\quad 0 < t < \infty,
\]

□ новые ГУ находятся интерполяцией ГУ на кривой

РИС, 37 2.

\begin{itemize}
\tightlist
\item
  \begin{enumerate}
  \def\labelenumi{\arabic{enumi}.}
  \setcounter{enumi}{6}
  \tightlist
  \item
    В некоторых журналах публикуются тексты программ для решения
    уравнений с частными производными на ЭВМ. Ниже приведен краткий
    перечень этих журналов:
  \end{enumerate}
\item
  \begin{enumerate}
  \def\labelenumi{\alph{enumi})}
  \tightlist
  \item
    ACM Transactions on Mathematical Software,
  \end{enumerate}
\item
  \begin{enumerate}
  \def\labelenumi{\arabic{enumi})}
  \setcounter{enumi}{5}
  \tightlist
  \item
    Computer Journal,
  \end{enumerate}
\item
  \begin{enumerate}
  \def\labelenumi{\Alph{enumi})}
  \setcounter{enumi}{1}
  \tightlist
  \item
    Numerische Mathematik,
  \end{enumerate}
\item
  \begin{enumerate}
  \def\labelenumi{\alph{enumi})}
  \setcounter{enumi}{17}
  \tightlist
  \item
    BIT
  \end{enumerate}
\end{itemize}

Кроме того, недавно был разработан пакет прикладных программ ELLPACK,
предназначенный для решения достаточно широкого класса краевых
эллиптических задач. Этот пакет позволит решать разнообразные двух- и
трехмерные задачи в

различных системах координат для произвольных границ и общих граничных
условий. Пользователю представляется большой ассортимент различных
методов решения 1).

\section{ЗАДАЧИ}\label{ux437ux430ux434ux430ux447ux438-5}

\begin{itemize}
\item
  1 Получите аппроксимацию (39.4) для второй производной f'\,`(x)
  f'\,'(x) = \{\(f(x+h)\) 2f(x) + \(f(x-h)\)\}/h\^{}2 .
\item
  \begin{enumerate}
  \def\labelenumi{\arabic{enumi}.}
  \setcounter{enumi}{1}
  \tightlist
  \item
    Проделайте две итерации для задачи Дирихле (37.5), используя
    итерационный метод Либманна.
  \end{enumerate}
\item
  \begin{enumerate}
  \def\labelenumi{\arabic{enumi}.}
  \setcounter{enumi}{2}
  \tightlist
  \item
    К какой алгебранческой системе сведется задача Дирихле для уравнений
    Пуассона в квадрате
  \end{enumerate}
\end{itemize}

\(u_{xx} + u_{yy} = f(x,y), \quad 0 < x < 1, \quad 0 < y < 1,\)
\((\Gamma Y)\) \(u(x,y) = g(x,y)\) на границе, если ее решать ме-

тодом конечных разностей? 4. Решите задачу 3, если

(УЧП)

, 0 \textless{} x \textless{} 1 , 0 \textless{} y \textless{} 1 , (ГУ)
\(u(x, y)\) = \(g(x, y)\) на границе.

\begin{enumerate}
\def\labelenumi{\arabic{enumi}.}
\setcounter{enumi}{4}
\tightlist
\item
  Как бы вы решали задачу Неймана внутри квадрата
\end{enumerate}

(УЧП)

\[
\begin{cases} u(x, 0) = \varphi(x), \\ u_t(x, 0) = \psi(x), \end{cases}
\]

, 0 \textless{} x \textless{} 1 , 0 \textless{} y \textless{} 1 , на
верхней, нижней и левой стороне квадрата,

методом конечных разностей?

\begin{enumerate}
\def\labelenumi{\arabic{enumi}.}
\setcounter{enumi}{5}
\tightlist
\item
  Постройте блок-схему алеоритма решения задачи Дирихле в квадрате
\end{enumerate}

\[
\quad 0 \leqslant x \leqslant 1,
\]

\(u_{xx} + u_{yy} = f(x, y), \quad 0 < x < 1, \quad 0 < y < 1,\)
\((\text{ГУ})\) \(u(x, y) = g(x, y)\) на границе,

\begin{itemize}
\tightlist
\item
  з довільною кількістю вузлів сітки. Якщо ви знаєте мову програмування,
  напишіть програму для виконання цих обчислень.
\item
  \begin{enumerate}
  \def\labelenumi{\arabic{enumi})}
  \tightlist
  \item
    Російською мовою публікуються статті з чисельних методів розв'язання
    диференціальних рівнянь у численних журналах і збірках. Наприклад:
    «Журнал обчислювальної математики та математичної фізики»,
    видавництво «Наука», «Журнал інженерної фізики», видавництво
    Академії наук БРСР, збірки серії «Обчислювальні методи і
    програмування», Видавництво Московського державного університету.
    Багато публікацій, переважно статей, присвячені наборам прикладних
    проблем. Як приклад, можна вказати: Горбунов-Посадов М. М., Карпов
    В. Я, Корягін та ін. Пакет SAFRA. Програмне забезпечення для
    обчислювальних експериментів. У: Пакети прикладних програм.
    Обчислювальний експеримент. Москва, Наука, 1983. Примітка. Переклад.
  \end{enumerate}
\end{itemize}

\section{ЯВНЫЕ РАЗНОСТНЫЕ
СХЕМЫ}\label{ux44fux432ux43dux44bux435-ux440ux430ux437ux43dux43eux441ux442ux43dux44bux435-ux441ux445ux435ux43cux44b}

МЕТА ЛЕКЦІЇ: Ввести концепцію явної різницької схеми та показати, як її
можна використовувати для розв'язання гіперболічних і параболічних
задач. Основна ідея полягає в тому, що після заміни рівняння типу

\[
u_{tt} \simeq \frac{1}{h^2} [u(x, t+k) - 2u(x, t) + u(x, t-k)],
\]

его конечно-разностной аппроксимацией получаются формулы, явно
выражающие значения решения для одного момента времени через значения
решения в предыдущий момент времени. Таким образом, смешанную задачу для
параболического или гиперболического уравнения можно решать,
последовательно вычисляя решения для всех последующих моментов времени.

Явна схема не позбавлена недоліків. Якщо необхідно підвищити точність
наближення похідних, то при зміні сегмента збільшується не лише обсяг
обчислень, а й похибки округлення.

Попередня лекція була присвячена розв'язанню еліптичних задач на межі
лінії (стаціонарних задач). У таких задачах необхідно знайти розв'язок
диференціального рівняння в похідних у заданій області простору, якщо
розв'язок або його похідна задано на межі області визначення Приблизне
розв'язання еліптичних задач зводиться до розв'язку системи алгебраїчних
рівнянь для значень функції у внутрішніх вузлах сітки. Іншими словами,
значення розв'язку у всіх внутрішніх вузлах визначаються одночасно.

У цій лекції ми розглянемо диференціальні методи розв'язання задач,
залежних від часу. Основна ідея полягає в тому, що якщо розв'язок
відомий у початковий момент часу, то можливо знайти розв'язок за
допомогою схеми підрахунку тика, коли
\(t=\Delta t,\ 2\Delta t,\ 3\Delta t,\ldots\) Замінюючи часткові похідні
часу та просторову змінну на скінченні різниці, можна отримати явні
вирази для \(u_{tj}\) за допомогою значень функції u у попередні часи.
Такий процес називається явною схемою підрахунку з тикуванням.

Чтобы показать этот метод в действии, рассмотрим типичную задачу
теплопроводности.

\section{Явная схема для уравнения
теплопроводности}\label{ux44fux432ux43dux430ux44f-ux441ux445ux435ux43cux430-ux434ux43bux44f-ux443ux440ux430ux432ux43dux435ux43dux438ux44f-ux442ux435ux43fux43bux43eux43fux440ux43eux432ux43eux434ux43dux43eux441ux442ux438}

Рассмотрим задачу теплопроводности в стержне, начальная температура
которого равна нулю. Пусть температура левого конца фиксирована, а на
правом конце происходит теплообмен с окружающей средой, так что тепловой
полок пропорционален Р И

\[
u_{xx} \simeq \frac{1}{h^2} [u(x+h, t) - 2u(x, t) + u(x-h, t)].
\]

С. 38.1. Сетка для уравнения теплопроводности.

разности температур конца стержия и среды. Пусть температура среды
определяется функцией \(g\left(t\right)\) . Другими словами, мы решаем
задачу

задачу (УЧП)

\[
u_t(x, 0) \approx \frac{1}{k} [u(x, k) - u(x, 0)] = \frac{1}{k} [u(x, k) - \varphi(x)].
\]

(38.1) (ГУ)

\[
u(x, t+k) = 2u(x, t) - u(x, t-k) + + (k/h)^2 [u(x+h, t) - 2u(x, t) + u(x-h, t)].
\]

Для решения этой задачи методом конечных разностей построим
прямоугольную сетку, узлы которой определяются формулами (см. рис. 38.1)

\[
\frac{1}{k}\left[u\left(x,\,k\right)-\varphi\left(x\right)\right]=\psi\left(x\right),
\]

, \(j = 0, 1, 2, ..., n\) , \(y_i = ik\) , \(i = 0, 1, 2, ..., m\) .

Отметим, что значения \(u_{ij}\) на левой и нижней сторонах сетки на
рис. 38.1 известны из граничных и начальных условий и наша вадача
состоит в отыскании остальных значений \(u_{ij}\) . Для решения этой
задачи заменим частные производные в уравнении теплопроводности их
конечно-разностными аппроксимациями

\[
u_t = u_{xx}
\begin{cases} u(0, t) = 0, \\ u(1, t) = 0, \\ u(x, 0) = \sin(\pi x), \quad 0 \le x \le 1. \end{cases}
\]

\pandocbounded{\includegraphics[keepaspectratio]{images/_page_293_Figure_6.jpeg}}

\begin{verbatim}
\begin{array}{ccccccccccccccccccccccccccccccccccc
\end{verbatim}

\[
\frac{\Delta t}{(\Delta x)^2} \leqslant 0.5.
\]

\begin{verbatim}
\begin{array}{ccccccccccccccccccccccccccccccccccc
\end{verbatim}

\[
\begin{split} u_t(x,\ t) &= \frac{1}{k} \big[ u(x,\ t+k) - u(x,\ t) \big], \\ u_{xx}(x,\ t) &= \frac{\lambda}{h^2} \big[ u(x+h,\ t+k) - 2u(x,\ t+k) + u(x-h,\ t+k) \big] + \\ &\quad + \frac{(1-\lambda)}{h^2} \big[ u(x+h,\ t) - 2u(x,\ t) + u(x-h,\ t) \big], \end{split}
\frac{1}{k} (u_{i+1, j} - u_{i, j}) = \frac{\lambda}{h^2} (u_{i+1, j+1} - 2u_{i+1, j} + u_{i+i, j-1}) + \frac{(1-\lambda)}{h^2} (u_{i, j+1} - 2u_{i, j} + u_{i, j-1}),
\]

(39.3)

\[
(\Gamma Y) \begin{cases} u_{i, 1} = 0, & i = 1, 2, ..., m, \\ u_{i, n} = 0, & i = 1, 2, ..., m, \\ u_{i, j} = 1, & j = 2, ..., n-1. \end{cases}
-\lambda r u_{i+1, j+1} + (1+2r\lambda) u_{i+1, j} - \lambda r u_{i+1, j-1} = = r(1-\lambda) u_{i, j+1} + [1-2r(1-\lambda)] u_{i, j} + r(1-\lambda) u_{i, j-1},
\]

\pandocbounded{\includegraphics[keepaspectratio]{images/_page_296_Figure_5.jpeg}}

\pandocbounded{\includegraphics[keepaspectratio]{images/_page_296_Figure_7.jpeg}}

\[
\begin{array}{l} -u_{21}+3u_{22}-u_{23}=u_{11}-u_{12}+u_{13}=1,\\ -u_{22}+3u_{23}-u_{24}=u_{12}-u_{13}+u_{14}=1,\\ -u_{23}+3u_{24}-u_{25}=u_{13}-u_{14}+u_{15}=1,\\ -u_{24}+3u_{25}-u_{26}=u_{14}-u_{15}+u_{16}=1. \end{array}
\begin{cases} u(0, t) = g_1(t), \\ u(1, t) = g_2(t), \end{cases}
\]

(HY)

\[
\begin{bmatrix} 3 & -1 & 0 & 0 \\ -1 & 3 & -1 & 0 \\ 0 & -1 & 3 & -1 \\ 0 & 0 & -1 & 3 \end{bmatrix}
\quad 0 \leqslant x \leqslant 1,
\begin{bmatrix} u_{2\bar{z}} \\ u_{\bar{z}a} \\ u_{24} \\ u_{\bar{z}x} \end{bmatrix}
\]

а начальное условие--- по схеме

\[
=
\]

В результате для вычисления величины \(u(x, t+k)\) получаем следующую
явную схему:

(38.5)

З (38.5) видно, що для обчислення розв'язку на заданому часовому шарі
необхідно знати розв'язок на двох попередніх шарах. Отже, щоб почати
рахувати, потрібно використовувати P И С

\[
\begin{bmatrix} 1 \\ 1 \\ 1 \\ 1 \end{bmatrix}
\]

. 38.4. Использование массивов при реализации явной схемы.

начальным условнем на скорость

.

з якої ми виводимо \(u(x, k) = \varphi(x) + k\psi(x)\) , тобто цінність
розв'язку на \(t = \Delta t\) . Розв'язок у наступні часи можна знайти
за явною формулою (38.5).

\section{ЗАДАЧИ}\label{ux437ux430ux434ux430ux447ux438-6}

\begin{enumerate}
\def\labelenumi{\arabic{enumi}.}
\tightlist
\item
  Постройте явную конечно-разностную схему для задачи
\end{enumerate}

(УЧП)

\[
,  0 < x < 1  ,  0 < t < \infty  ,
\]

(ГУ)

\[
\begin{bmatrix} b_1 & c_1 & 0 & 0 & \dots & & & 0 \\ a_1 & b_2 & c_2 & 0 & \dots & & & 0 \\ 0 & a_{\bar{z}} & b_3 & c_3 & \dots & & & 0 \\ & & & & & & & & \\ \vdots & & & & & & & & \\ 0 & 0 & 0 & \dots & & & & a_{n-1} & b_n \end{bmatrix}
\]

\begin{verbatim}
(HУ)  u(x, 0) = \sin(\pi x), \quad 0 \le x \le 1.  Полагая  h = \Delta x = 0.1  , найдите решение при  t_1 = 0.005  ,  t_2 = 0.010  ,  t_3 = 0.015  . Постройте график полученного решения на сетке x = 0; 0.1; 0.2; 0.3; . .; 0.9: 1 при t = 0.015.
\end{verbatim}

РИС. 38.5. Зависимость полной ошибки от шага сетки.

\begin{itemize}
\tightlist
\item
  \begin{enumerate}
  \def\labelenumi{\arabic{enumi}.}
  \setcounter{enumi}{1}
  \tightlist
  \item
    Решите задачу 1 аналитически (методом разделения переменных).
    Найдите значение решения в узлах сетки x=0; 0,1; 0,2; 0,3; \ldots;
    0,9; 1 при t=0,015. Сравните эти результаты с результатами
    численного решения задачи 1. (Для табулирования аналитического
    решения можно воспользоваться калькулятором или написать маленькую
    программу для ЭВМ.)
  \end{enumerate}
\item
  \begin{enumerate}
  \def\labelenumi{\arabic{enumi}.}
  \setcounter{enumi}{2}
  \tightlist
  \item
    В соответствии с замечанием 3 постройте блок-схему решения
    гиперболической задачи.
  \end{enumerate}
\item
  \begin{enumerate}
  \def\labelenumi{\arabic{enumi}.}
  \setcounter{enumi}{3}
  \tightlist
  \item
    Решите задачу 1 с новым граничным условием при x=1
    \(u_x(1, t) = -[u(1, t) 1].\) \# НЕЯВНЫЕ РАЗНОСТНЫЕ СХЕМЫ (СХЕМА
    КРАНКА---НИКОЛЬСОНА)
  \end{enumerate}
\end{itemize}

ЦЕЛЬ ЛЕКЦИИ: Показать, как можно решать задачи, содержащие функции,
зависящие от времени, с помощью неявной разностной схемы. В этой схеме,
как и рачьше, частные производные заменяются конечно-разностными
аппроксимациями, но \(u_{l+1,\ J}\) уже не выражаются в явном виде через
значения на предыдущих слоях. Теперь для определения величин
\(u_{l+1,\ J}\) необходимо решить систему уравнений. Другими словами,
для определения решения на каждом временном наге необходимо решать
систему уравнений.

Перевага неявних схем над явними полягає в тому, що в неявних схемах
танг сітки можна зробити достатньо великим без страху, що помилки
округлення «знищать» розв'язок.

Для решения параболической задачи мы воспользуемся широко известной
\emph{схемой Кранка} --- \emph{Никольсона}.

Как уже говорилось в предыдущей лекции, для надежной работы явной схемы
необходимо, чтобы временной шаг сетки был мал. В частности, если простая
задача для уравнения теплопроводности

(39.1)

\[
\begin{bmatrix} x_1 \\ x_{\bar{z}} \\ x_{\bar{s}} \\ \vdots \\ x_n \end{bmatrix}
\]

решается по явной схеме, то шаги сетки \(\Delta t\) и \(\Delta x\)
должны удовлетворять условию

\begin{verbatim}
Інакше метод буде чисельно нестабільним (а помилки округлення зростуть нескінченно). (Питання чисельної стабільності розглядаються у [53].) Іншими словами, якщо відстань сітки на осі x вибрана  \Delta x=0,1  , то часовий крок  \Delta t  не може бути більшим за  0.5\Delta x^2=0.005  . Це означає, що ви можете пройти від t=0 до t=1, зробивши 200 кроків.
\end{verbatim}

Существуют, однако, такие схемы (неявные схемы), которые позволяют вести
вычисления с гораздо большим шагом, при этом,

правда, растет объем вычислений, совершаемых на одном шаге. По таким
схемам можно вести вычисления с достаточно большим шагом, но на каждом
шаге приходится решать систему алгебраимеских уравнений. Для иллюстрации
этого метода решим следующую задачу теплопроводности.

\section{Неявная схема для уравнения
теплопроводности}\label{ux43dux435ux44fux432ux43dux430ux44f-ux441ux445ux435ux43cux430-ux434ux43bux44f-ux443ux440ux430ux432ux43dux435ux43dux438ux44f-ux442ux435ux43fux43bux43eux43fux440ux43eux432ux43eux434ux43dux43eux441ux442ux438}

Рассмотрим задачу

(39 2)

\[
=
\]

Воспользуемся следующими конечно-разностными аппроксимациями для частных
производных \(u_t\) и \(u_{xx}\) :

где \(\lambda\) выбирается из отрезка {[}0,1{]}. Отметим, что \(u_{xx}\)
аппроксимируется взвешенным средним центральных разностных производных в
момент t и t+k. При \(\lambda=0.5\) получается обычное среднее этих двух
центральных производных, а при \(\lambda=0.75\) ---одна из разностных
производных берется с весом 0.75, а вторая---с весом 0.25. При
\(\lambda=0\) получается обычная явная схема, о которой шла речь в
предыдущей лекции.

После замены частных производных \(u_t\) и \(u_{xx}\) в задаче (39.2) мы
получаем разностную задачу (см. рис. 39.1).

Разностное уравнение

\[
\begin{bmatrix} d_t \\ d_{\bar{z}} \\ d_{\bar{s}} \\ \vdots \\ \vdots \\ d_n \end{bmatrix}
\]

(39.3)

,

298

Перенесем все неизвестные значения u с верхнего временного слоя (с
индексом i+1) в левую часть уравнения (39.3) и получим (39.4)

де вводиться позначення r = k/h\^{}2 . Зверніть увагу, що якщо i
фіксоване, а j змінюється від 2 до n-1, співвідношення (39.4) визначають
систему n-2 рівнянь з n-2 невідомим \(u_{t+1,2}\), \(u_{t+1,3}\),

\[
\begin{bmatrix} 1 & c_1^* & 0 & 0 & \dots & 0 \\ 0 & 1 & c_2^* & 0 & \dots & 0 \\ 0 & 0 & 1 & c_2^* & \dots & 0 \\ \vdots & \vdots & & & \vdots & \vdots \\ 0 & 0 & \dots & & 1 \end{bmatrix}
\]

РИС. 39.1. Сетка для неявной разностной схемы.

РИС. 39.2. Шаблон для неявной схемы.
\(u_{i+1, 4}, \ldots, u_{i+1, n-1}\) , которые являются решением задачи
во внут-

ренних узлах сетки на временном слое \(t=(i+1)\Delta t\) .

Рис. 39.2 дает наглядное представление о ссруктуре каждого уравнения
системы (39.4). Перейдем теперу к решению системы (39.4).

\section{Алгоритм решения задачи
(39.2)}\label{ux430ux43bux433ux43eux440ux438ux442ux43c-ux440ux435ux448ux435ux43dux438ux44f-ux437ux430ux434ux430ux447ux438-39.2}

ШАГ 1. Выбираем некоторое значение \(\lambda\) ( \(0 \le \lambda \le 1\)
). Если \(\lambda = 0\) , то уравнения (39.4), переходят в явные формулы
из лекции 38.

ШАГ 2. Пусть, например, \(h=\Delta x=0,2\) и \(k=\Delta t=0,08\) (при
этом \(r=k/h^2=2\) ). В данном случае сетка содержит 6 узлов вдоль осн x
(4 внутренних узла) (см. рис. 39.1). Возьчем весовой параметр
\(\lambda=0,5\) (получающаяся при этом схема называется схемой
Кранка---Никольсона). В соответствии с вычислительной схемой (ее принято
называть шаблоном), изображенной на рис. 39.2, двигаясь слева направо
(j=2,3,4,5) по первым двум слоям (i=1), получаем следующие четыре
уравнения:

\[
\begin{bmatrix} x_1 \\ x_2 \\ x_3 \\ \vdots \\ x_n \end{bmatrix}
\]

Перепишем их в матричной форме

(39.5)

\[
=
\begin{bmatrix} u_{2\bar{z}} \\ u_{\bar{z}a} \\ u_{24} \\ u_{\bar{z}x} \end{bmatrix}
\begin{bmatrix} d_1^* \\ d_2^* \\ d_3^* \\ \vdots \\ d_n^* \end{bmatrix}
\]

.

Матрица этой системы называется трехдиагональной. Для того чтобы решить
трехдиагональную систему

\[
c_1^* = c_1/b_1
x_n = d_n^*, \quad x_j = d_j^* - c_j^* x_{j+1}, \quad j = n-1, \ n-2, \ \dots, \ 2, \ 1.
\begin{bmatrix} x_1 \\ x_{\bar{z}} \\ x_{\bar{s}} \\ \vdots \\ x_n \end{bmatrix}
\]

\begin{verbatim}
\begin{array}{ll} (\mathrm{Y}\mathrm{H}\Pi) & u_t = u_{xx}, \quad 0 < x < 1, \quad 0 < t < \infty, \\ (\Gamma\mathrm{Y}) & \left\{ \begin{array}{ll} u_{\,0}(0, \ t) = 1, \\ u_{\,x}(1, \ t) + u_{\,}(1, \ t) = g_{\,}(t), \\ u_{\,x}(1, \ t) = 0, \quad 0 \leqslant \kappa \leqslant 1. \end{array} \end{array}
\end{verbatim}

\[
\begin{bmatrix} d_t \\ d_{\bar{z}} \\ d_{\bar{s}} \\ \vdots \\ \vdots \\ d_n \end{bmatrix}
\]

\begin{verbatim}
\right. & 0 < t < \infty, \end{array}
\end{verbatim}

преобразуем ее в эквивалентную систему вида

\[
u_t = u_{xx} + u, \quad 0 < x < 1,
\begin{bmatrix} x_1 \\ x_2 \\ x_3 \\ \vdots \\ x_n \end{bmatrix}
u(x, t) = e^{-\pi^2 t} \sin(\pi x).
\begin{bmatrix} d_1^* \\ d_2^* \\ d_3^* \\ \vdots \\ d_n^* \end{bmatrix}
\]

где

Розв'язок якого показано на рис. 40.1. Що краще --- мати аналітичне
рішення цієї проблеми?

(40.2)

\[
,  c_{j+1}^* = \frac{c_{j+1}}{b_{j+1} - a_j c_j^*}  ,  j = 1, 2, \ldots, n-2  , и  d_1^* = d_1/b_1  ,  d_{j+1}^* = \frac{d_{j+1} - a_j d_j^*}{b_{j+1} - a_j c_j^*}  ,  j = 1, 2, \ldots, n-1  .
\]

В этом преобразовании нет ничего таниственного, поскольку вторая система
полностью эквивалентна первой. Матрица новой системы устроена так, что
эта система очень легко решается. Решая уравнения последовательно снизу
вверх, получаем

или численное, представленное при \(\alpha = 1\) в табл. 40.1?

Вопрос хороший, а ответ на него зависит от того, что мы хотим делать с
этим решением дальше. Однако каждому из типов решений присущи очевидные
преимущества.

\section{Преимущества аналитического
решения}\label{ux43fux440ux435ux438ux43cux443ux449ux435ux441ux442ux432ux430-ux430ux43dux430ux43bux438ux442ux438ux447ux435ux441ux43aux43eux433ux43e-ux440ux435ux448ux435ux43dux438ux44f}

\begin{enumerate}
\def\labelenumi{\arabic{enumi}.}
\tightlist
\item
  Очевидно, що розв'язок (40.2) є більш інформативним, ніж таблиця
  чисел. Якщо ми хочемо обчислити розв'язок у певній точці (x, t), ми
  можемо зробити це так точно, як хочемо, просто збільшуючи кількість
  членів ряду, які розглядаються. Водночас легко отримати оцінку зверху
  щодо значення допущеної помилки.
\end{enumerate}

\pandocbounded{\includegraphics[keepaspectratio]{images/_page_301_Figure_4.jpeg}}

РИС. 40.1. Решение уравнения теплопроводности (40.1) в различные моменты
времени.

\begin{itemize}
\tightlist
\item
  \begin{enumerate}
  \def\labelenumi{\arabic{enumi}.}
  \setcounter{enumi}{1}
  \tightlist
  \item
    Аналітичне розв'язання завжди дозволяє обчислити значення розв'язку
    в одній точці (x, t) без необхідності обчислювати значення
    розв'язків в інших точках, як це відбувається при розв'язанні задач
    із явними або неявними різницею схем.
  \end{enumerate}
\item
  \begin{enumerate}
  \def\labelenumi{\arabic{enumi}.}
  \setcounter{enumi}{2}
  \tightlist
  \item
    Аналітичне рішення дозволяє визначити розв'язок у будь-якій точці, а
    не лише на вузлах сітки.
  \end{enumerate}
\item
  \begin{enumerate}
  \def\labelenumi{\arabic{enumi}.}
  \setcounter{enumi}{3}
  \tightlist
  \item
    Найважливішою перевагою аналітичного розв'язку для нас є можливість
    простежити вплив фізичних параметрів, початкових і граничних умов на
    природу розв'язку.
  \end{enumerate}
\end{itemize}

Чисельні методи не виявляють цих регулярностей, оскільки дозволяють
знаходити розв'язок лише за заданих параметрів --- початкових і
граничних умовах. Іноді важливо знати взаємозв'язок між параметрами
моделі та розв'язком, особливо коли йдеться про оцінку фізичних
параметрів за типом розв'язку. Наприклад, припустимо, що розв'язок і
визначається експериментально, і ми знаємо аналітичний розв'язок

u = функция от параметров.

Тогда можно поставить вопрос об определении параметров нак функции
входных данных по схеме

Параметры = функция от u = функция исходных данных.

Такая задача называется параметрической идентификацией. Для
параметрической идентификации необходимо уметь решать уравнение с
частными производными. Несколько позже в этой лекции мы приведем пример
параметрической идентификации в биологии, но сначала посмотрим, чем же
хороши численные решения.

\section{Преимущества численных
решений}\label{ux43fux440ux435ux438ux43cux443ux449ux435ux441ux442ux432ux430-ux447ux438ux441ux43bux435ux43dux43dux44bux445-ux440ux435ux448ux435ux43dux438ux439}

Головна перевага чисельних розв'язків полягає в тому, що їх можна
отримати навіть якщо аналітичне рішення неможливо. Майже всі нелінійні
диференціальні рівняння в похідних похідних потрібно розв'язувати
чисельними методами, і більшість реальних моделей фізики, хімії,
біології тощо є нелінійними. Як правило, лінійні моделі наближаються до
нелінійних лише якщо в них відкинути нелінійні члени. Наприклад, ось
кілька нелінійних рівнянь:

\begin{enumerate}
\def\labelenumi{\arabic{enumi})}
\item
  нелинейное волновое уравнение \(u_{tt} = u_{xx} + f(u)\) ;
\item
  уравнение реакции с нелинейной диффузией \(u_t = u_{xx} + f(u)\) ;
\item
  система Ходжкина---Хаксли
\end{enumerate}

\begin{verbatim}
Застосовуючи цей метод до системи чотирьох рівнянь (39,5), отримуємо розв'язок  u_{22}=0,60  ;  u_{23}=0,80  ;  u_{24}=0,80  ;  u_{23}=0,60  . Ми отримали наближені значення розв'язку у внутрішніх точках сітки на  t=\Delta t  . Тепер ми можемо зробити наступний крок у часі, але для цього нам доведеться розв'язати нову систему рівнянь.
\end{verbatim}

В неявной схеме объем вычислений на каждом шаге больше, чем в явной, но
хорошую точность можно получить даже при гораздо большем шаге.

\section{ЗАДАЧИ}\label{ux437ux430ux434ux430ux447ux438-7}

\begin{itemize}
\tightlist
\item
  \begin{enumerate}
  \def\labelenumi{\arabic{enumi}.}
  \tightlist
  \item
    Из дифференциального уравнения (39.3) получите уравнение (39.4).
  \end{enumerate}
\item
  \begin{enumerate}
  \def\labelenumi{\arabic{enumi}.}
  \setcounter{enumi}{1}
  \tightlist
  \item
    Постройте неявную конечно-разностную схему для задачи
  \end{enumerate}
\end{itemize}

\[
u(x, 0) = Ae^{-x^2/a}
\begin{aligned} (\text{УЧ}\Pi) & u_t = Du_{xx} - Vu_x, & -\infty < x < \infty, \\ (\text{HY}) & u(x, 0) = Ae^{-x^2, a}, & -\infty < x < \infty. \end{aligned}
u(x, t) = \frac{A}{2\sqrt{\pi Dt}} \int_{-\infty}^{+\infty} e^{-\xi^2/a} e^{-(x-\xi)^2/4Dt} d\xi =
= \frac{A\sqrt{a}}{\sqrt{a+4Dt}} e^{-[x^2/(a+4Dt)]},
\]

\texttt{-\ 4.\ Как\ выглядит\ шаблон\ для\ уравнения\ (39.4),\ когда}
\lambda = 1 ? - 5. Постройте блок-схему розв'язки рівняння
теплопровідності (39.2). Якщо есть возможность, то напишите программу
для ЭВМ Було бы очень полезно решить эту задачу численно при различных
значениях λ с начальным условием и (x, 0) = sin (πx). Сравните
полученные при различных λ численные результаты с аналитическим
решением, яке в данном случае имеет вид

,

\begin{verbatim}
\overline{C} = 0   (по той же причине, что и   \overline{A} = 0  ),
\end{verbatim}

6 Воспользовавшись формулами, приведенными в тексте, решите систему
алгебраических уравнений (39.5). Лекція 40 \# СРАВНЕНИЕ АНАЛИТИЧЕСКИХ
РЕШЕНИЙ С ЧИСЛЕННЫМИ

ЦЕЛЬ ЛЕКЦИИ: Сравнить достоинства и недостатки аналитических и численных
методов розв'язки уравнений з частинними похідними. Приведена важная
математическая задача об идентификации физических величин
(параметрическая идентификация). В связи с цієї задачей рассматривается
один важный пример из биологии.

Вероятно, пришло время обсудить достоинства и недостатки аналитических и
численных решений уравнений з частинними похідними. Сначала выясним, що
мы подразумеваем под этими двумя типами решений.

\section{Аналитические
розв'язки}\label{ux430ux43dux430ux43bux438ux442ux438ux447ux435ux441ux43aux438ux435-ux440ux43eux437ux432ux44fux437ux43aux438}

Под аналитическими решениями мы подразумеваем такие розв'язки, в которых
неизвестная функція и выражена через независимые змінні и параметры
системи в виде формул, бесконечных рядов и интегралов.

\section{Численные
розв'язки}\label{ux447ux438ux441ux43bux435ux43dux43dux44bux435-ux440ux43eux437ux432ux44fux437ux43aux438}

Під чисельним ми зрозуміємо чисельні розв'язки, отримані чисельно після
приблизної заміни початкового рівняння іншим, простішим рівнянням.
Наприклад, у методі скінченної різниці похідні виражаються через
скінченні різниці, а розв'язок диференціального рівняння в похідних у
похідних апроксимується розв'язком різницького рівняння. Результатом
такої процедури зазвичай є таблиця значень розв'язків, а для деяких
значень пояснювальних змінних.

Таблина 40.1

{\def\LTcaptype{none} % do not increment counter
\begin{longtable}[]{@{}
  >{\raggedright\arraybackslash}p{(\linewidth - 22\tabcolsep) * \real{0.1788}}
  >{\raggedright\arraybackslash}p{(\linewidth - 22\tabcolsep) * \real{0.0861}}
  >{\raggedright\arraybackslash}p{(\linewidth - 22\tabcolsep) * \real{0.1192}}
  >{\raggedright\arraybackslash}p{(\linewidth - 22\tabcolsep) * \real{0.0728}}
  >{\raggedright\arraybackslash}p{(\linewidth - 22\tabcolsep) * \real{0.0662}}
  >{\raggedright\arraybackslash}p{(\linewidth - 22\tabcolsep) * \real{0.0662}}
  >{\raggedright\arraybackslash}p{(\linewidth - 22\tabcolsep) * \real{0.0331}}
  >{\raggedright\arraybackslash}p{(\linewidth - 22\tabcolsep) * \real{0.0331}}
  >{\raggedright\arraybackslash}p{(\linewidth - 22\tabcolsep) * \real{0.0331}}
  >{\raggedright\arraybackslash}p{(\linewidth - 22\tabcolsep) * \real{0.0728}}
  >{\raggedright\arraybackslash}p{(\linewidth - 22\tabcolsep) * \real{0.1192}}
  >{\raggedright\arraybackslash}p{(\linewidth - 22\tabcolsep) * \real{0.1192}}@{}}
\toprule\noalign{}
\begin{minipage}[b]{\linewidth}\raggedright
\end{minipage} & \begin{minipage}[b]{\linewidth}\raggedright
×
\end{minipage} & \begin{minipage}[b]{\linewidth}\raggedright
\end{minipage} & \begin{minipage}[b]{\linewidth}\raggedright
\end{minipage} & \begin{minipage}[b]{\linewidth}\raggedright
\end{minipage} & \begin{minipage}[b]{\linewidth}\raggedright
\end{minipage} & \begin{minipage}[b]{\linewidth}\raggedright
\end{minipage} & \begin{minipage}[b]{\linewidth}\raggedright
\end{minipage} & \begin{minipage}[b]{\linewidth}\raggedright
\end{minipage} & \begin{minipage}[b]{\linewidth}\raggedright
\end{minipage} & \begin{minipage}[b]{\linewidth}\raggedright
\end{minipage} & \begin{minipage}[b]{\linewidth}\raggedright
\end{minipage} \\
\midrule\noalign{}
\endhead
\bottomrule\noalign{}
\endlastfoot
& 0 & 0,1 & 0,2 & 0,3 & 0,4 & 0,5 & 0,6 & 0,7 & 0,8 & 0,9 & 1 \\
00,010,020,03 & 100 & 10,20,15 & 10,34 & l: & : & : & ! & : & 10,34 &
10,20,15 & 1000 \\
; & & : & : & : & : & : & : & : & : & : & : \\
\end{longtable}
}

Теперь, когда мы знаем, що подразумевается под решением каждого из цих
типов, давайте сравним их между собой.

\section{\texorpdfstring{\textbf{Сравнительный анализ численного} и
аналитического
решений}{Сравнительный анализ численного и аналитического решений}}\label{ux441ux440ux430ux432ux43dux438ux442ux435ux43bux44cux43dux44bux439-ux430ux43dux430ux43bux438ux437-ux447ux438ux441ux43bux435ux43dux43dux43eux433ux43e-ux438-ux430ux43dux430ux43bux438ux442ux438ux447ux435ux441ux43aux43eux433ux43e-ux440ux435ux448ux435ux43dux438ux439}

Рассмотрим простейшую смещанную параболическую задачу

(40.1)

\[
\bar{D} = A\xi_{xx} + B\xi_{xy} + C\xi_{yy} + D\xi_x + E\xi_y = -2(x^2 + y^2)
\]

Розв'язок якого показано на рис. 40.1. Що краще --- мати аналітичне
рішення цієї проблеми?

(40.2)

\[
\widetilde{E} = A\eta_{xx} + B\eta_{xy} + C\eta_{yy} + D\eta_x + E\eta_y = 2(y^2 - x^2),
\]

или численное, представленное при \(\alpha = 1\) в табл. 40.1?

Вопрос хороший, а ответ на него зависит от того, що мы хотим делать с
этим решением дальше. Однако каждому из типов решений присущи очевидные
преимущества.

\section{Преимущества аналитического
розв'язки}\label{ux43fux440ux435ux438ux43cux443ux449ux435ux441ux442ux432ux430-ux430ux43dux430ux43bux438ux442ux438ux447ux435ux441ux43aux43eux433ux43e-ux440ux43eux437ux432ux44fux437ux43aux438}

\begin{enumerate}
\def\labelenumi{\arabic{enumi}.}
\tightlist
\item
  Очевидно, що розв'язок (40.2) є більш інформативним, ніж таблиця
  чисел. Якщо ми хочемо обчислити розв'язок у певній точці (x, t), ми
  можемо зробити це так точно, як хочемо, просто збільшуючи кількість
  членів ряду, які розглядаються. Водночас легко отримати оцінку зверху
  щодо значення допущеної помилки.
\end{enumerate}

\pandocbounded{\includegraphics[keepaspectratio]{images/_page_301_Figure_4.jpeg}}

Рис. 40.1. Розв'язок рівняння теплопровідності (40.1) в различные
моменты времени.

\begin{itemize}
\tightlist
\item
  \begin{enumerate}
  \def\labelenumi{\arabic{enumi}.}
  \setcounter{enumi}{1}
  \tightlist
  \item
    Аналітичне розв'язання завжди дозволяє обчислити значення розв'язку
    в одній точці (x, t) без необхідності обчислювати значення
    розв'язків в інших точках, як це відбувається при розв'язанні задач
    із явними або неявними різницею схем.
  \end{enumerate}
\item
  \begin{enumerate}
  \def\labelenumi{\arabic{enumi}.}
  \setcounter{enumi}{2}
  \tightlist
  \item
    Аналітичне рішення дозволяє визначити розв'язок у будь-якій точці, а
    не лише на вузлах сітки.
  \end{enumerate}
\item
  \begin{enumerate}
  \def\labelenumi{\arabic{enumi}.}
  \setcounter{enumi}{3}
  \tightlist
  \item
    Найважливішою перевагою аналітичного розв'язку для нас є можливість
    простежити вплив фізичних параметрів, початкових і граничних умов на
    природу розв'язку.
  \end{enumerate}
\end{itemize}

Чисельні методи не виявляють цих регулярностей, оскільки дозволяють
знаходити розв'язок лише за заданих параметрів --- початкових і
граничних умовах. Іноді важливо знати взаємозв'язок між параметрами
моделі та розв'язком, особливо коли йдеться про оцінку фізичних
параметрів за типом розв'язку. Наприклад, припустимо, що розв'язок і
визначається експериментально, і ми знаємо аналітичний розв'язок

u = функція от параметров.

Тогда можна поставить вопрос об определении параметров нак функції
входных данных по схеме

Параметры = функція от u = функція исходных данных.

Такая задача називається параметрической идентификацией. Для
параметрической идентификации необходимо уметь решать уравнение з
частинними похідними. Несколько позже в цієї лекции мы приведем пример
параметрической идентификации в биологии, но сначала посмотрим, чем же
хороши численные розв'язки.

\section{Преимущества численных
решений}\label{ux43fux440ux435ux438ux43cux443ux449ux435ux441ux442ux432ux430-ux447ux438ux441ux43bux435ux43dux43dux44bux445-ux440ux435ux448ux435ux43dux438ux439-1}

Головна перевага чисельних розв'язків полягає в тому, що їх можна
отримати навіть якщо аналітичне рішення неможливо. Майже всі нелінійні
диференціальні рівняння в похідних похідних потрібно розв'язувати
чисельними методами, і більшість реальних моделей фізики, хімії,
біології тощо є нелінійними. Як правило, лінійні моделі наближаються до
нелінійних лише якщо в них відкинути нелінійні члени. Наприклад, ось
кілька нелінійних рівнянь:

\begin{enumerate}
\def\labelenumi{\arabic{enumi})}
\item
  нелинейное хвильове рівняння \(u_{tt} = u_{xx} + f(u)\) ;
\item
  уравнение реакции с нелинейной диффузией \(u_t = u_{xx} + f(u)\) ;
\item
  система Ходжкина---Хаксли
\end{enumerate}

\[
u(x, t) = \frac{A\sqrt{a}}{\sqrt{a+4Dt}}e^{-[(x-Vt)^{2}/(a+4Dt)]}.
\]

\(Аналитические розв'язки для цих уравнений не известны ни при каких нелинейных функциях f и\)
\(g^{(1)}\) . Поэтому общий подход к решению нелинейных (а в ряде
случаев и линейных) задач обычно базируется на численных решениях.

Тепер давайте розглянемо приклад того, як можна використати аналітичне
рішення для визначення фізичних параметрів. Детальніший опис цих завдань
можна знайти у рекомендованій літературі.

\section{Параметрическая идентификация (в
биологии)}\label{ux43fux430ux440ux430ux43cux435ux442ux440ux438ux447ux435ux441ux43aux430ux44f-ux438ux434ux435ux43dux442ux438ux444ux438ux43aux430ux446ux438ux44f-ux432-ux431ux438ux43eux43bux43eux433ux438ux438}

Припустимо, що біолог намагається визначити швидкість, з якою \((K^+)\)
іони калію дифундують у екзоплазмовому розчині. Якщо ви знаєте цей
коефіцієнт дифузії, можна багато сказати про те, як нервові імпульси
передаються вздовж аксонів. Проблема в тому, що це співвідношення майже
неможливо знайти прямим вимірюванням. Однак можливо знайти математичний

1) Це не так. В качестве примера можна привести уравнение
\(u_{ii} = u_{xx} - \sin u_x - \Pi pum\) , ped.

связь между концентрацией поташа \(u\left(x,\,t\right)\) и коэффициентом
диффузии \(D,\,\,a\,\,\) затем по измерениям величины
\(u\left(x,\,t\right)\) определить величину D.

Біологи Ходжкін і Кейз виявили, що після розміщення гігантського аксону
у спеціальному соляному розчині концентрація радіохктивного калію (**2K)
вздовж аксону приблизно описується кривою в початковий момент

,

.

Отже, параметри A і A слід обирати так, щоб крива відповідала
експериментальним даним (наприклад, у сенсі методу найменших квадратів).
Також було встановлено, що P И С

\pandocbounded{\includegraphics[keepaspectratio]{images/_page_303_Figure_6.jpeg}}

. 40.2. Начальная концентрация ионов ( \(^{42}K\) ): a---теоретическая
кривая; \(\delta\) ---экспериментальные точки.

с течением времени ноны поташа под действием диффузии и конвекции
растекаются вдоль аксона. Ходжкин и Кейс предположили, що концентрация и
является решением следующей задачи с начальными данными для уравнения
диффузии:

(40.3)

\[
\overline{G} = G = 0
\]

Цю проблему вирішили у лекції 15, переключившись на рухому систему
координат. Якщо читач пам'ятає, загальна схема методу виглядає так.
Спочатку відкидаємо конвективний член і отримуємо розв'язок суто
дифузійної задачі (при V = 0):

,

подставим их в уравнение

\[
= \frac{A\sqrt{a}}{\sqrt{a+4Dt}} e^{-[x^2/(a+4Dt)]},
\overline{A}u_{\xi\xi} + \overline{B}u_{\xi\eta} + \overline{C}u_{\eta\eta} + \overline{D}u + \overline{E}u_{\eta} + \overline{F}u = \overline{G}
u(x, t) = \frac{A\sqrt{a}}{\sqrt{a+4Dt}}e^{-[(x-Vt)^{2}/(a+4Dt)]}.
\]

и получим

\[
SS = \sum_{i=1}^{n} [y_i - (a + bx_i)]^2.
u_{\xi\eta} = \frac{-(x^2 + y^2) u_{\xi} + (y^2 - x^2) u_{\eta}}{8x^2y^3}.
u_t = \frac{1}{r} \frac{\partial}{\partial r} \left( Dr u_r - s \omega^2 r^2 u \right), \quad 0 < r < 1,
\]

\begin{verbatim}
ШАГ 2. Выражая, наконец, x и y через   \xi   и   \eta  , окончательно получаем
\end{verbatim}

(23.8)

\[
Au_{xx} + Bu_{xy} + Cu_{yy} + Du_x + Eu_y + Fu = G
u_{\xi\eta} = \frac{\eta u_{\xi} - \xi u_{\eta}}{2(\xi^2 - \eta^2)}.
\begin{cases} 0, & t < x/V, \\ P, & t \ge x/V. \end{cases}
\]

Задача решена. Теперь мы знаем, что нужно сделать, чтобы найти новые
координаты и привести уравнение к канонической форме.

\section{ЗАМЕЧАНИЯ}\label{ux437ux430ux43cux435ux447ux430ux43dux438ux44f-3}

\begin{enumerate}
\def\labelenumi{\arabic{enumi}.}
\tightlist
\item
  На самом деле для гиперболических уравнений существует две
  канонические формы. Вторую можно получить из первов заменой переменных
  вида
\end{enumerate}

Конечно, все це очень просто! Не сложнее, чем сбрасывать що-нибудь на
ленту транспортера, а затем наблюдать за его движением. Однако все
становится очень интересным, якщо примесь диффундирует в среде. Чтобы
разобраться в том, що про-исходит с движущейся волной при наличии
диффузии, решим

следующую задачу

(15.2)

\[
\alpha = \alpha (\xi, \eta) = \xi + \eta,
\]

З формальною заміною змінних \(\xi = \xi(x, y)\) та
\(\eta = \eta(x, y)\) (41.1) зводиться до тієї ж форми

(41.2)

\begin{verbatim}
\beta = \beta (\xi, \eta) = \xi - \eta.  Воспользуемся этой заменой для преобразования уравнения (23.8). Так как
\end{verbatim}

где, как обычно, \(H(x)\) --- функція Хевисайда Начальное распределение
концентрации изображено на рис 15 4.

\pandocbounded{\includegraphics[keepaspectratio]{images/_page_114_Figure_4.jpeg}}

Рис. 15.3 Конвективная волна a---передиції фронт примеси, \(\delta\)
---на одну единицу вперед волнового фронта

Обратим внимание на то, що в задаче (15.2) мы отодвинули границу в ---
\(\infty\) (теперь решается задача Коши), так що она не оказывает теперь
влияние на оценку розв'язок. Чтобы решить задачу (15.2), можна
воспользоваться преобразованием Лапласа Р И С

\pandocbounded{\includegraphics[keepaspectratio]{images/_page_114_Figure_7.jpeg}}

. 15 4 початкові умови в задаче конвективной диффузии.

ло переменной t или преобразованием Фурье по переменной x. Однако в
цьому случае гораздо интереснее поступить совсем по-другому Введем новую
систему координат, яка движется вдоль старой со скоростью V. Другими
словами, вместо системи координат, привязанной к берегу реки, мы
рассмотрим систему координат, яка движется со скоростью фронта примеси
(конечно, при налични диффузии фронт буде «смазан»). С точки зрения
математики це означает, що мы заменяем пространственную координату x на
новую \(\xi = x - Vt\) .

Теперь ясно, що

когда \(\xi = 0\) , мы находимся па фронте распространяющейся примеси.

когда \(\xi = 1\) , мы находимся на одну единицу длины впереди фронта.

когда \(\xi = -1\) , мы находимся на одну единицу длины позади фронта.

Наша задача --- преобразовать исходную задачу с НУ

\[
\overline{D} = A\xi_{xx} + B\xi_{xy} + C\xi_{yy} + D\xi_{x} + E\xi_{y},
[\xi_x/\xi_y] = -B/2A
\]

TO

(23 9)

\[
\xi(x, y) = c
u_{\alpha\alpha} - u_{\beta\beta} = \frac{-\beta u_{\alpha} - \alpha u_{\beta}}{2\alpha\beta}.
\frac{dy}{dx} = -\left[\xi_x/\xi_y\right] = B_t 2A.
\]

\begin{verbatim}
При желании можно выразить   \alpha   и   \beta   через исходные координаты x и y. В результате получаем
\end{verbatim}

\[
\overline{B} = 2A\xi_x\eta_x + B(\xi_x\eta_y + \xi_y\eta_x) + 2C\xi_y\eta_y,
\alpha = \xi + \eta = (y^2 - x^2) + (y^2 + x^2) = 2y^3,
\]

1) Завдяки рівності

І з \(B^2 - 4AC = 0\) році , то

\[
\beta = \xi - \eta = (y^2 - x^2) - (y^2 + x^2) = -2x^2.
\]

Оскільки
\(\xi_x/\xi_y = -B \ 2A = -2 \sqrt{A\tilde{C}}/2A = -\sqrt{C/A}\)
останній вираз для B буде записаний як

\[
u_{\xi\xi} - u_{\eta\eta} = \Psi(\xi, \, \eta, \, u_{\xi}, \, u, u_{\eta}).
\]

І оскільки \(\overline{A}\) дорівнює нулю, то \(\overline{B}\) , звісно,
також дорівнює нулю.

Оскільки при нашому виборі \(\xi\) обидва коефіцієнти \(\overline{A}\) і
\(\overline{B}\) повертаються до нуля, ми можемо взяти будь-яку функцію
як \(\eta\) (якщо вона не пропорційна координаті \(\xi\) ). У прикладі
можна вибрати \(\eta = y\) .

Залишається знайти канонічну форму рівняння в нових координатах. Для
цього потрібно просто підставити формули \(\xi\) і \(\eta\) (41.3) і
знайти всі шанси \(\overline{A}\) , \(\overline{B}\) , \(\overline{C}\)
, \(\overline{D}\) , \(\overline{E}\) , \(\overline{F}\) і
\(\overline{G}\) . На цьому вивчення параболічного випадку завершується,
але перш ніж перейти до іншої теми, розглянемо простий приклад.

\section{\texorpdfstring{Перетворення параболічного рівняння
\(u_{xx} + 2u_{xy} + u_{yy} = 0\) у канонічну
форму}{Перетворення параболічного рівняння u\_\{xx\} + 2u\_\{xy\} + u\_\{yy\} = 0 у канонічну форму}}\label{ux43fux435ux440ux435ux442ux432ux43eux440ux435ux43dux43dux44f-ux43fux430ux440ux430ux431ux43eux43bux456ux447ux43dux43eux433ux43e-ux440ux456ux432ux43dux44fux43dux43dux44f-u_xx-2u_xy-u_yy-0-ux443-ux43aux430ux43dux43eux43dux456ux447ux43dux443-ux444ux43eux440ux43cux443}

Розглянемо рівняння

\[
3u_{xx} + 7u_{xy} + 2u_{yy} = 0
\]

Його коефіцієнти мають такі значення: A=1, B=2, C=1, D=E=F=G=0. Отже,
\(B^2-4AC=0\) для всіх значень x і y. Щоб знайти нові координати \(\xi\)
і \(\eta\) та канонічну форму рівняння, ми рухаємося наступним чином.

КРОК 1. Запишемо \emph{характеристичне рівняння} (тепер лише одне, а не
два)

\[
u_{\xi\eta} = \Phi(\xi, \eta, u, u_1, u_2).
\]

Розв'язуючи це рівняння відносно y, отримуємо

\[
u_{\alpha\alpha} - u_{\beta\beta} = \Psi(\alpha, \beta, u, u_{\alpha}, u_{\beta}).
\]

,

Де --- характеристика \(\xi = y - x\) визначена. Використання цієї
характеристики забезпечує досягнення рівності \(\overline{A} = 0\) .
Координату \(\eta\) можна вибрати довільно (якщо вона не залежить від
\(\xi\) ). Ми оберемо його наступним чином:

\[
u_{xx} + 4u_{xy} = 0.
\]

.

Нові координати

\[
u_{tt} = e^{2} (u_{xx} + u_{yy} + u_{zz}),
\]

изображены на рис. 41 1.

\pandocbounded{\includegraphics[keepaspectratio]{images/_page_309_Figure_3.jpeg}}

Рис. 41.1. Новая система координат \(\xi = y - x\) , \(\eta = y\) .

ШАГ 2. Осталось найти каноническую форму уравнения. Подставляя \(\xi\) и
\(\eta\) в формулы для коефіцієнтів \(\overline{A}\) , \(\overline{B}\)
, \(\overline{C}\) , \(\overline{D}\) , \(\overline{E}\) ,
\(\overline{F}\) и \(\overline{G}\) , получаем \(\overline{A} = 0\) (так
и должно быть; мы находили \(\xi\) из условия равенства нулю цього
коэффициента), \(\overline{B}=0\) (мы уже раньше показали, що этог
коефіцієнт також обращается в нуль),

(41.5)

\[
\overline{D} = A \xi_{xx} + B \xi_{xy} + C \eta_{yy} + D \xi_x + E \xi_y = 0,
\]

Отже, нове рівняння

(НУ)

\[
\xi = x - Vt
u(x, t) =
\]

є загальним розв'язком рівняння

в трехмерном пространстве и показать, что полученное решение
удовлетворяет \emph{принципу Гюйгенса}. Для решения соответствующей
двумерной задачи

(УЧП)

Перед нами розв'язок задачи конвективной диффузии (15.2). Его очень
легко интерпретировать, якщо представить, що мы движемся относительно
графика, изображенного на рис. 15.6.

\pandocbounded{\includegraphics[keepaspectratio]{images/_page_117_Figure_3.jpeg}}

Рис. 15.6. Диффузня из области высокой концентрации в область низкой
концентрации. Чем больше коефіцієнт дифузії, тем быстрее установится
стационарное значение.

Другими словами, в зависимости от относительной величины D (коэффициента
диффузии) и V (скорости потока) розв'язок движется слева направо со
скоростью V и в то же самое время передний фронт расплывается со
скоростью, определяемой величиной D.

\pandocbounded{\includegraphics[keepaspectratio]{images/_page_117_Figure_6.jpeg}}

Рис. 15.7. Розв'язок задачи конвективнои диффузии. Вещество одновременно
движется и диффундирует.

(На рис. 15.7 показано, как расплывается передний фронт концентрации.)

\section{ЗАМЕЧАНИЯ}\label{ux437ux430ux43cux435ux447ux430ux43dux438ux44f-4}

Преобразование координат является важным методом розв'язки уравнений з
частинними похідними. Выбрав подходящую систему координат, можна
существенно упростить уравнение.

\section{ЗАДАЧИ}\label{ux437ux430ux434ux430ux447ux438-8}

\begin{enumerate}
\def\labelenumi{\arabic{enumi}.}
\tightlist
\item
  Решите задачу Коши:
\end{enumerate}

\[
u_{tt} = c^2 (u_{xx} + u_{yy}), \quad
\]

Але тепер у \(B^2-4AC<0\) . Переходячи до нових пояснювальних змінних,
ми хочемо перетворити її у нову форму

\[
\begin{cases} -\infty < x < \infty, \\ -\infty < y < \infty, \end{cases}
\]

\begin{enumerate}
\def\labelenumi{\arabic{enumi}.}
\setcounter{enumi}{1}
\tightlist
\item
  Найдите розв'язок задачи Коши для уравнения конвективной диффузии:
\end{enumerate}

\[
\begin{cases} -\infty < x < \infty, \\ -\infty < y < \infty, \\ -\infty < z < \infty, \end{cases}
\]

Це можливо, але лише шляхом введення складних координат. Щоб знайти ці
комплексні координати \(\xi\) і \(\eta\) , розв'яжемо характеристичні
рівняння, так само як у гіперболічному випадку:

\[
\begin{cases} u(x, y, z, 0) = \varphi(x, y, z), \\ u_{t}(x, y, z, 0) = \psi(x, y, z). \end{cases}
\]

Воспользуйтесь преобразованием из лекции 8.

\begin{enumerate}
\def\labelenumi{\arabic{enumi}.}
\setcounter{enumi}{2}
\tightlist
\item
  Найдите розв'язок задачи переноса:
\end{enumerate}

\[
P_{2}
\frac{dy}{dx} = \sqrt{-4x^2} = 2ix
\begin{aligned} u_{tt} &= c^2 \Delta u, \\ (HY) & \begin{cases} u(x, y, z, 0) = 0, \\ u_t(x, y, z, 0) = \psi(x, y, z), \end{cases} \end{aligned}
\\ -\infty < x < \infty, \quad 0 < t < \infty,
\]

\begin{verbatim}
Цю задачу можна розв'язати методом  \Phi  Юр'єра  ^{1}  ), і виявляється, що розв'язок можна представити у вигляді
\end{verbatim}

\[
u_t = Du_{xx} - Vu_x, \quad -\infty < x < \infty, \quad 0 < t < \infty,
P_{1}
u(x, 0) = e^{-x^2}, \quad -\infty < x < \infty.
\]

\begin{verbatim}
где  \overline{\psi}  — *среднее* значение начального распределения  \psi  *по сфере* радиуса ct с центром в точке (x, y, z), т. е.
\end{verbatim}

\[
y^2 u_{xx} + x^2 u_{yy} = 0,
\]

\pandocbounded{\includegraphics[keepaspectratio]{images/_page_177_Picture_3.jpeg}}

\[
\frac{dy}{dx} = \frac{B - \sqrt{B^2 - 4AC}}{2A} = -\frac{\sqrt{-4x^2y^2}}{2y^2} = -i\frac{x}{y},
\]

\pandocbounded{\includegraphics[keepaspectratio]{images/_page_177_Picture_7.jpeg}}

\[
u(\xi, \tau) = \frac{1}{2\sqrt{D\pi\tau}} \int_{-\infty}^{+\infty} e^{-\beta^2} e^{-(\xi-\beta)^2/4D\tau} d\beta.
\begin{aligned} (\text{YAII}) & u_{tt} = c^2 \Delta u, \\ (\text{HY}) & \begin{cases} u = 0, \\ u_t = \varphi, \end{cases} \end{aligned}
u_{tt} = \alpha^2 u_{xx}
u_h(r, \theta) = \sum_{n=0}^{\infty} r^n [a_n \cos(n\theta) + b_n \sin(n\theta)].
u_{\xi\eta} = \psi(\xi, \eta, u, u_{\xi}, u_{\eta}).
\Delta u = -\frac{r^2}{2} - \frac{r^3}{2} \cos{(2\theta)}
\xi(x, y) = y^2, \quad \eta(x, y) = x^2.
u = \frac{\partial}{\partial t} [t\overline{\varphi}].
T \sin \theta_2 - T \sin \theta_1 =
u_p(r, \theta) = Ar^4 + Br^4 \cos(2\theta).
\Delta x \rho u_{tt} = T \left[ u_x \left( x + \Delta x, \ t \right) - u_x \left( x, \ t \right) \right] + \Delta x F \left( x, \ t \right) - \frac{1}{\Delta x \rho u_t} \left( x, \ t \right) - \Delta x \gamma \dot{u} \left( x, \ t \right)
u(x, t) = \frac{1}{2c} \int_{x-ct}^{x+ct} \varphi(s) ds.
u_{tt} = \alpha^{3} u_{xx} - \beta u_{t} - \gamma u + F(x, t).
u_p(r, \theta) = -\frac{r^4}{32} - \frac{r^4}{24} \cos(2\theta).
(16.2) u_{tt} = \alpha^2 u_{xx},
u_t(x, t) = \frac{1}{2} [\varphi(x + ct) + \varphi(x - ct)].
u_{tt} = ku_{xx}
\sum_{n=0}^{\infty} [a_n \cos(n\theta) + b_n \sin(n\theta)] - \frac{1}{32} - \frac{1}{24} \cos(2\theta) = 0.
u_{tt} = \frac{\partial}{\partial x} [\alpha^2(x) u_x],
u_1(r, \theta) = \frac{1}{32} + \frac{1}{24}r^2\cos(2\theta) - \frac{r^4}{24}\cos(2\theta) - \frac{r^4}{32} = \frac{(r^4 - 1)}{32} - \frac{(r^4 - r^2)}{24}\cos(2\theta).
u(x, 0) = f(x)
u = u_0 + u_1 = r \cos \theta - \frac{(r^4 - 1)}{32} - \frac{(r^4 - r^2)}{24} \cos (2\theta).
i_x + Cv_t + Gv = 0, \\ v_x + Li_t + Ri = 0,
i_{xx} + Gv_x - CLi_{tt} - CRi_t = 0.
v_{x} = -Li_{t} - Ri,
P_{2}
i_{t,s} = CLi_{tt} + (CR + GL)i_t + GRi.
\Delta u = 0, \quad 0 < r < 1 + \frac{1}{4} \sin \theta,
v_{xx} = CLv_{tt} + (CR + GL)v_t + GRv.
\]

записывается в виде

\[
v_{tt} = \alpha^2 v_{xx}, \\ i_{tt} = \alpha^2 i_{xx}, \qquad \alpha^2 = 1/CL,
u \left( 1 + \frac{1}{4} \sin \theta, \ \theta \right) = \cos \theta, \quad 0 \le \theta \le 2\pi.
u_{tt} = u_{xx} - u_t
\]

\begin{verbatim}
где  \overline{\phi}  и  \overline{\psi}  — средние значения функций  \phi  и  \psi  на сфере раднусом ct с центром в точке (x, y, z).
\end{verbatim}

\begin{verbatim}
Цей розв'язок зазвичай називають формулою Пуассона для хвильового рівняння у тривимірному просторі. Це природне узагальнення формули д'Аламбера на тривимірний випадок. Найважливішим у формулі Пуассона є те, що обидва інтеграли в  \bar{\phi}  і  \bar{\psi}  взяті з поверхні сфери. Цей факт дозволяє нам надати таку важливу інтерпретацію рішення.
\end{verbatim}

\begin{verbatim}
У момент  t=t_1  розв'язок u у (x,y,z) залежить лише від початкових розподілів сфери з радіусом  ct_1  центрованим на (x,y,z) (рисунок 24.3). Тепер припустимо, що початкові розподіли  \phi  і  \psi  дорівнюють нулю всюди, крім малої сфери (див. рисунок 24.3). З часом радіус P И С
\end{verbatim}

\pandocbounded{\includegraphics[keepaspectratio]{images/_page_179_Picture_3.jpeg}}

. 24.3. Схема поясияет, как начальное возмущение влияет на точку (к, ч,
г).

\begin{verbatim}
Сфера, що оточує точку (x, y, z), зростає зі швидкістю c, так що через  t_{\hat{z}}  секунди ця сфера починає перетинати область початкового збурення, і тому розв'язок u(x, y, z, t) стає ненульовим. У  t_{\hat{z}} < t < t_{\hat{a}}  розв'язок U у (x, y, z) залишається ненульовим, оскільки сфера продовжує перетинати область початкового збурення, але на  t=t_{\hat{a}}  розв'язок знову стає нульовим. Іншими словами, розв'язок, що поширюється з початкової області збурення, має чітко визначену задню кромку. Цей загальний принцип відомий як принцип Гюйгенса для тривимірного простору. Виявляється, що передній край хвилі завжди чітко окреслений, але задній край чітко окреслений лише в просторах розмірності 3, 5, 7.... Ми вже знаємо з формули д'Аламбера, що початкові умови
\end{verbatim}

\pandocbounded{\includegraphics[keepaspectratio]{images/_page_349_Figure_5.jpeg}}

Що ви можете сказати про точність?

\begin{enumerate}
\def\labelenumi{\arabic{enumi}.}
\setcounter{enumi}{3}
\tightlist
\item
  Як отримати послідовність випадкових точок, що лежать всередині
  трикутника T?
\end{enumerate}

\pandocbounded{
\includegraphics[keepaspectratio]{images/_page_319_Figure_13.jpeg}}

\pandocbounded{\includegraphics[keepaspectratio]{images/_page_319_Figure_14.jpeg}}

\begin{enumerate}
\def\labelenumi{\arabic{enumi}.}
\setcounter{enumi}{4}
\tightlist
\item
  Як отримати послідовність випадкових чисел, розподіл яких показаний на
  рисунку?
\end{enumerate}

\pandocbounded{\includegraphics[keepaspectratio]{images/_page_319_Figure_16.jpeg}}

Іншими словами, як отримати послідовність чисел \{0, 1, 2\}, якщо
ймовірність виникнення 0 і 2 дорівнює 0,25, а ймовірність появи 1 ---
0,5?

Лекція 43

\section{РОЗВ'ЯЗАННЯ ДИФЕРЕНЦІАЛЬНИХ РІВНЯНЬ У ПОХІДНИХ ПОХІДНИХ
МОНТЕ-КАРЛО}\label{ux440ux43eux437ux432ux44fux437ux430ux43dux43dux44f-ux434ux438ux444ux435ux440ux435ux43dux446ux456ux430ux43bux44cux43dux438ux445-ux440ux456ux432ux43dux44fux43dux44c-ux443-ux43fux43eux445ux456ux434ux43dux438ux445-ux43fux43eux445ux456ux434ux43dux438ux445-ux43cux43eux43dux442ux435-ux43aux430ux440ux43bux43e}

МЕТА ЛЕКЦІЇ: Показати, як побудувати таку гру, результатом якої є
приблизне розв'язання диференціального рівняння. Одна з цих ігор (модель
випадкової ходьби) веде до скінченної різниці апроксимації задачі
Діріхле в квадраті. Цю модель можна узагальнити для отримання розв'язків
інших задач.

У попередній лекції було сказано, що можливо винайти такі ігри на удачу,
результатом яких є розв'язки (можливо, приблизні) задач для рівнянь з
частковими коефіцієнтами проникнення. У цій лекції ми покажемо, як
побудувати таку гру для наближеного розв'язання задачі Діріхле (PPP)
\(u_{xx} + u_{yy} = 0\) , 0 \textless{} x \textless{} 1, 0 \textless{} y
\textless{} 1,

(GI)

\[
u\left(1+\frac{1}{4}\sin\theta,\ \theta\right)=\cos\theta
R(A) = g_1 P_A(p_1) + g_2 P_A(p_2) + \dots + g_{12} P_A(p_{12}).
\]

\begin{verbatim}
приводят, вообще говоря, к нерезкому заднему фронту (поскольку формула Даламбера подразумевает интегрирование функции  \psi(x)  в пределах от x-ct до x+ct).
\end{verbatim}

Покажем теперь, что принцип Гюйгенса неприменим к цилиндрическим волнам.
Такого рода решения описывают волны на воде от точечного источника,
когда задний фронт волны не резкий, а постепенно затухает до нуля.

\section{Двумерное волновое
уравнение}\label{ux434ux432ux443ux43cux435ux440ux43dux43eux435-ux432ux43eux43bux43dux43eux432ux43eux435-ux443ux440ux430ux432ux43dux435ux43dux438ux435}

Чтобы решить двумерную волновую задачу

(УЧП)

\[
R(A) = 1 \cdot (0.04) + 1 \cdot (0.15) + 1 \cdot (0.03) + 0 \cdot (0.06) + \dot + 0 \cdot (0.04).
u(1+\epsilon\sin\theta, \theta) = \cos\theta, \quad 0 \le \epsilon \le 1/4.
R(A) = \frac{1}{4} [R(B) + R(C) + R(D) + R(E)]
f(x+h) = f(x) + f'(x)h + f''(x)\frac{h^2}{2!} + \dots,
u_{i,j} = \frac{1}{4} \left[ u_{i-1,j} + u_{i+1,j} + u_{i,j-1} + u_{i,j+1} \right],
\]

(PPP)

\[
u_{xx}=u_{l,\ j+1}-2u_{l,\ j}+u_{l,\ j-1},
\]

\pandocbounded{\includegraphics[keepaspectratio]{images/_page_350_Figure_2.jpeg}}

\[
u_{i, j} = \frac{u_{i, j+1} + u_{i, j-1} + \sin x_j (u_{i+1, j} + u_{i-1, j})}{2 (1 + \sin x_j)}.
u(1+\varepsilon\sin\theta,\ \theta)=u(1,\ \theta)+u_r(1,\ \theta)(\varepsilon\sin\theta)+u_{rr}(1,\ \theta)\frac{(\varepsilon\sin\theta)^2}{21}+\ldots
\]

(i, j+1)

\[
\Delta u = 0, \quad 0 < r < 1 + \varepsilon \sin \theta,
u(A) = g_1 P_A(p_1) + g_2 P_A(p_2) + \dots + g_N P_A(p_N),
u(1, \theta) + u_r(1, \theta) (\varepsilon \sin \theta) + u_{rr}(1, \theta) \frac{(\varepsilon \sin \theta)^2}{2!} + \dots = \cos \theta.
u_{xx} + u_{yy} = 0
(46.8) u = u_0 + \varepsilon u_1 + \varepsilon^2 u_2 + \dots.
u_{xx} + x^2 u_{yy} = 0
P_{0}
(\text{UBP}) \qquad
\begin{cases} \Delta u_{0} = 0, & 0 < r < 1 \text{ (внутри круга),} \\ u_{0}(1, \theta) = \cos \theta, & u_{0}(r, \theta) = r \cos \theta, \end{cases}
\]

\begin{enumerate}
\def\labelenumi{\arabic{enumi}.}
\setcounter{enumi}{4}
\tightlist
\item
  Разработайте схему метода Монте-Карло для розв'язки смешанной
  параболической задачи
\end{enumerate}

\[
\begin{cases} u(x, \ 0) = \varphi(x), \\ u_t(x, \ 0) = 0 \end{cases}
\]

\begin{verbatim}
\right. \end{array}
\end{verbatim}

\[
P_{1}
J[y] = \int_{b}^{a} F(x, y, y') dx,
\begin{cases} \Delta u_{1} = 0, & 0 < r < 1 \text{ (внутри круга),} \\ u_{1}(1, \theta) = -\sin \theta \frac{\partial u_{0}(1, \theta)}{\partial r} = -\sin \theta \cos \theta \end{cases}
J[y] = \int_{0}^{1} [y^{2}(x) + y'^{2}(x)] dx.
\frac{df(x)}{dx} = 0
T = \int_{0}^{T} dt = \int_{a}^{L} \frac{dt}{ds} ds = \int_{0}^{L} \frac{ds}{v} = \frac{1}{\sqrt{2mg}} \int_{0}^{L} \frac{ds}{\sqrt{y}} = \frac{1}{\sqrt{2mg}} \int_{a}^{b} \sqrt{\frac{1+y'^{2}}{y}} dy
u = u_0 + \frac{1}{4}u_1 + \left(\frac{1}{4}\right)^2 u_1 + \cdots
J[y] = \int_{a}^{b} F(x, y, y') dx.
u_u + \frac{1}{4} u_t
\frac{df(x)}{dx} = 0
\]

\begin{enumerate}
\def\labelenumi{\arabic{enumi}.}
\setcounter{enumi}{1}
\tightlist
\item
  Результаты задачи 1 используйте для решения следующей задачи:
\end{enumerate}

(УЧП)

\[
J[y] = \int_{a}^{y} F(x, y, y') dx
u_t = (1+x) u_{xx}, \\ u(x, 0) = q(x), \quad -\infty < x < \infty,
J[y] = \int_{a}^{b} F(x, y, y') dx
\]

(НУ)

\[
y(a) = A
u_1 = (1 + \varepsilon_1) u_{xx}
J[\overline{y}] \leqslant J[\overline{y} + \varepsilon \eta]
u = u_1 + \varepsilon u_1 + \varepsilon^2 u_2 + \cdots
\varphi(\varepsilon) = J[\bar{y} + \varepsilon \eta]
\frac{d\varphi\left(\varepsilon\right)}{d\varepsilon} = \frac{d}{d\varepsilon} J\left[\overline{y} + \varepsilon\eta\right]|_{\varepsilon=0} =
\]

(Читач має робити це самостійно.) Інтегруючи на частини, ми отримуємо

\[
P_{0}
\]

Оскільки цей інтеграл є нульовим у будь-якій функції \(\eta(x)\) , що
задовольняє граничні умови \(\eta(a) = \eta(b) = 0\) , ми Р и с

\pandocbounded{\includegraphics[keepaspectratio]{images/_page_330_Figure_6.jpeg}}

. 44 3. Граф функції \(J[y+] \in \eta\) {]} в околиці
\(\varepsilon = 0\) .

Ми приходимо до висновку, що друга частина підоб'єктивного виразу має
дорівнювати нулю, тобто

(44 2)

\begin{verbatim}
,  (x,  y, z) \in R^3  ,  u = 0   u_t =
\end{verbatim}

(Рівняння Ейлера--Лагранжа).

Рівняння (44 2) називається рівнянням Ейлера-Лагранжа, і хоча загалом
воно здається комплексним, коли в нього підставляється конкретна функція
F(x, y, y'), воно стає звичайним диференціальним рівнянням другого
порядку відносно невідомої функції \(\overline{y}(x)\) . Отже, щоб
визначити мінімізуючу функцію \(\overline{y}\) , необхідно розв'язати
рівняння Ейлера-Лагранжа.

Отже, ми показали, що

якщо функція \(y(x)\) мінімізує функціональний
\(J[y] = \int_a^b F(x, y, y') dx\) (у класі гладких функцій з граничними
умовами \(y(a)\) = A; \(y(b)\) = B), то вона повинна задовольняти
рівняння

\[
g(x) = e^{-x^2}.
\]

(з граничною умовою).

(Щоб спростити позначення, ми прибрали лінію над \(\bar{y}\) .) Щоб
проілюструвати теорію, наведемо приклад.

Пошук мінімальної функціональності
\(J[y] = \int_{0}^{1} [y^2 + y'^2] dx\) .

Спробуємо знайти функцію \(y(x)\), яка проходить через точки (0,0) і
(1,1) і мінімізує функціональність

\begin{enumerate}
\def\labelenumi{\arabic{enumi}.}
\setcounter{enumi}{3}
\tightlist
\item
  Решите аналогичную (см. задачу 3) задачу для двумерного уравнения
\end{enumerate}

(УЧП)

З контексту зрозуміло, що шукана функція має бути диференційовною,
оскільки інтегративний вираз залежить від y'. Щоб визначити
\(\bar{y}(x)\) , запишемо рівняння Ейлера--Лагранжа (разом із граничною
умовою \(y(0)\)=0 та \(y(1)\)=1).

Зверніть увагу, що в перших двох задачах коефіцієнти є сталими. Читач
повинен розуміти, що параметр b має бути достатньо малим, а необережний
ряд може відрізнятися

\section{ЗАДАЧИ}\label{ux437ux430ux434ux430ux447ux438-9}

\begin{verbatim}
- 1. Подставьте разложение (46.8) в задачу (46.7) и получите последовательность задач P_0,\ P_1,\ P_2,\ \dots . 2. Покажите, что нелинейную задачу
\end{verbatim}

\(y(0) = 0,\) \(y(1) = 1.\) Але \(F(x, y, y') = y^2 + y'^2\) і, отже,

(НУ)

(Продифференцировали \(F\) по \(y\) и \(y'.\) )

Уравнение Эйлера --- Лагранжа принимает вид

\[
u(r, \theta) = u_0(r, \theta) + \frac{1}{4}u_1(r, \theta)
\]

откуда

\[
\Delta u = 0,
\]

Решая це простое дифференциальное уравнение с граничными условиями
\(y(0)\)=0 и \(y(1)\)=1, получаем

\[
q_{xx} + \varphi_{uu} = 0
\]

.

График цієї функції изображен на рис. 44.4. Ясно, що любая другая
гладкая функція, удовлетворяющая тем же граничным условиям, даст большее
значение функционала \(J(y)\).

\pandocbounded{\includegraphics[keepaspectratio]{images/_page_332_Figure_5.jpeg}}

Рис. 44 4. Все допустимые гладкие кривые удовлетворяют граничным
условиям \(y(0)\) = 0 и \(y(1)\) = 1: a---допустимая кривая; \(\delta\)
---минимизирующая функція \(y(x) = 0.42e^x - 0.42e^{-x}\) .

\section{ЗАМЕЧАНИЯ}\label{ux437ux430ux43cux435ux447ux430ux43dux438ux44f-5}

\begin{enumerate}
\def\labelenumi{\arabic{enumi}.}
\item
  Уравнение Эйлера---Лагранжа аналогично равенству нулю производной в
  дифференциальном исчислении. Якщо читатель помнит, то на самом деле
  равенство нулю производной не является достаточным условием
  экстремума. Наприклад, производная функції \(f(x) = x^3\) при x = 0
  равна нулю, но ця точка не является нп точкой минимума, ни точкой
  максимума. Так же и с уравнением Эйлера---Лагранжа: оно является лише
  необходимым условием экстремума, но не достаточным. Значит, экстремум
  может достигаться на других функциях. Очень часто, однако, розв'язок
  уравнения Эйлера---Лагранжа по существу задачи дает локальный или даже
  глобальный минимум.
\item
  Якщо у (х) минимизирует функционал
\end{enumerate}

\[
q_{uu} + q_{vv} = 0
\]

то площадь поверхности, образованной вращением кривой \(y(x)\) вокруг
оси x, буде минимальной (рис. 44.5). Решением уравнения Эйлера ---
Лагранжа в цьому случае буде

цепная линия г д

\pandocbounded{\includegraphics[keepaspectratio]{images/_page_353_Figure_9.jpeg}}

е константы \(\alpha\) и \(\beta\) определяются из условий \(y(a)\) = A
и \(y(b)\) = B (т. е. цепная линия должна проходить через заданные
точки). Решить уравнение Эйлера---Лагранжа в данном случае достаточно
трудно, поэтому мы рекомендуем читателю лише проверить, що цепная линия
ему удовлетворяет.

\begin{enumerate}
\def\labelenumi{\arabic{enumi}.}
\setcounter{enumi}{2}
\tightlist
\item
  В цієї лекции мы нашли, що функція
\end{enumerate}

\[
w = f(z)
\]

доставляет минимум функционалу П р

\pandocbounded{\includegraphics[keepaspectratio]{images/_page_355_Figure_5.jpeg}}

и цьому оказывается, що \(J[\overline{y}] = 0.46\) . Якщо подставить в
функционал любую другую гладкую функцию, график которой проходит через
точки (0, 0) и (1, 1), то значение функционала J{[}y{]} станет больше.

\pandocbounded{\includegraphics[keepaspectratio]{images/_page_333_Picture_8.jpeg}}

Рис. 44 5. Минимальная поверхность вращения.

\begin{enumerate}
\def\labelenumi{\arabic{enumi}.}
\setcounter{enumi}{3}
\tightlist
\item
  Основные законы физики чаще всего формулируются на языке вариационных
  принципов, а не дифференциальных уравнений. В качестве примеров можна
  привести принцип Ферма (свет при распространении из одной точки в
  другую выбирает путь, которому соответствует наименьшее время
  распространения) или принцип Гамильтона (в консервативном поле частица
  движется так, що интеграл действия
\end{enumerate}

\pandocbounded{\includegraphics[keepaspectratio]{images/_page_355_Figure_7.jpeg}}

(кинетическая энергия--- потенциальная энергия) \(dt\) буде
минимальным). Таким образом, физические явления развиваются лише так, що
эти функционалы принимают минимальные значения.

Основные идеи вариационного исчисления можна распространить на
функционалы, зависящие от функції нескольких независимых змінних, вида

\[
u + iv = (x + iy)^2 = x^2 - y^2 + 2ixy
\]

Соответствующее уравнение Эйлера --- Лагранжа для такого функционала
буде уже уравнением з частинними похідними

\begin{verbatim}
\left\{
\end{verbatim}

Такого рода функционалы мы рассмотрим в следующей лекции. Однако общая
философия метода буде иной. Некоторым новым способом мы найдем функцию
\(u(x, y)\), яка минимизирует функционал, а поскольку она является
решением уравнения Эйлера --- Лагранжа, то тем самым мы найдем розв'язок
уравнения з частинними похідними. Другими словами, мы будем искать
розв'язок уравнения з частинними похідними методом минимизации
функционала. В цієї лекции все було наоборот: задачу минимизации
функционала, зависящего от функції одной переменной, мы сводили к
решению уравнения Эйлера --- Лагранжа. Методы розв'язки дифференциальных
уравнений путем минимизации соответствующих функционалов принято
называть прямыми методами вариационного исчисления. Методы минимизации
функционалов путем розв'язки соответствующих уравнений Эйлера ---
Лагранжа називаються непрямыми методами вариационного исчисления.
Следующая лекция посвящена хорошо известному прямому методу Ритца.

\section{ЗАДАЧИ}\label{ux437ux430ux434ux430ux447ux438-10}

\begin{enumerate}
\def\labelenumi{\arabic{enumi}.}
\tightlist
\item
  Среди всех кривых, удовлетворяющих условиям \(y(0)\) = 0, и \(y(1)\) =
  1, найдите ту, яка минимизирует функционал
\end{enumerate}

Чи знаєте ви, як виглядає графік цієї функції з сегментом {[}0, 1{]}?
Подумай про це. Зверніть увагу, що синусоїдальне перетворення функції
\(f(x)\) визначається лише для додатних цілих чисел n.~Іншими словами,
скінченний синус і косинус перетворюють функції у чисельні
послідовності.

\pandocbounded{\includegraphics[keepaspectratio]{images/_page_185_Figure_3.jpeg}}

РИС. 25.1. Графики функции \(f(x)\) = 1 и ее преобразования.

\section{Свойства
преобразований}\label{ux441ux432ux43eux439ux441ux442ux432ux430-ux43fux440ux435ux43eux431ux440ux430ux437ux43eux432ux430ux43dux438ux439-1}

Прежде чем приступить к решению задач, мы должны получить некоторые
свойства этих преобразований.

\begin{verbatim}
Якщо you(x, t) — функція  \partial \theta yx  змінних, і ми виконуємо перетворення змінної x, то
\end{verbatim}

Как вы интерпретируете получелный результат? Чему равна величина
\(J[\overline{y}]\) ? Каков смысл величины \(J[\overline{y}]\) ?

\begin{enumerate}
\def\labelenumi{\arabic{enumi}.}
\setcounter{enumi}{1}
\tightlist
\item
  Кинетическая энергия колеблющейся материальной точки определяется
  выражением \(KE = \frac{1}{2} m\dot{y}^2\) , где y = dy/dt, а
  потенциальная равна \(PE = \frac{1}{2} ky^2\) .
\end{enumerate}

\pandocbounded{\includegraphics[keepaspectratio]{images/_page_335_Picture_2.jpeg}}

Принцип Гамильтона утверждает, що материальная точка буде двигаться так,
що интеграл

(Зверніть увагу, що перетворення було виконано на змінній x, і отримана
послідовність залежить лише від часу t.)

А как быть с производными? Ниже приводятся несколько полезных формул

буде минимален. Якщо цей закоп справедлив, то получите дифференциальное
уравнение движения материальной точки.

\begin{enumerate}
\def\labelenumi{\arabic{enumi}.}
\setcounter{enumi}{2}
\tightlist
\item
  Покажите, що в классе гладких функцій, удовлетворяющих граничным
  условиям \(y(0)\)=0 и \(y(\pi/2)=1\) , минимум функционала
\end{enumerate}

\begin{verbatim}
достигается на функції  \overline{y}(x) = \sin x  . Вычислите  J(\sin x)  . 4. Исходя из функционала
\end{verbatim}

получите уравнение Эйлера --- Лагранжа

\section{ВАРИАЦИОННЫЕ МЕТОДЫ РОЗВ'ЯЗКИ УРАВНЕНИЙ З ЧАСТИННИМИ
ПОХІДНИМИ}\label{ux432ux430ux440ux438ux430ux446ux438ux43eux43dux43dux44bux435-ux43cux435ux442ux43eux434ux44b-ux440ux43eux437ux432ux44fux437ux43aux438-ux443ux440ux430ux432ux43dux435ux43dux438ux439-ux437-ux447ux430ux441ux442ux438ux43dux43dux438ux43cux438-ux43fux43eux445ux456ux434ux43dux438ux43cux438}

ЦЕЛЬ ЛЕКЦИИ: Показать, как можна решать дифференциальное уравнение,
рассматривая его как уравнение Эйлера---Лагранжа для некоторого
функционала и находя некоторым новым методом функцию, минимизирующую цей
функционал. Тогда минимизирующая функція буде решением уравнения з
частинними похідними. Задача, конечно, состоит в том, чтобы найти такой
функционал, для которого исходное уравнение является уравнением
Эйлера---Лагранжа. Хорошо известен результат (теорема о минимуме
энергии), який гласит, що нахождение розв'язки и некоторой эллиптической
краевой задачи

, в области \(D\) , \(u = 0\) , на границе \(D\) ,

эквивалентно нахождению такой функції u (тоже обращающейся в нуль на
границе области D), яка минимизирует функционал потенинальной энергии

\begin{verbatim}
То есть,  \Delta u = f  —уравнение Эйлера—Лагранжа для J[u]. Функція, приближенно минимизирующая функционал J[u], ищется методом Ритца; тем самым мы получаем розв'язок (приближениое) уравнения з частинними похідними. В цієї лекции мы познакомимся с методом Ритца и посмотрим, как с помощью цього метода можна и и и и и и и уравнения.
\end{verbatim}

Очень хорошо решать граничные задачи (как, наприклад, задачу о
закрепленной мембране) путем отыскания гладкой поверхности, яка
минимизирует потенциальную энергию мембраны. То есть, якщо мы
рассматриваем исходное уравнение как уравнение Эйлера --- Лагранжа для
искоторого функционала J{[}u{]}, то розв'язок дифференциального
уравнения можна получить путем минимизации функционала (поскольку
минимизирующая функція функционала является решением соответствующего
уравнения

Эйлера --- Лагранжа). В лекции 44 мы рассматривали такие функционалы, у
которых уравнение Эйлера --- Лагранжа було обыкновенным дифференциальным
уравнением. В цієї лекции рассматриваются функционалы, для которых
уравнение Эйлера --- Лагранжа относится к классу уравнений з частинними
похідними. Наприклад для функционала

\[
(\text{УЧП}) \quad u_{tt} = u_{xx} + \sin(\pi x), \quad 0 < x < 1, \quad 0 < t < \infty,
\]

уравнение Эйлера---Лагранжа (подобно случаю обыкновенного
дифференциального уравнения в последней лекции) имеет вид

\[
(\text{ГУ}) \quad
\]

и, отже, для розв'язки задачи Дирихле в единичном квадрате

(УЧП)

, \(0 < x < 1\) , \(0 < y < 1\) , (ГУ) \(u = g\) , на границе квадрата,

можна, наоборот, найти такую функцию \(u(x, y)\), яка минимизирует
J{[}u{]} и равна g на границе. Вероятно, не стоит удивляться тому, що
функционал

представляет собой потенциальную эпергию мембраны и фактически мы
находим поверхность с минимальной потенциальной энерешей. Вопрос,
конечно, состоит в следующем: якщо дано дифференциальное уравнение, как
мы находим функционал, выражающий потенциальную энергию розв'язки? Ответ
на цей вопрос дает хорошо известная теорема (теорема о минимуме
энергии), яка утверждает:

Розв'язок и задачи Дирихле

\(^{\Lambda}u = f\) в области \(^{D}\) \((ГУ)\) \(^{U}=0\) на границе
\(^{D}\) является той же функциен, яка минимизирует (среди функцій,
удовлетворяющих граничному условию u=0) энергегический функционал

\[
\begin{cases} u_{tx}, \quad 0 = \begin{cases} 1, \quad -1 < x < 1, \\ 0 \quad \text{в остальных точках,} \end{cases} \end{cases}
\]

и, отже, для розв'язки задачи (45.1) достаточно найти функцию
\(\overline{u}\) (среди функцій, обращающихся в нуль на граннце),
минимизирующую \(J\{u\}\) . Этим замечанием завершается первая часть
лекции. В оставшейся части мы покажем, как находить минимизирующую
функцию \(\overline{u}\) для J{[}u{]}, используя метод Ритца.

\section{Метод Ритца минимизации
функционалов}\label{ux43cux435ux442ux43eux434-ux440ux438ux442ux446ux430-ux43cux438ux43dux438ux43cux438ux437ux430ux446ux438ux438-ux444ux443ux43dux43aux446ux438ux43eux43dux430ux43bux43eux432}

Цей метод --- один из многих вариационных методов, описанных в
литературе по вариационному исчислению. Идея, предложенная математиком
В. Ритцем, очень проста. Метод состоит из следующих шагов.

ШАГ 1. Выбираем n и функцию u, минимизирующую функционал

ищем в виде

где функції \(\phi_1, \phi_2, \phi_3, \ldots, \phi_n\) принадлежат к
классу достаточно хороших функцій и все обращаются в нуль на границе,
так що из них можна построить разумное приближение для розв'язки задачи.
Эти функції принято называть пробными функциями. Типичный набор пробных
функцій для задачи Дирихле в единичном квадрате приведен ниже:

\begin{verbatim}
— обращаются в нуль на границе,  \varphi_2(x, y) = x\varphi_1(x, y),   \varphi_3(x, y) = y\varphi_1(x, y),   \varphi_4(x, y) = x^2\varphi_1(x, y),   \varphi_5(x, y) = xy\varphi_1(x, y),   \varphi_6(x, y) = y^2\varphi_1(x, y),   \vdots   \vdots  Иными словами, первые четыре приближения имеют вид
\end{verbatim}

Решение задачи выполним за три следующих шага:

ШАГ І (Определение подходящего преобразования.)

Поскольку переменная x меняется от 0 до 1, мы воспользуемся конечным
преобразованием, а именно синус-преобразованием; позже станет ясно
почему. Мы могли бы решить эту задачу с помощью преобразования Лапласа
по переменной t (по трудоемкости этот метод не отличается от конечного
синус-преобразования).

ШАГ 2 (Выполнение преобразования).

Для удобства будем использовать обозначение \(S_n(t) = S[u]\) . Применим
синус-преобразование к исходному уравнению

\[
\begin{aligned} & u_{ti} = c^{2}u_{xx}, & -\infty < x < +\infty, & 0 < t < \infty, \\ & (\text{НУ}) & \begin{cases} u_{t}(x, 0) = 0, \\ u_{t}(x, 0) = \begin{cases} 1, & -1 < x < 1, \\ 0, & \text{в остальных случаях.} \end{cases} \end{aligned} \end{cases}
\]

и воспользуемся тождеством для синус-преобразования, получим

в координатах (x, t).

\pandocbounded{\includegraphics[keepaspectratio]{images/_page_134_Figure_7.jpeg}}

Рис. 18.4. Розв'язок задачи (18.4) в xt-плоскости.

Задача (18.4) описывает поведение струны, которой сообщается рачальная
единичная скорость на отрезке -1 \textless{} x \textless{} 1. Смещение .
.

описывается формулой Даламбера

является теперь функцией коефіцієнтів \(a_1a_2\dots a_n\) . Отже, для
того чтобы найти минимум функционала J, приравняем нулю частные
производные

\[
Aa = b
b_i = -\int_0^1 \int_0^1 f(x, y) \varphi_i(x, y) dx dy,
\]

где

Це розв'язок в различные моменты времени изображено на рис. 18.5.

\pandocbounded{\includegraphics[keepaspectratio]{images/_page_135_Figure_5.jpeg}}

Рис. 18.5. Розв'язок задачи (18.4) в различные моменты времени.

Этим завершается наша интерпретация формулы Даламбера в плоскости
змінних \(\boldsymbol{x}\) и t. В оставшейся части лекции мы займемся
решением смешанной задачи для полубесконечной струны:

(18.6)

\begin{verbatim}
Единственное отличие в том, що аппроксимирующие функції  \phi_1, \phi_2, \phi_3, \ldots, \phi_n  повинні удовлетворять граничным условиям:
\end{verbatim}

является иравнение Лапласа, но це доказательство проводится аналогично
тому, как це делалось в лекции 44 для функционала, зависящего от функції
одной переменной.

\begin{enumerate}
\def\labelenumi{\arabic{enumi}.}
\setcounter{enumi}{2}
\item
  В книгах по вариационному исчислению показано, как строить
  энергетический функционал не лише для уравнения Пуассона с граничным
  условием Дирихле u=0, но и для многих других дифференциальных
  уравнений и многих других типов граничных условий. Отже, большое число
  краевых задач можна решить, минимизируя соответствующие энергетические
  функционалы.
\item
  Чем большее значение п мы выбираем, тем меньшее значение
  \(J[\bar{u}_n]\) получается. Значит, с увеличением n растет точность л
  приближенного розв'язки дифференциального уравнения. Один из методов
  определения и, основан на последовательном вычислении функционала
  \(J[u_n]\) для возрастающих значений n.~Обчислення прекращаются, якщо
  с ростом п функционал перестает практически уменьшаться.
\item
  Якщо и велико, то проведение вычислений по методу Ритца требует
  привлечения ЭВМ. На рис. 45.2 (с. 345) приведена блок-
\end{enumerate}

схема розв'язки краевой задачи методом Ритца.

Пользователь повинен позаботиться о подпрограмме, яка должна вычислять
значения пробных функцій. Такая подпрограмма может иметь вид:

SUBROUTINE BC (X.Y.PHI)

\begin{itemize}
\tightlist
\item
  SUBROUTINE PROVIDED BY THE USER TO C - EVALUATE THE - FUNCTIONS
  PHI(1), PHI(2), \ldots, PHI(N) DIMENSION PHI(20)
\end{itemize}

PHI(1) = X*Y*(1-X)*(1-Y) \(PHI(2) = X^*PHI(1)\) (используемые функції)

\[
PHI(3) = Y\*PHI(1)
\]

PHI(N) = (що получится)RETURN END

\section{\texorpdfstring{\textbf{ЗА}ДАЧИ}{ЗАДАЧИ}}\label{ux437ux430ux434ux430ux447ux438-11}

Найдите энергетический функционал для задачи

(УЧП)

Если теперь преобразовать начальные условия краевой задачи, то получим
начальные условия для обыкновенного дифференциального уравнения

, \(0 < x < 1\) , \(0 < y < 1\) , (ГУ) \(u = 0\) на границе.

\begin{enumerate}
\def\labelenumi{\arabic{enumi}.}
\setcounter{enumi}{1}
\tightlist
\item
  Как минимизировать функционал
\end{enumerate}

методом Ритца?

УКАЗАНИЕ. Якщо ввести новую функцию \(z(x)\):

то можна заметить, що она удовлетворяет граничным условиям \(z(0)\) = 0
и \(z(1)\) = 0.1

\begin{itemize}
\tightlist
\item
  \begin{enumerate}
  \def\labelenumi{\arabic{enumi}.}
  \setcounter{enumi}{2}
  \tightlist
  \item
    Напишите программу для проведения вычислений по блоксхеме,
    изображенной на рис. 45 2.
  \end{enumerate}
\item
  \begin{enumerate}
  \def\labelenumi{\arabic{enumi}.}
  \setcounter{enumi}{3}
  \tightlist
  \item
    Покажите, що уравнение
  \end{enumerate}
\end{itemize}

является уравнением Эйлера --- Лагранжа для функционала

\begin{enumerate}
\def\labelenumi{\arabic{enumi}.}
\setcounter{enumi}{4}
\tightlist
\item
  Задачу Дирихле
\end{enumerate}

Итак, решим теперь новое семейство задач Кощи:

(ОДУ)

где A=0.06 и B=0.04. Как найти потенциальную энергию цього розв'язки? Мы
рекомендуем читателю вспомнить синуспреобразование и самостоятельно
получить розв'язок \(u\left(x,y\right)\) .

*) При такой замене возникает особенность в т. x=1. Лучше ввести функцию
\(z(x)\) по формуле \(z(x)\)=\(y(x)\)=x. Прим. ред.

\pandocbounded{\includegraphics[keepaspectratio]{images/_page_343_Figure_2.jpeg}}

Далсе. Приближенное розв'язок
\(u_n(x,y) = a_1 \varphi_1 + a_2 \varphi_2 + \dots + a_n \varphi_n\) . В
це время можна численно найти эначение функционала \(J[u_n]\) . Якщо
необходимо найти значение розв'язки \(u_n(x,y)\) в некоторой точке,
можна написать небольшую программу эля быполнения отой операции. При
необходимости можна повторить обчислення по цієї программе с другими
значениями п

РИС 452. Блок-схема метода Ритца.

\section{РОЗВ'ЯЗОК УРАВНЕНИЙ З ЧАСТИННИМИ ПОХІДНИМИ МЕТОДАМИ ТЕОРИИ
ВОЗМУЩЕНИЙ}\label{ux440ux43eux437ux432ux44fux437ux43eux43a-ux443ux440ux430ux432ux43dux435ux43dux438ux439-ux437-ux447ux430ux441ux442ux438ux43dux43dux438ux43cux438-ux43fux43eux445ux456ux434ux43dux438ux43cux438-ux43cux435ux442ux43eux434ux430ux43cux438-ux442ux435ux43eux440ux438ux438-ux432ux43eux437ux43cux443ux449ux435ux43dux438ux439}

ЦЕЛЬ ЛЕКЦИИ: Показать, как различные сложные задачи (нелинейные
уравнения с змінними коэффициентами, области неправильной формы и т. д.)
можна решить методом возмущения более простых задач. Иными словами, мы
хотим показать, как следует изменить розв'язок простой задачи, чтобы оно
давало приближенное розв'язок трудной задачи.

Наприклад, мы покажем, как нелинейную задачу

(УЧП)

\[
,  0 < r < 1  ,
\]

\begin{verbatim}
(ГУ)  u(1, \theta) = \cos \theta  ,  0 \le \theta \le 2\pi  .
\end{verbatim}

Очень часто задачу, розв'язок которой не известно, можна путем
непрерывного изменения каких-то параметров свести к близкой, но легко
решаемой задаче. Тогда и розв'язок исходной задачи буде близко решению
модифицирующей задачи.

Наприклад, рівняння Лапласа

в помощью семейства уравнений

\[
,  0 \le \varepsilon \le 1  ,
\]

, \(0 < r < 1\) , \(0 \le \theta < 2\pi\) , \(u(1, \theta) = g(\theta)\)
, \(0 \le \theta < 2\pi\) ,

при одних граничных условиях (наприклад, \(g(\theta) = \cos \theta\) )
имеет два различных вещественных розв'язки, а при других (наприклад,
\(g(\theta) = A \cos \theta\) с A \textgreater{} 20,65) ця задача не
имеет ни одного вещественного розв'язки. --- Прим. ред.

4) Краевая задача

можна непрерывно модифицировать в нелинейное уравнение
\(\Delta u = u^2 = 0\) 1).

На рис. 46.1 изображена схема розв'язки нелинейной задачи путем сведения
ее к решению уравнения Лапласа2)

\[
\right.
\]

Рис. 46.1. Схема метода возмущении.

Чтобы найти розв'язок возмущенного уравнения Лапласа u

\pandocbounded{\includegraphics[keepaspectratio]{images/_page_362_Figure_1.jpeg}}

\[
(0, t) = f(t),
u_{x}(0, t) = 0.
x + ct = \text{const},
\]

\pandocbounded{\includegraphics[keepaspectratio]{images/_page_362_Figure_2.jpeg}}

\(0 < r < 1,\) \(0 \le \theta < 2\pi,\) \(u(1, 0) = 0,\)
\(0 \le \theta < 2\pi,\) нмеет два розв'язки (одно --- нулсьое, другое
--- нетривиальное), в то время как задача u

\pandocbounded{\includegraphics[keepaspectratio]{images/_page_362_Figure_5.jpeg}}

\[
(1, \theta) = 0, \quad 0 \le \theta < 2\pi,
\]

\pandocbounded{\includegraphics[keepaspectratio]{images/_page_362_Figure_7.jpeg}}

\begin{verbatim}
\begin{array}{ccccccccccccccccccccccccccccccccccc
\end{verbatim}

\[
\begin{array}{ccc} (\text{УЧ}\Pi) & \Delta u = 0, & 0 < r < 1, \\ (\text{ГУ}) & u(1, \theta) = \cos \theta, & 0 \leqslant \theta \leqslant 2\pi, \end{array}
u_{tt} = c^2 u_{xx}, \quad 0 < x < L, \quad 0 < t < \infty,
H(x-a) =
u(0, t) = g_1(t),
\begin{cases} 0, x < a, \\ 1, x \geqslant a, \end{cases}
u_x(0, t) = g_1(t),
\]

В результате получаем

\[
u_x(0, t) - \gamma_1 u(0, t) = g_1(t),
H(x-a) =
u(0, t) = g_1(t),
\begin{cases} 0, & x < a, \\ 1, & x \geqslant a, \end{cases}
u_x(0, t) = g_1(t),
\begin{cases} 0, & n = 2, 4, \dots, \\ (4/n\pi)\cos(n\pi t), & n = 3, 5, 7, \dots \end{cases}
u_x(0, t) - \gamma_1 u(0, t) = g_1(t),
f(x) = \mathcal{F}^{-1}[F] = \frac{1}{\sqrt{2\pi}} \int_{-\infty}^{+\infty} F(\omega) e^{-i\omega x} d\omega \qquad F(\omega) = \mathcal{F}[f] = \frac{1}{\sqrt{2\pi}} \int_{-\infty}^{+\infty} f(x) e^{-i\omega x} dx
\begin{cases} \Delta u_{1} = -r^{2} \cos^{2} \theta = -\frac{r^{2}}{2} \left[1 + \cos(2\theta)\right] = -\frac{r^{2}}{2} - \frac{r^{2}}{2} \cos(2\theta), \\ u_{1}(1, \theta) = 0. \end{cases}
\]

Следовательно, решение нашей задачи можно представить в виде

\[
\Delta u = -\frac{r^2}{2} - \frac{r^3}{2} \cos{(2\theta)}
\begin{cases} 1, |x| < a \\ 0, |x| > a \end{cases}
u_p(r, \theta) = Ar^4 + Br^4 \cos(2\theta).
\]

\section{ЗАМЕЧ
АНИЯ}\label{ux437ux430ux43cux435ux447-ux430ux43dux438ux44f-1}

\begin{enumerate}
\def\labelenumi{\arabic{enumi}.}
\tightlist
\item
  Для решения задачи методом конечного синус- или косинуспреобразования
  краевые условия при x=0 и x=L должны иметь вид
\end{enumerate}

\[
u_p(r, \theta) = -\frac{r^4}{32} - \frac{r^4}{24} \cos(2\theta).
\begin{cases} 1, |x| < 1 \\ 0, |x| > 1 \end{cases}
\sum_{n=0}^{\infty} [a_n \cos(n\theta) + b_n \sin(n\theta)] - \frac{1}{32} - \frac{1}{24} \cos(2\theta) = 0.
\]

Отметим, що якщо величина \(u(0, t)\) положительна, то величина
\(u_x\)(0, t) також положительна, якщо \(u(L, t)\) положительна, то
\(u_x\)(L, t)

\pandocbounded{\includegraphics[keepaspectratio]{images/_page_144_Picture_4.jpeg}}

Рис. 19.6. Возникновение неоднородных граничных условий упругого
закрепления: h---коефіцієнт упругости пружины.

отрицательна. Однородные граничні умови (19.2) можна записать в виде

(19.3)

\[
P_{2}
\delta(x-a)
\Delta u = 0, \quad 0 < r < 1 + \frac{1}{4} \sin \theta,
\]

\begin{verbatim}
\begin{array}{ll} (\mathrm{Y} \mathrm{U} \Pi) & u_t = u_{xx}, \quad 0 < x < 1, \quad 0 < t < \infty, \\ (\mathrm{\Gamma} \mathrm{Y} & \left\{ \begin{array}{ll} u_x \left( 0, \ t \right) = 0 \\ u_x \left( 1, \ t \right) = 0, \end{array} \end{array}
\end{verbatim}

досягається в кінцевій точці сімейства граничних умов виду (див. рисунок
46 3)

\[
(1+x^2)^{-1}
Tu_x(0, t) = -\beta u_t(0, t).
xe^{-a+x^4}, a > 0
Tu_{x}(0, t) = \varphi[u(0, t)],
\]

\(H(x+a)\) - \(H(x-a)\)

\[
Tu_{x}(0, t) = -hu^{3}(0, t)
\frac{a}{x^2+a^2}
mu_{tt}(L, t) = -ku_x(L, t) + mg.
\]

\begin{verbatim}
\right. & 0 < t < \infty, \\ (\text{НУ}) & u(x, 0) = 0, \quad 0 \leqslant x \leqslant 1. \end{array}
\end{verbatim}

\[
P_{0}
\frac{2ax}{(x^2+a^2)^2}
\]

Отже, розв'язуючи кожну з задач Дірікла всередині кола для функцій
\(u_0\) , \(u_1\) , \(u_2\) , отримуємо розв'язок

(46.9)

\[
\begin{cases} \cos(ax), |x| > \pi/2a \\ 0, |x| > \pi/2a \end{cases}
\]

Задачі Діріхле в деформованій області (зверніть увагу, що тут

Ми вже взяли \(\epsilon = 1/4\) :)

Як вправу, читача просять знайти пертурбацію першого порядку \(u_1\) і
перевірити, наскільки добре є приблизне розв'язання я к

\pandocbounded{\includegraphics[keepaspectratio]{images/_page_189_Picture_10.jpeg}}

що розглянути рівняння з параметром

(46.11)

\begin{enumerate}
\def\labelenumi{\arabic{enumi}.}
\setcounter{enumi}{17}
\tightlist
\item
\end{enumerate}

і побачити його розв'язок у вигляді

\[
\begin{cases} 1-|x|, |x| < 1 \\ 0, |x| > 1 \end{cases}
\]

Підставляючи цю декомпозицію на (46.11), отримуємо таку послідовність
завдань

\[
u_{x}(0, t) = \frac{h}{T} [u(0, t) - \theta_{1}(t)]
u_{tt} = u_{xx}, \quad 0 < x < 1, \quad 0 < t < \infty,
\cos(ax)
\Delta u = 0,
\frac{2}{\pi} \frac{\sin a\omega}{\omega}
q_{uu} + q_{vv} = 0
\]

Перед нами неоднородное уравнение теплопроводности, так что метод
разделения переменных неприменим. Можпо, конечно, воспользоваться
конечным синус-преобразованием по переменной x или преобразованием
Лапласа по переменной t, но вместо этого мы рассмотрим две подзадачи

\[
w = f(z)
\frac{2}{\pi} \frac{\sin a\omega}{(\omega^2+a^2)^2}
\]

\begin{verbatim}
\left\{
\end{verbatim}

\[
\frac{\pi}{2} e^{-a+\omega}
\Rightarrow.
\begin{aligned} (Y & \Pi) & \varphi_{xx} + \varphi_{yy} = 0, & -\infty < x < \infty, & 0 < y < \infty, \\ (47.1) & \varphi(x, 0) = \begin{cases} 0, & |x| > 1, \\ 1, & |x| \le 1. \end{cases} \end{aligned}
w = \ln\left\{\frac{z-1}{z+1}\right\}
\frac{\pi}{2} e^{-a+\omega}
\]

\begin{verbatim}
\begin{cases} (\text{УЧП}) & \varphi_{uu} + \varphi_{vv} = 0, \quad -\infty < u < \infty, \quad 0 < v < \pi, \\ (\text{47.2}) & \begin{cases} \varphi(u, 0) = 0, \\ \varphi(u, \pi) = 1. \end{cases}
\end{verbatim}

\[
\frac{\pi}{2} e^{-a+\omega}
u(x, t) = \sum_{n=1}^{\infty} c_n X_n(x) T_n(t),
\frac{\pi}{2} e^{-a+\omega}
u_{tt} = \alpha^2 u_{xx}, \quad 0 < x < L, \quad 0 < t < \infty,
\frac{\pi}{2} e^{-a+\omega}
(\Gamma Y)
\]

\begin{verbatim}
Каждая из этих задач легко решается и в то же время ясно, что  \mathit{suma}  решений задач  P_1  и  P_2  является решением исходной вадачи, т. е.
\end{verbatim}

Щоб візуально зрозуміти, як працює ця функція, ми радимо читачеві
побудувати її за кількома лініями y=c.

У другому прикладі ми перетворюємо область між двома неконцентричними
колами у кільце.

Задача Діріхле в області між двома некопоцентричними колами

Припустимо, ви хочете знайти потенціал між двома колами

\[
\frac{2ax}{(x^2+a^2)^2}
\quad 0 < t < \infty,
\]

Решение задачи \(P_{i}\) Розв'язання задачі \(P_i\) Розв'язання задачі
\(P_i\) Можна показати, що в загальному випадку розв'язок задачі

(HY)

\[
\frac{\cos(ax), |x| > \pi/2a}{0, |x| > \pi/2a}
u = \gamma [(x-s)(x-t)-\gamma^2],
\frac{1-|x|, |x| < 1}{0, |x| > 1}
\quad 0 \le x \le L,
u(x, t) = X(x) T(t).
T'' - \alpha^2 \lambda T = 0,
\cos(ax)
X'' - \lambda X = 0,
\frac{\pi}{2} [\delta(\omega+a) + \delta(\omega-a)]
u(x, t) = [C \sin(\beta x) + D \cos(\beta x)][A \sin(\alpha \beta t) + B \cos(\alpha \beta t)]
\beta_n = \frac{n\pi}{L}, \quad n = 0, 1, 2, \ldots
u_n(x, t) = X_n(x) T_n(t) = \sin(n\pi x/L) [a_n \sin(n\pi at/L) + b_n \cos(n\pi at/L)],
f(x) = \int_{0}^{\infty} F(\omega) \sin(\omega x) d\omega
u_n(x, t) = R_n \sin(n\pi x/L) \cos(n\pi\alpha (t - \delta_n)/L],
F(\omega) = \frac{?}{\pi} \int_{0}^{\infty} f(x) \sin(\omega x) dx
u(x, t) = \sum_{n=1}^{\infty} \sin(n\pi x/L) \left[ a_n \sin(n\pi \alpha t/L) + b_n \cos(n\pi \alpha t/L) \right]
0 < x < \infty
u(x, \theta) = f(x),
\sum_{n=1}^{\infty} b_n \sin(n\pi x/L) = f(x),
\sum_{n=1}^{\infty} a_n (n\pi \alpha/L) \sin(n\pi x/L) = g(x).
0 < \omega < x
\int_{0}^{L} \sin(m\pi x/L) \sin(n\pi x/L) dx =
-\omega^2 F(\omega) + \frac{2}{\pi} \omega f(0)
\begin{cases} 1, |x| < 1 \\ 0, |x| > 1 \end{cases}
\frac{1}{\alpha} F\left(\frac{\omega}{\alpha}\right)
\]

\(f(x)\) * \(g(x)\)

проследим, как оно решается методом конечного синус-преобразования. Вот
что мы должны сделать: разложить еходное воздействие \(f(x, t)\) на
компоненты, найти отклик \(U_n\) на каждую компоненту, а затем сложить
эти отклики. Математически это может быть и не совсем очевидно, но
давайте внимательно проследим за ходом решения. Разложим входящие в
уравнение с частными производными члены в ряды Фурье по синусам:

\[
(1+x^2)^{-1}
\frac{2\omega}{\pi (\alpha^2 + \omega^2)}
xe^{-a+x^4}, a > 0
x^{-1/2}
\]

\(H(x+a)\) - \(H(x-a)\)

\[
[2/\pi\omega]^{1/2}
\frac{a}{x^2+a^2}
\]

where

\[
\frac{2ax}{(x^2+a^2)^2}
\frac{2}{\pi\omega} [1-\cos(\omega a)]
\begin{cases} \cos(ax), |x| > \pi/2a \\ 0, |x| > \pi/2a \end{cases}
\]

1

\[
\begin{cases} 1-|x|, |x| < 1 \\ 0, |x| > 1 \end{cases}
e^{-1}a + \omega
\cos(ax)
\]

Зверніть увагу, по-перше, що коефіцієнти \(A_n\) , \(B_n\) \(F_n\)
дійсно є функціями часу t, оскільки ми застосували синусоїдальне
перетворення до функції двох змінних x, t, u, а по-друге, ми розклали
вхідний ефект \(f(x,t)\) на прості компоненти \(F_n\)(t) . Тепер
знайдіть відповіді \(U_n\)(t) для кожного компонента \(F_n\)(t) , потім
додайте всі відповіді \(U_n\)(t) і отримайте розв'язок \(u(x,t)\).

Щоб знайти відповіді \(U_n\)(t) , необхідно робити невеликі розрахунки з
коефіцієнтами \(A_n\)(t) і \(B_n\)(t) , щоб у вирази інтегрування була
включена лише функція u, а не її похідні \(u_i\) і \(u_{xx}\) .
Інтегруючи на частини, ми отримуємо

\[
\frac{2}{\pi} \frac{\sin a\omega}{\omega}
\frac{1}{2}e^{-\alpha}\sin\omega
\frac{2}{\pi} \frac{\sin a\omega}{(\omega^2+a^2)^2}
\frac{1-e^{-a\omega}}{\omega}
\frac{\pi}{2} e^{-a+\omega}
\frac{\pi}{2} e^{-a+\omega}
\frac{\pi}{2} e^{-a+\omega}
\frac{2}{\pi}F''(\omega)
\frac{\pi}{2} e^{-a+\omega}
f(x) = \int_{0}^{\infty} F(\omega) \cos(\omega x) d\omega \qquad \qquad F(x) = \frac{2}{\pi} \int_{0}^{\infty} f(x) \cos(\omega x) dx
\frac{\pi}{2} e^{-a+\omega}
0 < x < \infty \qquad \qquad 0 < \omega < \infty
\frac{\cos(ax), |x| > \pi/2a}{0, |x| > \pi/2a}
\frac{1-|x|, |x| < 1}{0, |x| > 1}
\cos(ax)
\]

j'\,'(x)

\[
u_{tt} = -\alpha^2 u_{xxxx},
S_n = \frac{2}{\pi} \int_0^{\pi} f(x) \sin(nx) dx.
u(0, t) = 0,
\]

в последнее соотношение, получаем

\[
u_{xx}(0, t) = 0,
y = \frac{\pi}{b - a} (x - a).
-\omega^2 F(\omega) + \frac{2}{\pi} \omega f(0)
\]

и, следовательно, (26.1) принимает вид

\[
\frac{2\omega}{\pi (\alpha^2 + \omega^2)}
f(x) = \sum_{n=1}^{\infty} S_n \sin(nx).
u(x, t) = X(x) [A \sin(\omega t) + B \cos(\omega t)].
\]

\begin{verbatim}
Поскольку это разложение является тождеством по x, его коэффициенты должны быть равны нулю,  \tau  . е.
\end{verbatim}

\[
X^{(1V)} - \omega^2 X = 0,
f(x) = \sum_{n=1}^{\infty} S_n \sin(nx)
X(x) = C\cos\sqrt{\omega}x + D\sin\sqrt{\omega}x + E \cosh\sqrt{\omega}x + F \sinh\sqrt{\omega}x.
\]

.

Таким образом, мы получили связь между входом и выходом системы, между
коэффициентами \(F_n\) и \(U_n\) . Прежде чем решать уравнения для
\(U_n\)(t) , обратимся к начальному условию

\begin{enumerate}
\def\labelenumi{\arabic{enumi}.}
\setcounter{enumi}{7}
\tightlist
\item
  \(\frac{x}{x^4+4}\)
\end{enumerate}

\[
S_n = \frac{2}{\pi} \int_{0}^{\pi} f(x) \sin(nx) dx
\]

\begin{enumerate}
\def\labelenumi{\arabic{enumi}.}
\setcounter{enumi}{9}
\tightlist
\item
  \(-x^2/(x)\)
\end{enumerate}

Если разложить \(u(x, 0)\) в ряд Фурье по синусам и приравнять его нулю:

\section{ТАБЛИЦЯ C. Перетворення
Фур'є}\label{ux442ux430ux431ux43bux438ux446ux44f-c.-ux43fux435ux440ux435ux442ux432ux43eux440ux435ux43dux43dux44f-ux444ux443ux440ux454}

\begin{verbatim}
то определяются начальные условия для функции  U_n\left(t\right)  , т. е.
\end{verbatim}

Решая их, находим

\[
n = 1, 2, ...
\]

, \(\sin \sqrt{\omega} = 0 \Rightarrow \omega = (n\pi)^2\) ,
\(n = 1, 2, \dots\) Другими словами, \emph{собственные частоты} свободно
опирающейся балки равны

\begin{verbatim}
,  n = 1,  2, \dots  Мы получили разложение нашей смешанной задачи на ряд простых задач об отклике системы на элементарные входные воздействия:
\end{verbatim}

(OДУ)

а фундаментальные розв'язки и, (удовлетворяют уравнению и граничным
условиям (21.2)) записываются в виде

(HУ) \(U_n\)(0) = 0, n = 1, 2, \ldots{} Каждое из этих уравнений легко
решается методом интегрирующего множителя (или с помощью преобразования
Лапласа). В любом случае получаем

А теперь, поскольку уравнение и граничні умови линейны и однородны,
можна утверждать, що сумма

(21.4)

\[
\frac{\pi}{2} e^{ax}
\]

також удовлетворяет уравнению и граничным условиям. Единственное, що нам
осталось---це выбрать коефіцієнти \(a_n\) и \(b_n\) так, чтобы
удовлетворить начальным условиям. Подставляя (21.4) в початкові умови,
получаем

(21.5)

Таким образом, мы получили отклики \(U_n\)(t) на простейшие входные
воздействия \(F_n\)(t) . Для получения решения исходной задачи осталось
совершить последний шаг---сложить все отклики

\[
n = 1, 2, ...
C_n = \frac{2}{\pi} \int_0^{\pi} f(x) \cos(nx) dx.
\]

f'\,'(x)

\begin{verbatim}
Читатель, должно быть, заметил, что каждый отклик  U_n(t)  умножается на весовой множитель  \sin(n\pi x)  . Конечно, в разложении функции f(x, t) на компоненты  F_n(t)  присутствуют те же весовые множители.
\end{verbatim}

\section{ЗАМЕЧАНИЕ}\label{ux437ux430ux43cux435ux447ux430ux43dux438ux435}

Використання скінченного синусоїдального перетворення дозволяє записати
розв'язок у вигляді ряду, але для більшості інших інтегральних
перетворень розв'язки представлені як інтеграли (неперервний розклад).
Читача можна попросити дати інтерпретацію проблеми

Таким образом, розв'язок задачи имеет вид (21.4), а коефіцієнти \(a_n\)
и \(b_n\) вычисляются по формулам (21.6)

Чтобы помочь читателю разобраться в цієї задаче, приведем простой
пример.

\section{Модель колеблющейся
балки}\label{ux43cux43eux434ux435ux43bux44c-ux43aux43eux43bux435ux431ux43bux44eux449ux435ux439ux441ux44f-ux431ux430ux43bux43aux438}

Рассмотрим свободно опирающуюся балку, изображенную на рис. 21.2. Пусть
початкові умови таковы:

\[
y = \frac{\pi}{b-a}(x-a).
\]

\(u_1(x, 0) = 0\) Мы могли бы найти розв'язок, подставив функції
\(f(x)\) и \(g(x)\) в формулы (21.6), однако легко видеть, що

\[
f(x) = \frac{C_0}{2} + \sum_{n=1}^{\infty} C_n \cos(nx).
\]

для всех \(n=1,\,2,\,\ldots,\) \(b_1=1,\,b_2=0,\,b_3=0,5\)
\(b_n=0,\,n=4,\,5,\,\ldots\) Отже, розв'язок имеет вид

\[
f(x) = \frac{C_0}{2} + \sum_{n=1}^{\infty} C_n \cos(nx) \qquad C_n = \frac{2}{\pi} \int_{0}^{\pi} f(x) \cos nx \, dx
\]

.

Интересно сравнить це розв'язок с решением для струны при тех же
начальных условиях. Якщо заглянуть в лекцию 20, там можна найти
розв'язок соответствующей задачи для струны:

\[
0 \le x \le \pi \qquad n = 0, 1, 2, \dots
\]

.

\pandocbounded{\includegraphics[keepaspectratio]{images/_page_158_Figure_1.jpeg}}

РИС 21.2. Простые колебания свободно опирающейся балки.

Другими словами, балка колеблется на более высоких частотах, чем струна.
Отметим, що в том и другом случаях высшие частоты кратны основным
частотам колебаний.

\section{ЗАМЕЧАНИЯ}\label{ux437ux430ux43cux435ux447ux430ux43dux438ux44f-6}

\begin{itemize}
\tightlist
\item
  \begin{enumerate}
  \def\labelenumi{\arabic{enumi}.}
  \tightlist
  \item
    Вообще говоря, балку можна закрепить тремя способами:
  \end{enumerate}
\item
  а) оставить свободной (незакрепленной),
\item
  б) свободно опереть.
\item
  в) жестко закрепить.
\end{itemize}

Схематически эти способы закрепления изображены на рис. 21.3.

\begin{enumerate}
\def\labelenumi{\arabic{enumi}.}
\setcounter{enumi}{1}
\tightlist
\item
  Еще одна важная задача --- колебания балки с заделанным концом (см.
  рис. 213). Благодаря нестандартным граничным условиям
\end{enumerate}

\(u_x\)(0, t) = 0, \(u_{xx}\)(1, t) = 0, \(u_{xx}\)(1, t) = 0, розв'язок
уже не буде выражаться в виде ряда из произведений синусов и косинусов.
Оно буде иметь более сложную структуру

где власні функції \(X_n(x)\) являются линейными комбинациями синусов,
косинусов, гиперболических синусов и гиперболических косинусов.
Розв'язок цієї задачи можна найта в руководствах по теории упругости.

\pandocbounded{\includegraphics[keepaspectratio]{images/_page_159_Figure_2.jpeg}}

Рис. 21.3. Типичные задачи для балки: a---защемленная балка; b---балка,
жестко заделанная с обоих концов; a---левый конец жестко заделан, правый
свободно опирается.

\section{ЗАДАЧИ}\label{ux437ux430ux434ux430ux447ux438-12}

\begin{enumerate}
\def\labelenumi{\arabic{enumi}.}
\tightlist
\item
  Решите задачу для балки с закрепленным концом.
\end{enumerate}

f'\,'(x)

\begin{enumerate}
\def\labelenumi{\arabic{enumi}.}
\setcounter{enumi}{5}
\tightlist
\item
\end{enumerate}

\$\$ \textbackslash begin\{cases\}

\end{document}
